% ============================================================
%  From Zero to Propagators
%  Orders, Lattices, and Propagator Networks in Lean 4
% ============================================================
\documentclass[12pt, a4paper, openright]{book}

% ── Fonts & encoding ────────────────────────────────────────
\usepackage{fontspec}
\setmainfont{Latin Modern Roman}
% NB: load DejaVu Sans Mono from the system files explicitly — the
% TeX-distribution copy is trimmed and lacks ⊓ (U+2293) and ⊔ (U+2294),
% which the Lean listings need.
\setmonofont[Scale=0.82,
             Path=/usr/share/fonts/truetype/dejavu/,
             Extension=.ttf,
             UprightFont=DejaVuSansMono,
             BoldFont=DejaVuSansMono-Bold,
             ItalicFont=DejaVuSansMono,
             ItalicFeatures={FakeSlant=0.2},
             BoldItalicFont=DejaVuSansMono-Bold,
             BoldItalicFeatures={FakeSlant=0.2}
            ]{DejaVuSansMono}

% ── Layout ──────────────────────────────────────────────────
\usepackage[top=1in, bottom=1.15in, left=1.1in, right=1.1in, headheight=14pt]{geometry}
\usepackage{microtype}
\usepackage{placeins}
\usepackage{setspace}
\setstretch{1.08}
\usepackage{parskip}
\setlength{\parskip}{6pt}

% ── Math ────────────────────────────────────────────────────
\usepackage{amsmath}
\usepackage{amssymb}
\usepackage{mathtools}
% unicode-math must come AFTER amssymb (it redefines symbols the
% ams packages would otherwise re-declare).
\usepackage{unicode-math}
\setmathfont{Latin Modern Math}

% ── Theorems ────────────────────────────────────────────────
\usepackage{amsthm}
\newtheoremstyle{defstyle}{8pt}{6pt}{\normalfont}{0pt}{\bfseries}{.}{.5em}{}
\newtheoremstyle{thmstyle}{8pt}{6pt}{\itshape}{0pt}{\bfseries}{.}{.5em}{}
\theoremstyle{thmstyle}
\newtheorem{theorem}{Theorem}[chapter]
\newtheorem{lemma}[theorem]{Lemma}
\newtheorem{corollary}[theorem]{Corollary}
\newtheorem{proposition}[theorem]{Proposition}
\theoremstyle{defstyle}
\newtheorem{definition}[theorem]{Definition}
\newtheorem{example}[theorem]{Example}
\newtheorem{remark}[theorem]{Remark}

% ── Code (minted) ───────────────────────────────────────────
\usepackage{minted}
\usemintedstyle{friendly}
\setminted[lean4]{
  bgcolor=gray!6,
  frame=lines,
  framesep=4pt,
  rulecolor=gray!40,
  fontsize=\small,
  breaklines=true,
  linenos=true,
  numbersep=6pt,
  xleftmargin=14pt,
}
\setminted[text]{
  bgcolor=yellow!8,
  frame=single,
  framesep=4pt,
  rulecolor=gray!30,
  fontsize=\small,
  linenos=false,
  xleftmargin=6pt,
}
% Inline lean
\newcommand{\lc}[1]{\mintinline{lean4}{#1}}

% ── Diagrams ────────────────────────────────────────────────
\usepackage{tikz}
\usetikzlibrary{positioning, arrows.meta, shapes.geometric, calc, decorations.pathreplacing}
\tikzset{
  node/.style={circle, draw=gray!80, fill=gray!10, inner sep=3pt, minimum size=7mm, font=\small},
  edge/.style={-{Stealth[scale=0.8]}, gray!60, thick},
  hasseedge/.style={draw=gray!70, thick},
}

% ── Coloured boxes ──────────────────────────────────────────
\usepackage{tcolorbox}
\tcbuselibrary{theorems, breakable, skins, most}

\tcbset{
  defbox/.style={
    enhanced, breakable,
    colback=blue!4, colframe=blue!55!black,
    fonttitle=\bfseries, title={Definition},
    attach boxed title to top left={xshift=8pt, yshift=-2pt},
    boxed title style={colback=blue!55!black, colframe=blue!55!black},
    boxrule=0.6pt, arc=3pt, left=6pt, right=6pt, top=8pt, bottom=6pt,
  },
  thmbox/.style={
    enhanced, breakable,
    colback=green!3, colframe=green!50!black,
    fonttitle=\bfseries,
    boxrule=0.6pt, arc=3pt, left=6pt, right=6pt, top=6pt, bottom=6pt,
  },
  exbox/.style={
    enhanced, breakable,
    colback=orange!4, colframe=orange!70!black,
    fonttitle=\bfseries, title={Example},
    attach boxed title to top left={xshift=8pt, yshift=-2pt},
    boxed title style={colback=orange!70!black, colframe=orange!70!black},
    boxrule=0.6pt, arc=3pt, left=6pt, right=6pt, top=8pt, bottom=6pt,
  },
  exbox_notitle/.style={
    enhanced, breakable,
    colback=orange!4, colframe=orange!70!black,
    boxrule=0.6pt, arc=3pt, left=6pt, right=6pt, top=6pt, bottom=6pt,
  },
  exercisebox/.style={
    enhanced, breakable,
    colback=purple!4, colframe=purple!60!black,
    fonttitle=\bfseries,
    boxrule=0.6pt, arc=3pt, left=6pt, right=6pt, top=6pt, bottom=6pt,
  },
  notebox/.style={
    enhanced,
    colback=yellow!8, colframe=yellow!70!black,
    boxrule=0.4pt, arc=3pt, left=6pt, right=6pt, top=4pt, bottom=4pt,
  },
}

\newenvironment{defbox}[1][]{%
  \begin{tcolorbox}[defbox, title={Definition\ifx\\#1\\\else\ --- #1\fi}]%
}{\end{tcolorbox}}

\newenvironment{thmbox}[1][]{%
  \begin{tcolorbox}[thmbox, title={\ifx\\#1\\Theorem\else#1\fi}]%
}{\end{tcolorbox}}

\newenvironment{exbox}[1][]{%
  \begin{tcolorbox}[exbox, title={Example\ifx\\#1\\\else\ --- #1\fi}]%
}{\end{tcolorbox}}

\newenvironment{notebox}{%
  \begin{tcolorbox}[notebox]%
}{\end{tcolorbox}}

\newcounter{exercise}[chapter]
\newenvironment{exercisebox}[1][]{%
  \refstepcounter{exercise}%
  \begin{tcolorbox}[exercisebox,
    title={Exercise \thechapter.\theexercise\ifx\\#1\\\else\ --- #1\fi}]%
}{\end{tcolorbox}}

% ── Fancy chapter headings ───────────────────────────────────
\usepackage{titlesec}
\titleformat{\chapter}[display]
  {\huge\bfseries}
  {\large\scshape\chaptertitlename\ \thechapter}
  {10pt}
  {\Huge\bfseries}
\titlespacing*{\chapter}{0pt}{20pt}{30pt}

% ── Headers & footers ───────────────────────────────────────
\usepackage{fancyhdr}
\pagestyle{fancy}
\fancyhf{}
\fancyhead[LE]{\small\itshape\leftmark}
\fancyhead[RO]{\small\itshape\rightmark}
\fancyfoot[C]{\small\thepage}
\renewcommand{\headrulewidth}{0.4pt}

% ── Hyperlinks ──────────────────────────────────────────────
\usepackage[colorlinks=true, linkcolor=blue!60!black,
            citecolor=green!50!black, urlcolor=blue!70!black,
            bookmarks=true, bookmarksnumbered]{hyperref}

% ── Misc ────────────────────────────────────────────────────
\usepackage{enumitem}
\usepackage{booktabs}
\usepackage{xspace}

% Lean-specific math notation
\newcommand{\lean}[1]{\texttt{#1}}

% ============================================================
\begin{document}
\frontmatter

% ── Title page ──────────────────────────────────────────────
\begin{titlepage}
  \centering
  \vspace*{2cm}
  {\Huge\bfseries From Zero to Propagators\par}
  \vspace{0.5cm}
  {\large\itshape Orders, Lattices, and Propagator Networks in Lean~4\par}
  \vspace{3cm}
  \begin{tcolorbox}[width=0.75\textwidth,
                    colback=gray!8, colframe=gray!50, arc=4pt,
                    boxrule=0.6pt]
    \small\centering
    \textit{``The propagator model is to functional programming\\
    what constraint logic programming is to logic programming:\\
    a radical generalisation that turns programs into networks\\
    of co-operating, partial knowledge sources.''}\\[4pt]
    --- Inspired by Radul \& Sussman (2009)
  \end{tcolorbox}
  \vfill
  {\large A Self-Contained Journey Through Order Theory\\
  to Verified Constraint Propagation\par}
  \vspace{1cm}
  {\normalsize Lean~4 code requires no external libraries beyond the standard prelude.\par}
  \vspace{0.5cm}
  {\small\color{gray} Edition 1.0\par}
\end{titlepage}

% ── Preface ──────────────────────────────────────────────────
\chapter*{Preface}
\addcontentsline{toc}{chapter}{Preface}

This book has a single grand ambition: to take you from the very first
principles of ordered structures all the way to building a working
\emph{propagator network} --- a powerful computational model that arises
naturally from the mathematics of lattices --- and to do it with
complete rigour in Lean~4.

\medskip
\textbf{Who this book is for.}
We assume you are comfortable reading mathematics at the level of a second-year
undergraduate: sets, functions, logical quantifiers, and basic proof
techniques are fair game. No prior knowledge of abstract algebra or Lean is
required, though some familiarity with functional programming will help when
reading the longer code sections.

\medskip
\textbf{The Lean~4 philosophy here.}
Every definition and theorem in this book has a corresponding Lean~4
encoding. We deliberately avoid Mathlib to keep the dependency footprint at
zero and to force you to understand every axiom rather than importing it from
a library. Where a Mathlib equivalent exists we point it out in a margin
note.

\medskip
\textbf{How to read this book.}
Part~I builds the algebraic foundations: relations, partial orders, and
lattices. Part~II shows how these structures give rise to a beautiful
computational model --- the propagator network. Part~III contains two
capstone projects that demonstrate the full power of the model: a Sudoku
solver driven entirely by constraint propagation, and a type-inference
engine for a small lambda calculus. By the time you reach the end of
Part~III the connection between the mathematics and the code should feel
inevitable.

\medskip
\textbf{A note on Lean Unicode.}
All code listings use the Unicode characters that Lean~4 accepts natively:
\texttt{∀ ∃ → ← ↔ ≤ ≥ ⊓ ⊔ ⊥ ⊤ ∈ ∉ α β γ ×} and so on. When
reading the code, treat these as single characters exactly as Lean treats
them. In Lean~4 you enter them with backslash sequences (e.g.\ \texttt{\textbackslash{}le} for \texttt{≤}); your editor should substitute them live.

\tableofcontents

% ============================================================
\mainmatter

% ============================================================
\part{The Algebra of Order}

% ============================================================
\chapter{Relations and Partial Orders}

\begin{quote}
  \itshape
  ``The concept of an ordering is one of the most basic in all of
  mathematics, and yet we rarely pause to examine what it actually requires.''
\end{quote}

Before we can study lattices or build propagator networks, we need a precise
vocabulary for talking about \emph{order}: what it means for one thing to be
``at most'' another, how such a concept can be captured in a type, and which
algebraic laws make it useful.

% ── 1.1 Binary Relations ────────────────────────────────────
\section{Binary Relations}

\begin{defbox}[Binary Relation]
  A \emph{binary relation} on a set $A$ is a subset $R \subseteq A \times A$.
  We write $a \mathrel{R} b$ (or $R(a, b)$) as shorthand for $(a, b) \in R$.
\end{defbox}

Lean~4 encodes a relation as a function returning \lc{Prop}:

\begin{minted}{lean4}
-- A binary relation on type α
def Rel (α : Type) : Type := α → α → Prop
\end{minted}

Using \lc{Prop} rather than \lc{Bool} is crucial: proofs are computationally
irrelevant (two proofs of the same proposition are definitionally equal), and
we can state facts like ``this relation is a partial order'' as types.

\begin{exbox}[Common relations]
  \begin{itemize}[nosep]
    \item $\leq$ on $\mathbb{N}$: $R = \{(m, n) \mid m \leq n\}$.
    \item Divisibility on $\mathbb{N}$: $m \mid n \iff \exists k,\; n = m \cdot k$.
    \item Subset inclusion on $\mathcal{P}(S)$: $A \subseteq B \iff \forall x, x \in A \to x \in B$.
    \item The equality relation $=$ on any set.
    \item The \emph{discrete} relation: $a \mathrel{R} b \iff a = b$.
    \item The \emph{trivial} relation: $a \mathrel{R} b$ for all $a, b$ (everything relates to everything).
  \end{itemize}
\end{exbox}

% ── 1.2 Properties ──────────────────────────────────────────
\section{Properties of Relations}

Four properties will occupy us throughout the book.

\begin{defbox}[Reflexivity, Symmetry, Antisymmetry, Transitivity]
  Let $R$ be a binary relation on $A$.
  \begin{align*}
    \text{Reflexive}    &\;:\iff\; \forall a \in A,\; a \mathrel{R} a \\
    \text{Symmetric}    &\;:\iff\; \forall a, b \in A,\; a \mathrel{R} b \to b \mathrel{R} a \\
    \text{Antisymmetric}&\;:\iff\; \forall a, b \in A,\; a \mathrel{R} b \to b \mathrel{R} a \to a = b \\
    \text{Transitive}   &\;:\iff\; \forall a, b, c \in A,\; a \mathrel{R} b \to b \mathrel{R} c \to a \mathrel{R} c
  \end{align*}
\end{defbox}

In Lean~4:

\begin{minted}{lean4}
namespace OrderBook

def Reflexive  {α : Type} (r : Rel α) : Prop :=
  ∀ (a : α), r a a

def Symmetric  {α : Type} (r : Rel α) : Prop :=
  ∀ (a b : α), r a b → r b a

def Antisymm   {α : Type} (r : Rel α) : Prop :=
  ∀ (a b : α), r a b → r b a → a = b

def Transitive {α : Type} (r : Rel α) : Prop :=
  ∀ (a b c : α), r a b → r b c → r a c

end OrderBook
\end{minted}

\begin{notebox}
  \textbf{Symmetry vs Antisymmetry.}
  These are not negations of each other! A relation can be both symmetric and
  antisymmetric (equality), neither (``$a$ is a parent of $b$''), or only one.
  Antisymmetry says: if you can go \emph{both} ways, the elements must be equal.
\end{notebox}

% ── 1.3 Partial Orders ──────────────────────────────────────
\section{Partial Orders}

\begin{defbox}[Partial Order]
  A \emph{partial order} on $A$ is a relation $\leq$ that is
  reflexive, antisymmetric, and transitive.
  The pair $(A, \leq)$ is called a \emph{partially ordered set} (or
  \emph{poset}).
\end{defbox}

The word \emph{partial} refers to the fact that not every two elements need be
comparable: there may exist $a, b$ with neither $a \leq b$ nor $b \leq a$.

In Lean~4 we capture this as a typeclass:

\begin{minted}{lean4}
namespace OrderBook

-- Our own partial order typeclass (mirrors core Lean / Mathlib).
class PartialOrder (α : Type) where
  le         : α → α → Prop
  le_refl    : ∀ (a : α), le a a
  le_antisymm: ∀ (a b : α), le a b → le b a → a = b
  le_trans   : ∀ (a b c : α), le a b → le b c → le a c

-- Attach the ≤ notation to any PartialOrder instance.
instance [PartialOrder α] : LE α := ⟨PartialOrder.le⟩

-- A convenience alias for strict less-than
def lt [PartialOrder α] (a b : α) : Prop :=
  PartialOrder.le a b ∧ ¬ (a = b)

instance [PartialOrder α] : LT α := ⟨lt⟩

end OrderBook
\end{minted}

\begin{notebox}
  \textbf{A naming convention.} Definitions like \lc{PartialOrder} live in
  our \lc{OrderBook} namespace so they cannot collide with core Lean's.
  Listings that sit outside an explicit \lc{namespace OrderBook … end}
  block assume \lc{open OrderBook} is in effect (occasionally written
  explicitly as \lc{open OrderBook in} before a single declaration), which
  lets us write \lc{PartialOrder.le\_antisymm} instead of the fully
  qualified \lc{OrderBook.PartialOrder.le\_antisymm}.
\end{notebox}

\subsection{The Natural Numbers}

The natural numbers with the usual $\leq$ form our first instance.
In Lean~4's core library, \lc{Nat.le} is defined inductively:

\begin{minted}{lean4}
-- Recap: how ≤ is defined in core Lean 4 (no Mathlib needed)
-- inductive Nat.le : Nat → Nat → Prop
--   | refl : Nat.le n n
--   | step : Nat.le n m → Nat.le n (m + 1)

instance : OrderBook.PartialOrder Nat where
  le          := (· ≤ ·)      -- uses Nat.le which is in core
  le_refl     := Nat.le_refl
  le_antisymm := fun a b h1 h2 => Nat.le_antisymm h1 h2
  le_trans    := fun a b c h1 h2 => Nat.le_trans h1 h2
\end{minted}

\begin{notebox}
  Lean~4's \emph{core} library already provides \lc{Nat.le_refl},
  \lc{Nat.le_antisymm}, and \lc{Nat.le_trans}. We reuse these rather than
  reprove them. In later chapters we will prove non-trivial instances by hand.
\end{notebox}

\subsection{Divisibility}

Divisibility is a partial order on $\mathbb{N}$: $n \leq m$ means ``$n$ divides
$m$'', i.e.\ $\exists k, m = n \cdot k$. This is a genuinely different order from
$\leq$: for instance $3 \leq 12$ (since $12 = 3 \cdot 4$) but $4 \not\leq 6$
(no natural number $k$ satisfies $6 = 4k$).

Core Lean already attaches an \lc{LE Nat} instance for the usual order, and
adding a second one for divisibility would make every \lc{≤} on \lc{Nat}
ambiguous. So we introduce a newtype wrapper and give \emph{it} the
divisibility order.

\begin{notebox}
  \textbf{Why \lc{structure} and not \lc{def DvdNat := Nat}?}
  Using \lc{def} creates a \emph{transparent alias}: Lean can and will see
  through it, meaning it occasionally treats \lc{DvdNat} as \lc{Nat} and
  tries to apply \lc{Nat}'s existing \lc{LE} or \lc{HMul} instances, causing
  ``typeclass instance stuck'' errors. A \lc{structure} creates a
  \emph{nominally distinct} type — Lean never conflates \lc{DvdNat} with
  \lc{Nat} unless you explicitly coerce. This is the reliable choice whenever
  you want a newtype for instance purposes.
\end{notebox}

\begin{minted}{lean4}
-- A nominally distinct wrapper type for Nat under divisibility
structure DvdNat where
  val : Nat

-- Coercion from DvdNat to Nat (one direction is enough for our proofs)
instance : Coe DvdNat Nat := ⟨DvdNat.val⟩

-- Numeric literals: (5 : DvdNat) desugars to DvdNat.mk 5
instance {n : Nat} : OfNat DvdNat n := ⟨⟨n⟩⟩

-- The divisibility relation, stated entirely in Nat to avoid ambiguity
-- in the multiplication on the right-hand side.
def Dvd' (n m : DvdNat) : Prop := ∃ k : Nat, m.val = n.val * k

instance : LE DvdNat := ⟨Dvd'⟩
\end{minted}

By using \lc{.val} explicitly rather than a coercion inside \lc{Dvd'}, we
ensure the multiplication \lc{n.val * k} is unambiguously \lc{Nat} multiplication.
Coercions are useful for moving single values across a boundary, but inside an
expression with multiple interacting terms they give Lean too many degrees of
freedom and can produce stuck typeclass searches.

\subsubsection{Building up the proof: the lemmas we need}

The reflexivity and transitivity proofs are short. Antisymmetry --- ``if $a
\mid b$ and $b \mid a$ then $a = b$'' --- is where the real mathematical content
lives. Rather than writing it as a single inscrutable block, we build it from
three small lemmas, each of which is worth understanding on its own.

\begin{notebox}
  \textbf{First contact with tactics.}
  The proofs below are our first \emph{tactic proofs}: everything after
  \lc{by} is a sequence of commands (\lc{obtain}, \lc{simp}, \lc{exact},
  \ldots) that transform the goal step by step. Every tactic and
  programming construct used anywhere in this book has an entry --- with a
  small worked example, and where illuminating, the term it desugars to ---
  in the primer of Appendix~\ref{app:quickref}. Keep it bookmarked; we will
  not re-explain a tactic after its first appearance.
\end{notebox}

\medskip
\textbf{Lemma 1: if $a = 0$ and $a \mid b$, then $b = 0$.}

If $a = 0$, then $b = 0 \cdot k = 0$ for any $k$. This is the zero case of
antisymmetry: the only thing that $0$ divides is $0$ itself.

\begin{minted}{lean4}
-- If 0 | m, then m.val = 0.
-- Proof: the witness k gives m.val = 0 * k.
-- We then reduce 0 * k to 0 using Nat.zero_mul.
theorem dvd_zero_left (m : DvdNat) (h : Dvd' ⟨0⟩ m) : m.val = 0 := by
  obtain ⟨k, hk⟩ := h
  -- hk : m.val = DvdNat.mk(0).val * k
  -- DvdNat.mk(0).val reduces to 0 by definition, so simp can see it
  simp at hk
  -- hk is now: m.val = 0
  exact hk
\end{minted}

\begin{notebox}
  \textbf{Why \lc{simp} instead of \lc{rw [Nat.zero\_mul]}?}
  With a \lc{structure}, the literal \lc{0 : DvdNat} is
  \lc{OfNat.ofNat 0 = DvdNat.mk 0}, so its \lc{.val} is
  \lc{DvdNat.mk(0).val}, not the bare \lc{0} that \lc{Nat.zero\_mul} expects
  as its left-hand side pattern. \lc{simp} performs the definitional unfolding
  first and then applies \lc{Nat.zero\_mul}, where a bare \lc{rw} cannot.
  If you want to avoid \lc{simp} you can write \lc{show m.val = 0} after an
  explicit \lc{have : DvdNat.mk(0).val = 0 := rfl}.
\end{notebox}

\medskip
\textbf{Lemma 2: if $k \cdot j = 1$ in $\mathbb{N}$, then $k = 1$.}

This is the key algebraic fact that drives the nonzero case. In $\mathbb{Z}$ the
units are $\pm 1$; in $\mathbb{N}$ the only unit is $1$. We prove it by
exhaustive case analysis: $k$ cannot be $0$ (product would be $0$) and cannot
be $\geq 2$ (product would be $\geq 2$).

\begin{minted}{lean4}
-- In ℕ, k * j = 1 forces k = 1 (and by symmetry, j = 1).
-- Proved by explicit case analysis on k; omega only closes the final
-- numeric contradiction. NB: the calc chain is written in the ≤
-- direction — core Lean ships a Trans instance for ≤ but not for ≥.
theorem mul_eq_one_left {k j : Nat} (h : k * j = 1) : k = 1 := by
  match k with
  | 0 =>
    -- 0 * j = 0 ≠ 1, so this case is impossible.
    -- simp knows 0 * j = 0, which contradicts h : 0 = 1.
    simp [Nat.zero_mul] at h
  | 1 =>
    -- k is already 1; no further argument needed.
    rfl
  | k + 2 =>
    -- (k+2) * j ≥ 2, contradicting h : (k+2) * j = 1.
    exfalso
    -- First establish j ≥ 1 (if j = 0, the product is 0, not 1).
    have hj : j ≥ 1 := by
      cases j with
      | zero   => simp [Nat.mul_zero] at h
      | succ j => exact Nat.succ_le_succ (Nat.zero_le j)
    -- Now show the product is at least 2.
    have hbig : 2 ≤ (k + 2) * j :=
      calc 2 ≤ k + 2         := Nat.le_add_left 2 k
        _ = (k + 2) * 1      := (Nat.mul_one (k + 2)).symm
        _ ≤ (k + 2) * j      := Nat.mul_le_mul_left (k + 2) hj
    -- h says the product is 1, but hbig says it is ≥ 2. Contradiction.
    omega

-- The right-factor version follows by commutativity.
theorem mul_eq_one_right {k j : Nat} (h : k * j = 1) : j = 1 :=
  mul_eq_one_left (by rw [Nat.mul_comm]; exact h)
\end{minted}

\medskip
\textbf{Lemma 3: if $a > 0$, $a \mid b$, and $b \mid a$, then $a = b$.}

Given witnesses $k$ and $j$ with $b = a \cdot k$ and $a = b \cdot j$, we
substitute to get $a = a \cdot (k \cdot j)$, cancel $a$ (valid because $a >
0$), apply Lemma~2, and conclude $k = 1$, hence $b = a$.

\begin{minted}{lean4}
-- All arithmetic here lives in Nat — no DvdNat involved.
theorem dvd_antisymm_pos
    {a b : Nat} (ha : 0 < a)
    (k : Nat) (hk : b = a * k)
    (j : Nat) (hj : a = b * j) : a = b := by
  -- Step 1: derive k * j = 1.
  -- Substitute hk into hj to get  a = a * (k * j),
  -- then cancel a (valid since a > 0).
  have hkj : k * j = 1 := by
    have hstep : a = a * (k * j) :=
      calc a = b * j       := hj
           _ = (a * k) * j := by rw [hk]
           _ = a * (k * j) := by rw [Nat.mul_assoc]
    exact Nat.eq_of_mul_eq_mul_left ha (by rw [Nat.mul_one]; exact hstep.symm)
  -- Step 2: j = 1 by Lemma 2 (right factor version).
  have hj1 : j = 1 := mul_eq_one_right hkj
  -- Step 3: substitute j = 1 into hj : a = b * j to get a = b * 1 = b.
  have h₂ : a = b * 1 := by rw [hj1] at hj; exact hj
  rw [Nat.mul_one] at h₂
  exact h₂
\end{minted}

\begin{notebox}
  \textbf{Why \lc{Nat.eq\_of\_mul\_eq\_mul\_left}?}
  This says: if $a > 0$ and $a \cdot x = a \cdot y$, then $x = y$.
  It is multiplication cancellation for $\mathbb{N}$, available in core Lean.
  The positivity hypothesis is essential: $0 \cdot x = 0 \cdot y$ for all
  $x, y$, so cancellation is unsound at zero.
\end{notebox}

\subsubsection{Assembling the instance}

With the lemmas in hand, each field is a short delegation.

\begin{minted}{lean4}
open OrderBook in
instance : PartialOrder DvdNat where

  -- The order itself: divisibility.
  le := Dvd'

  -- Reflexivity: n | n via witness k = 1, since n.val = n.val * 1.
  le_refl := fun n =>
    ⟨1, (Nat.mul_one n.val).symm⟩

  -- Antisymmetry: expose the divisibility witnesses, then case-split
  -- on whether n.val is zero or positive, delegating to Lemma 1 or
  -- Lemma 3 respectively.
  le_antisymm := fun n m hnm hmn => by
    obtain ⟨k, hk⟩ := hnm
    obtain ⟨j, hj⟩ := hmn
    -- Reduce  n = m  (DvdNat equality) to  a = b  on the .val fields:
    -- cases exposes the fields, congr 1 splits the constructor.
    cases n; cases m
    congr 1
    rename_i a b
    cases Nat.eq_zero_or_pos a with
    | inl ha =>
      -- a = 0: b = 0 * k = 0 by Lemma 1, and a = 0, so a = b.
      subst ha
      exact (dvd_zero_left ⟨b⟩ ⟨k, hk⟩).symm
    | inr ha =>
      -- a > 0: use Lemma 3 directly.
      exact dvd_antisymm_pos ha k hk j hj

  -- Transitivity: compose witnesses.  b.val = a.val * k and
  -- c.val = b.val * j, so c.val = a.val * (k * j).
  le_trans := fun a b c ⟨k, hk⟩ ⟨j, hj⟩ =>
    ⟨k * j, by rw [hj, hk, Nat.mul_assoc]⟩
\end{minted}

\begin{notebox}
  \textbf{Three details worth noticing.}
  (1) Our \lc{PartialOrder} class bundles the order itself as the \lc{le}
  field, so the instance must supply it even though we also declared an
  \lc{LE DvdNat} instance above --- that one only provides the \lc{≤}
  notation.
  (2) We destructure the hypotheses with \lc{obtain} \emph{inside} the
  tactic block rather than in the lambda binders: pattern-matching binders
  like \lc{fun n m ⟨k, hk⟩ ⟨j, hj⟩ =>} would leave the witnesses
  inaccessible to \lc{rename\_i} after \lc{cases n; cases m}.
  (3) \lc{dvd\_antisymm\_pos} already concludes \lc{a = b}, exactly the
  goal, so no \lc{.symm} is needed in the positive case.
\end{notebox}

\subsubsection{Reflection}

Notice that the structure of the proof exactly mirrors the structure of the
mathematics:
\begin{itemize}[nosep]
  \item Reflexivity is trivial because the witness $k = 1$ is explicit.
  \item Transitivity is trivial because multiplying witnesses gives a new
    witness: $b = a \cdot k$ and $c = b \cdot j$ compose to $c = a \cdot (kj)$.
  \item Antisymmetry requires real work because it cannot be witnessed
    directly --- it requires us to reason about what witnesses can possibly
    exist and conclude they must all be $1$.
\end{itemize}

The lemma \lc{mul\_eq\_one\_left} is the mathematical core. Everything else is
bookkeeping. When reading an unfamiliar proof, it is always worth asking: where
is the actual mathematical content? In this case it is entirely in that one
lemma.

\subsection{Sets}

Sets (represented as predicates) form a canonical partial order under inclusion.

\begin{minted}{lean4}
-- We represent sets as indicator functions
def Set (α : Type) : Type := α → Prop

-- Subset inclusion
def Set.subset {α : Type} (s t : Set α) : Prop :=
  ∀ (a : α), s a → t a

instance {α : Type} : OrderBook.PartialOrder (Set α) where
  le          := Set.subset
  le_refl     := fun s a ha => ha
  le_antisymm := fun s t h1 h2 =>
    funext (fun a => propext ⟨h1 a, h2 a⟩)
  le_trans    := fun s t u h1 h2 a ha => h2 a (h1 a ha)
\end{minted}

This is our first non-trivial Lean proof. Let us read it carefully:
\begin{itemize}
  \item \lc{le_refl}: given $a \in s$, produce $a \in s$ --- trivial.
  \item \lc{le_antisymm}: $s \subseteq t$ and $t \subseteq s$ imply $s = t$.
    We use \lc{funext} (function extensionality) and \lc{propext}
    (propositional extensionality) to conclude $s\,a \leftrightarrow t\,a$ for
    all $a$, hence $s = t$.
  \item \lc{le_trans}: chain the inclusions.
\end{itemize}

% ── 1.4 Hasse Diagrams ──────────────────────────────────────
\section{Hasse Diagrams}

Posets are best understood visually. A \emph{Hasse diagram} draws elements as
nodes and covers as upward edges: we say $a$ is \emph{covered by} $b$
(written $a \lessdot b$) if $a < b$ and there is no $c$ with $a < c < b$.

\begin{figure}[h]
  \centering
  \begin{tikzpicture}[node distance=1.4cm]
    % Powerset of {a,b,c}
    \node[node] (bot)  {\small $\varnothing$};
    \node[node] (a)    [above left=1.2cm and 1.6cm of bot] {\small $\{a\}$};
    \node[node] (b)    [above=1.8cm of bot]                {\small $\{b\}$};
    \node[node] (c)    [above right=1.2cm and 1.6cm of bot]{\small $\{c\}$};
    \node[node] (ab)   [above=1.8cm of a]                  {\small $\{a,b\}$};
    \node[node] (bc)   [above=1.8cm of c]                  {\small $\{b,c\}$};
    \node[node] (ac)   [above=1.8cm of b]                  {\small $\{a,c\}$};
    \node[node] (top)  [above=1.8cm of ac]                 {\small $\{a,b,c\}$};
    \draw[hasseedge] (bot)--(a); \draw[hasseedge] (bot)--(b); \draw[hasseedge] (bot)--(c);
    \draw[hasseedge] (a)--(ab);  \draw[hasseedge] (b)--(ab); 
    \draw[hasseedge] (b)--(bc);  \draw[hasseedge] (c)--(bc);
    \draw[hasseedge] (a)--(ac);  \draw[hasseedge] (c)--(ac);
    \draw[hasseedge] (ab)--(top);\draw[hasseedge] (ac)--(top);\draw[hasseedge] (bc)--(top);
  \end{tikzpicture}
  \caption{The powerset lattice $\mathcal{P}(\{a,b,c\})$ ordered by $\subseteq$.}
\end{figure}

\begin{figure}[h]
  \centering
  \begin{tikzpicture}[node distance=1.4cm]
    % Divisibility poset on {1,2,3,4,6,12}
    \node[node] (1)  {1};
    \node[node] (2)  [above left=1.4cm and 0.9cm of 1]  {2};
    \node[node] (3)  [above right=1.4cm and 0.9cm of 1] {3};
    \node[node] (4)  [above=1.4cm of 2]                  {4};
    \node[node] (6)  [above right=1.4cm and 0.0cm of 2]  {6};
    \node[node] (12) [above right=1.4cm and 0.0cm of 4]  {12};
    \draw[hasseedge] (1)--(2); \draw[hasseedge] (1)--(3);
    \draw[hasseedge] (2)--(4); \draw[hasseedge] (2)--(6);
    \draw[hasseedge] (3)--(6);
    \draw[hasseedge] (4)--(12);\draw[hasseedge] (6)--(12);
  \end{tikzpicture}
  \caption{The divisibility poset on $\{1, 2, 3, 4, 6, 12\}$.}
\end{figure}

% ── 1.4 Computation interlude ────────────────────────────────
\FloatBarrier
\section{Computation Interlude: Running the Divisibility Order}

So far every definition has lived in \lc{Prop}. This section bridges that gap
with a fully computable divisibility order on $\{1,\ldots,12\}$.

\begin{minted}{lean4}
def D12 : List Nat := [1, 2, 3, 4, 6, 12]
def dvd12 (a b : Nat) : Bool := b % a == 0

-- Supremum: smallest upper bound
def upperBounds (xs : List Nat) : List Nat :=
  D12.filter (fun u => xs.all (fun x => dvd12 x u))

def sup12 (xs : List Nat) : Option Nat :=
  let ubs := upperBounds xs
  ubs.find? (fun s => ubs.all (fun u => dvd12 s u))

#eval sup12 [4, 6]   -- some 12  (lcm)
#eval sup12 [2, 3]   -- some 6

-- Infimum: greatest lower bound
def lowerBounds (xs : List Nat) : List Nat :=
  D12.filter (fun l => xs.all (fun x => dvd12 l x))

def inf12 (xs : List Nat) : Option Nat :=
  let lbs := lowerBounds xs
  lbs.find? (fun s => lbs.all (fun l => dvd12 l s))

#eval inf12 [4, 6]   -- some 2  (gcd)
#eval inf12 [3, 4]   -- some 1

-- Print Hasse diagram edges
def covers (a b : Nat) : Bool :=
  dvd12 a b && a != b &&
  !D12.any (fun c => c != a && c != b && dvd12 a c && dvd12 c b)

def printHasse : IO Unit := do
  for a in D12 do
    for b in D12 do
      if covers a b then IO.println s!"{a} ──► {b}"

#eval printHasse  -- 1→2, 1→3, 2→4, 2→6, 3→6, 4→12, 6→12
\end{minted}

% ── 1.5 Exercises ───────────────────────────────────────────
\section{Exercises}

\begin{exercisebox}[Reflexivity and symmetry]
  Show that a relation that is both symmetric and antisymmetric must be a
  subset of the equality relation, i.e.\ $a \mathrel{R} b \to a = b$.
  Prove this in Lean~4.
  \begin{minted}{lean4}
theorem symm_antisymm_eq {α : Type} (r : Rel α)
    (hs : Symmetric r) (ha : Antisymm r) :
    ∀ a b : α, r a b → a = b := by
  intro a b h
  exact ha a b h (hs a b h)
  \end{minted}
\end{exercisebox}

\begin{exercisebox}[Divisibility on a small set]
  Consider the divisibility relation restricted to the set $\{1,2,3,4,6,12\}$.
  \begin{enumerate}[label=(\alph*), nosep]
    \item Draw its Hasse diagram.
    \item Which pairs of elements are incomparable?
    \item In Lean~4, define a type \lc{D12} with six constructors and prove
          that divisibility on it satisfies \lc{le\_refl}.
  \end{enumerate}
\end{exercisebox}

\begin{exercisebox}[Product order]
  Given two partial orders $(A, \leq_A)$ and $(B, \leq_B)$, define the
  \emph{product order} on $A \times B$ by
  $(a_1, b_1) \leq (a_2, b_2) \iff a_1 \leq_A a_2 \;\wedge\; b_1 \leq_B b_2$.
  Prove in Lean~4 that this is indeed a partial order by implementing
  the typeclass instance.
  \begin{minted}{lean4}
instance [OrderBook.PartialOrder α] [OrderBook.PartialOrder β] :
    OrderBook.PartialOrder (α × β) where
  le p q := PartialOrder.le p.1 q.1 ∧ PartialOrder.le p.2 q.2
  le_refl := fun ⟨a, b⟩ => ⟨PartialOrder.le_refl a,
                              PartialOrder.le_refl b⟩
  le_antisymm := fun ⟨a1,b1⟩ ⟨a2,b2⟩ ⟨h1,h2⟩ ⟨h3,h4⟩ => by
    have := PartialOrder.le_antisymm _ _ h1 h3
    have := PartialOrder.le_antisymm _ _ h2 h4
    simp_all
  le_trans := fun ⟨a1,b1⟩ ⟨a2,b2⟩ ⟨a3,b3⟩ ⟨h1,h2⟩ ⟨h3,h4⟩ =>
    ⟨PartialOrder.le_trans _ _ _ h1 h3,
     PartialOrder.le_trans _ _ _ h2 h4⟩
  \end{minted}
\end{exercisebox}

\begin{exercisebox}[Reverse order]
  If $(A, \leq)$ is a partial order, then so is $(A, \geq)$ where
  $a \geq b \iff b \leq a$.  This is called the \emph{dual} or \emph{opposite}
  order.  Implement a \lc{Dual} wrapper type and its \lc{PartialOrder} instance.
  We will use this construction crucially when encoding Sudoku domains in
  Chapter~8.
\end{exercisebox}

% ============================================================
\chapter{Special Elements and Monotone Maps}

In a poset, certain elements play privileged roles (least, greatest, bounds),
and certain functions are especially well-behaved (monotone maps).
This chapter equips us with the vocabulary for both.

\section{Extremal Elements}

\begin{defbox}[Minimum, Maximum, Minimal, Maximal]
  Let $(P, \leq)$ be a poset and $S \subseteq P$.
  \begin{itemize}[nosep]
    \item $m \in S$ is the \emph{minimum} of $S$ if $\forall s \in S,\; m \leq s$.
    \item $m \in S$ is the \emph{maximum} of $S$ if $\forall s \in S,\; s \leq m$.
    \item $m \in S$ is \emph{minimal} in $S$ if $\nexists s \in S,\; s < m$.
    \item $m \in S$ is \emph{maximal} in $S$ if $\nexists s \in S,\; m < s$.
  \end{itemize}
\end{defbox}

A minimum, if it exists, is unique (by antisymmetry). A minimal element need
not be unique. When a poset has a unique minimum we call it $\bot$ (bottom),
and a unique maximum $\top$ (top).

The distinction between minimum and minimal is subtle but important.
A \emph{minimum} beats every other element; a \emph{minimal} element
merely has no element strictly below it.

\begin{figure}[H]
  \centering
  \begin{tikzpicture}
    \begin{scope}[xshift=0cm]
      \node[font=\small\itshape, above] at (1, 3.8) {has a minimum};
      \node[node] (bot) at (1,0) {$a$};
      \node[node] (b)   at (0,1.5) {$b$};
      \node[node] (c)   at (2,1.5) {$c$};
      \node[node] (top) at (1,3)   {$d$};
      \draw[hasseedge] (bot)--(b); \draw[hasseedge] (bot)--(c);
      \draw[hasseedge] (b)--(top); \draw[hasseedge] (c)--(top);
      \node[right=0.15cm of bot, gray, font=\scriptsize] {minimum $= a$};
      \node[right=0.15cm of bot, gray, font=\scriptsize, yshift=-0.35cm]
        {(also minimal)};
    \end{scope}
    \begin{scope}[xshift=5cm]
      \node[font=\small\itshape, above] at (1, 3.8)
        {two minimals, no minimum};
      \node[node] (p) at (0,0) {$p$};
      \node[node] (q) at (2,0) {$q$};
      \node[node] (r) at (0,1.8) {$r$};
      \node[node] (s) at (2,1.8) {$s$};
      \node[node] (t) at (1,3)  {$t$};
      \draw[hasseedge] (p)--(r); \draw[hasseedge] (p)--(s);
      \draw[hasseedge] (q)--(s);
      \draw[hasseedge] (r)--(t); \draw[hasseedge] (s)--(t);
      \node[left=0.1cm of p, gray, font=\scriptsize] {minimal};
      \node[right=0.1cm of q, gray, font=\scriptsize] {minimal};
      \node[gray, font=\scriptsize] at (1,-0.7)
        {$p \not\leq q$ and $q \not\leq p$};
      \node[gray, font=\scriptsize] at (1,-1.1)
        {so neither is a minimum};
    \end{scope}
  \end{tikzpicture}
  \caption{Left: a poset with a unique minimum $a$ (every element is above $a$).
           Right: a poset where $p$ and $q$ are both minimal (nothing is below
           either) but there is no minimum (neither $p \leq q$ nor $q \leq p$).}
\end{figure}

\begin{minted}{lean4}
namespace OrderBook

-- Bottom and top elements
class Bot (α : Type) where bot : α
class Top (α : Type) where top : α

notation "⊥" => Bot.bot
notation "⊤" => Top.top

-- Being a bounded partial order
class BoundedPartialOrder (α : Type) extends PartialOrder α, Bot α, Top α where
  bot_le : ∀ (a : α), (⊥ : α) ≤ a
  le_top : ∀ (a : α), a ≤ (⊤ : α)

end OrderBook
\end{minted}

\begin{exbox}[Bounded orders]
  \begin{itemize}[nosep]
    \item $(\mathcal{P}(S), \subseteq)$: $\bot = \varnothing$, $\top = S$.
    \item $(\{0,1\}, \leq)$: $\bot = 0$, $\top = 1$.
    \item $(\mathbb{N}, \leq)$: $\bot = 0$, no $\top$ (unbounded above).
    \item The divisibility poset on $\{1,\ldots,12\}$: $\bot = 1$ (divides everything), no single $\top$.
  \end{itemize}
\end{exbox}

\section{Upper and Lower Bounds}

\begin{defbox}[Bounds]
  Let $(P, \leq)$ be a poset and $S \subseteq P$.
  \begin{itemize}[nosep]
    \item $u \in P$ is an \emph{upper bound} of $S$ if $\forall s \in S,\; s \leq u$.
    \item $l \in P$ is a \emph{lower bound} of $S$ if $\forall s \in S,\; l \leq s$.
    \item The \emph{supremum} (or \emph{least upper bound}) of $S$, written $\sup S$, is the
          smallest upper bound of $S$ when it exists.
    \item The \emph{infimum} (or \emph{greatest lower bound}) of $S$, written $\inf S$, is the
          largest lower bound of $S$ when it exists.
  \end{itemize}
\end{defbox}

\begin{figure}[H]
  \centering
  \begin{tikzpicture}
    \node[node] (a)   at (0,0)   {$a$};
    \node[node] (b)   at (2,1)   {$b$};
    \node[node,fill=blue!15,draw=blue!60!black] (c) at (0,2) {$c$};
    \node[node,fill=blue!15,draw=blue!60!black] (d) at (2,2) {$d$};
    \node[node] (e)   at (0,4)   {$e$};
    \node[node] (f)   at (2,4)   {$f$};
    \node[node] (top) at (1,5.5) {$g$};
    \draw[hasseedge] (a)--(b);
    \draw[hasseedge] (b)--(c);
    \draw[hasseedge] (b)--(d);
    \draw[hasseedge] (c)--(e); \draw[hasseedge] (d)--(e);
    \draw[hasseedge] (c)--(f); \draw[hasseedge] (d)--(f);
    \draw[hasseedge] (e)--(top);\draw[hasseedge] (f)--(top);
    \node[blue!70!black, font=\small\bfseries] at (-1.1, 2)
      {$S = \{c,d\}$};
    \node[gray, font=\scriptsize, right=2.5cm of c, yshift=0.5cm]
      (ann1) {$e, f, g$ are upper bounds of $S$};
    \node[gray, font=\scriptsize, below=0.05cm of ann1]
      {$\sup S = e$? Only if $e \leq f$.};
    \node[gray, font=\scriptsize, below=0.05cm of ann1, yshift=-0.4cm]
      {Here $e$ and $f$ incomparable $\Rightarrow$};
    \node[gray, font=\scriptsize, below=0.05cm of ann1, yshift=-0.75cm]
      {$\sup S = g$ (least upper bound).};
    \node[gray, font=\scriptsize, right=2.5cm of b, yshift=-0.2cm]
      {$a, b$ are lower bounds of $S$};
    \node[gray, font=\scriptsize, right=2.5cm of b, yshift=-0.6cm]
      {$\inf S = b$ (greatest lower bound).};
  \end{tikzpicture}
  \caption{Illustration of upper bounds, lower bounds, supremum and infimum
           for the subset $S = \{c, d\}$ (shaded blue) in a poset.
           The upper bounds of $S$ are all elements $\geq$ both $c$ and $d$:
           here $e$, $f$, and $g$. The \emph{supremum} is the \emph{least}
           of these; since $e$ and $f$ are incomparable, the least upper
           bound is $g$. The lower bounds are $a$ and $b$; the infimum
           is $b$ since $a \leq b$.}
\end{figure}

\begin{thmbox}[Uniqueness of sup/inf]
  In a poset, if $\sup S$ exists, it is unique. Similarly for $\inf S$.
\end{thmbox}
\begin{proof}
  Suppose $u$ and $v$ are both least upper bounds of $S$. Since $u$ is an
  upper bound and $v$ is a \emph{least} upper bound, $v \leq u$. By symmetry
  $u \leq v$. Antisymmetry gives $u = v$.
\end{proof}

\begin{minted}{lean4}
-- The supremum of two elements (binary version)
-- (We'll generalise to arbitrary subsets in Chapter 4.)
structure IsSupOf {α : Type} [OrderBook.PartialOrder α]
    (a b sup : α) : Prop where
  ge_a  : a ≤ sup
  ge_b  : b ≤ sup
  least : ∀ (u : α), a ≤ u → b ≤ u → sup ≤ u

-- Uniqueness of the binary sup
theorem sup_unique {α : Type} [OrderBook.PartialOrder α]
    {a b s t : α}
    (hs : IsSupOf a b s) (ht : IsSupOf a b t) : s = t := by
  apply PartialOrder.le_antisymm
  · exact hs.least t ht.ge_a ht.ge_b
  · exact ht.least s hs.ge_a hs.ge_b
\end{minted}

\section{Monotone Maps}

\begin{defbox}[Monotone Map]
  A function $f : (P, \leq_P) \to (Q, \leq_Q)$ between posets is
  \emph{monotone} (or \emph{order-preserving}) if
  $\forall a, b \in P,\; a \leq_P b \to f(a) \leq_Q f(b)$.
\end{defbox}

Monotone maps are the morphisms of the category of posets. They will be
central to propagators: every propagator is a monotone function on a lattice.

\begin{figure}[H]
  \centering
  \begin{tikzpicture}
    \begin{scope}[xshift=0cm]
      \node[font=\small\bfseries] at (1, 4.2) {$P$};
      \node[node] (p1) at (0,0) {$a$};
      \node[node] (p2) at (2,0) {$b$};
      \node[node] (p3) at (1,2) {$c$};
      \node[node] (p4) at (1,3.5) {$d$};
      \draw[hasseedge] (p1)--(p3); \draw[hasseedge] (p2)--(p3);
      \draw[hasseedge] (p3)--(p4);
    \end{scope}
    \draw[-{Stealth},thick,blue!60!black] (3.2,1.8)
      -- node[above,font=\small,blue!60!black]{$f$ (monotone)} (4.8,1.8);
    \begin{scope}[xshift=6cm]
      \node[font=\small\bfseries] at (1, 4.2) {$Q$};
      \node[node,minimum size=10mm,font=\tiny] (q1) at (1,0)
        {$f(a){=}f(b)$};
      \node[node] (q3) at (1,2)   {$f(c)$};
      \node[node] (q4) at (1,3.5) {$f(d)$};
      \draw[hasseedge] (q1)--(q3); \draw[hasseedge] (q3)--(q4);
      \node[gray,font=\scriptsize,right=0.15cm of q1,yshift=0.7cm]
        {collapsing $a,b$ is fine $\checkmark$};
    \end{scope}
    \draw[dashed,gray!50] (p1) -- (q1);
    \draw[dashed,gray!50] (p2) -- (q1);
    \draw[dashed,gray!50] (p3) -- (q3);
    \draw[dashed,gray!50] (p4) -- (q4);
    \begin{scope}[yshift=-3.5cm]
      \begin{scope}[xshift=0cm]
        \node[font=\small\bfseries] at (1, 2.7) {$P$};
        \node[node] (r1) at (0,0) {$x$};
        \node[node] (r2) at (1,2) {$y$};
        \draw[hasseedge] (r1)--(r2);
        \node[gray,font=\scriptsize,below=0.05cm of r1]{$x \leq y$};
      \end{scope}
      \draw[-{Stealth},thick,red!60!black] (3.2,1)
        -- node[above,font=\small,red!60!black]{$g$ (not monotone)} (4.8,1);
      \begin{scope}[xshift=6cm]
        \node[font=\small\bfseries] at (1, 2.7) {$Q$};
        \node[node] (s1) at (1,2) {$g(x)$};
        \node[node] (s2) at (1,0) {$g(y)$};
        \draw[hasseedge] (s2)--(s1);
        \node[gray,font=\scriptsize,right=0.1cm of s2,yshift=1cm]
          {\color{red!70!black}$g(x) > g(y)$: order violated!};
      \end{scope}
      \draw[dashed,gray!50] (r1) -- (s1);
      \draw[dashed,gray!50] (r2) -- (s2);
    \end{scope}
  \end{tikzpicture}
  \caption{Top: a monotone map $f$ preserves the order --- whenever $a \leq_P b$
           in the source, $f(a) \leq_Q f(b)$ in the target. Here $f$ collapses
           the incomparable elements $a$ and $b$ to the same image, which is
           permitted. Bottom: a non-monotone map $g$ where $x \leq y$ but
           $g(x) > g(y)$, violating the order.}
\end{figure}

\begin{minted}{lean4}
namespace OrderBook

structure Monotone {α β : Type}
    [PartialOrder α] [PartialOrder β] (f : α → β) : Prop where
  map_le : ∀ (a b : α), a ≤ b → f a ≤ f b

-- Composition of monotone maps is monotone
theorem Monotone.comp {α β γ : Type}
    [PartialOrder α] [PartialOrder β] [PartialOrder γ]
    {f : α → β} {g : β → γ}
    (hf : Monotone f) (hg : Monotone g) : Monotone (g ∘ f) where
  map_le := fun a b h => hg.map_le _ _ (hf.map_le _ _ h)

end OrderBook
\end{minted}

\begin{exbox}[Monotone functions]
  \begin{itemize}[nosep]
    \item $f(n) = n + 5$ is monotone on $(\mathbb{N}, \leq)$.
    \item Set union $S \mapsto S \cup \{x\}$ is monotone on $(\mathcal{P}(A), \subseteq)$.
    \item Set intersection $S \mapsto S \cap T$ is monotone on $(\mathcal{P}(A), \subseteq)$ for fixed $T$.
    \item Any constant function $f(x) = c$ is monotone.
    \item The identity function is always monotone.
  \end{itemize}
\end{exbox}

\section{Exercises}

\begin{exercisebox}[Monotone composition]
  Give an example showing that the composition of two \emph{antitone} maps
  (where $a \leq b \to f(b) \leq f(a)$) is monotone.
\end{exercisebox}

\begin{exercisebox}[Fixed points]
  A \emph{fixed point} of $f : P \to P$ is an element $x$ with $f(x) = x$.
  Show that if $f$ is monotone and $a \leq f(a)$, then the chain
  $a \leq f(a) \leq f^2(a) \leq \cdots$ is indeed a chain (all elements
  pairwise comparable with the given order).
\end{exercisebox}

\begin{exercisebox}[Characterising PartialOrder via Monotone]
  Show that the identity map $\mathrm{id} : (P,\leq) \to (P,\leq)$ is always
  monotone, and prove this in Lean~4 using your \lc{Monotone} structure.
\end{exercisebox}

\FloatBarrier
\section{Computation Interlude: Running Monotone Maps}

Building on the computable divisibility order from Chapter~1, we can now compute
with monotone maps and verify their properties at runtime.

\textbf{Checking monotonicity}

Given a function on \lc{D12}, we can decide whether it is monotone by checking
all pairs.

\begin{minted}{lean4}
-- Check all comparable pairs to see if f preserves the order
def isMonotone (f : Nat → Nat) : Bool :=
  D12.all (fun a =>
  D12.all (fun b =>
    -- if a | b then f(a) | f(b)
    !dvd12 a b || dvd12 (f a) (f b)))

-- The identity is monotone
#eval isMonotone id                         -- true

-- Multiplying by 2: if a | b then 2a | 2b
#eval isMonotone (· * 2)                   -- true

-- A non-monotone map: send every element to its "complement" 12/n
-- 1↦12, 2↦6, 3↦4, 4↦3, 6↦2, 12↦1 — this reverses the order
#eval isMonotone (fun n => 12 / n)         -- false (it's antitone)

-- But composing it with itself gives monotone again
#eval isMonotone (fun n => 12 / (12 / n)) -- true
\end{minted}

\textbf{Computing fixed points}

A fixed point of $f$ is an element where $f(x) = x$. We can find them all.

\begin{minted}{lean4}
def fixedPoints (f : Nat → Nat) : List Nat :=
  D12.filter (fun x => f x == x)

#eval fixedPoints id                    -- [1, 2, 3, 4, 6, 12]
#eval fixedPoints (fun _ => 1)          -- [1]
#eval D12.filter (fun n => 12 / n == n) -- [] (no fixed points for n ↦ 12/n)
\end{minted}

\textbf{Iterating a monotone map}

From Exercise~2.2, we know that if $f$ is monotone and $a \leq f(a)$, the chain
$a, f(a), f^2(a), \ldots$ is ascending. We can watch it converge.

\begin{minted}{lean4}
-- Apply f repeatedly until fixpoint or step limit
def iterToFixpoint (f : Nat → Nat) (start : Nat) (limit : Nat) : List Nat :=
  let rec go (x : Nat) (n : Nat) (acc : List Nat) : List Nat :=
    if n = 0 then acc.reverse
    else let y := f x
         if y == x then (acc.reverse ++ [x])
         else go y (n - 1) (x :: acc)
  go start limit []

-- f(n) = if 2*n ∈ D12 then 2*n else n
def double12 (n : Nat) : Nat :=
  if D12.contains (n * 2) then n * 2 else n

#eval iterToFixpoint double12 1 10   -- [1, 2, 4] ascending chain!
#eval iterToFixpoint double12 3 10   -- [3, 6, 12]
#eval iterToFixpoint double12 12 10  -- [12] already a fixed point
\end{minted}

The output directly illustrates the ascending chain idea: each application of
\lc{double12} moves strictly upward in the divisibility order until it hits a
fixed point. Note \emph{which} fixed point: starting from 1 the chain stops at
4, because $8 \notin D_{12}$ makes 4 a fixed point of \lc{double12} --- a
monotone map can have many fixed points, and iteration finds the one above
where you started. Starting from 3 instead, the chain climbs all the way to
the top element 12.

% ============================================================
\chapter{Lattices}

A poset in which every pair of elements has both a supremum and an infimum
is called a \emph{lattice}. This apparently modest requirement turns out to
yield a rich algebraic structure that shows up in logic, topology, computer
science, and --- as we will see in Part~II --- concurrency and constraint solving.

\section{Meet and Join}

\begin{defbox}[Lattice]
  A \emph{lattice} is a poset $(L, \leq)$ in which every two elements
  $a, b \in L$ have:
  \begin{itemize}[nosep]
    \item a \emph{meet} (greatest lower bound) $a \sqcap b$, and
    \item a \emph{join} (least upper bound) $a \sqcup b$.
  \end{itemize}
\end{defbox}

\begin{figure}[H]
  \centering
  \begin{tikzpicture}
    \node[node] (bot) at (1,0)   {$a \sqcap b$};
    \node[node] (a)   at (0,1.8) {$a$};
    \node[node] (b)   at (2,1.8) {$b$};
    \node[node] (top) at (1,3.6) {$a \sqcup b$};
    \draw[hasseedge] (bot)--(a); \draw[hasseedge] (bot)--(b);
    \draw[hasseedge] (a)--(top); \draw[hasseedge] (b)--(top);
    \node[right=0.2cm of top, gray, font=\scriptsize, text width=3.5cm]
      {join $a \sqcup b$:\\least upper bound of $a$ and $b$};
    \node[right=0.2cm of bot, gray, font=\scriptsize, text width=3.5cm]
      {meet $a \sqcap b$:\\greatest lower bound of $a$ and $b$};
    \node[left=0.1cm of a, gray, font=\scriptsize]
      {$a \sqcap b \leq a \leq a \sqcup b$};
    \node[right=0.1cm of b, gray, font=\scriptsize]
      {$a \sqcap b \leq b \leq a \sqcup b$};
  \end{tikzpicture}
  \caption{The meet $a \sqcap b$ sits below both $a$ and $b$; the join
           $a \sqcup b$ sits above both. In a lattice every pair of elements
           has both.}
\end{figure}

The meet and join satisfy the following axioms (which can be taken as an
equivalent \emph{algebraic} definition of a lattice):

\begin{align}
  a \sqcap b &= b \sqcap a   & a \sqcup b &= b \sqcup a \tag{commutativity}\\
  (a \sqcap b) \sqcap c &= a \sqcap (b \sqcap c) &
  (a \sqcup b) \sqcup c &= a \sqcup (b \sqcup c) \tag{associativity}\\
  a \sqcap a &= a  &  a \sqcup a &= a \tag{idempotence}\\
  a \sqcap (a \sqcup b) &= a &  a \sqcup (a \sqcap b) &= a \tag{absorption}
\end{align}

The order relation can be recovered from meet or join:
$$a \leq b \iff a \sqcap b = a \iff a \sqcup b = b.$$

In Lean~4:

\begin{minted}{lean4}
namespace OrderBook

class Lattice (α : Type) extends PartialOrder α where
  inf : α → α → α    -- meet  (⊓)
  sup : α → α → α    -- join  (⊔)
  -- meet is a lower bound
  inf_le_left  : ∀ (a b : α), inf a b ≤ a
  inf_le_right : ∀ (a b : α), inf a b ≤ b
  -- meet is the greatest lower bound
  le_inf : ∀ (a b c : α), a ≤ b → a ≤ c → a ≤ inf b c
  -- join is an upper bound
  sup_le_left  : ∀ (a b : α), a ≤ sup a b
  sup_le_right : ∀ (a b : α), b ≤ sup a b
  -- join is the least upper bound
  le_sup : ∀ (a b c : α), a ≤ c → b ≤ c → sup a b ≤ c

-- Notation
infixl:70 " ⊓ " => Lattice.inf
infixl:65 " ⊔ " => Lattice.sup

end OrderBook
\end{minted}

\subsection{Derived Properties}

We can immediately derive commutativity and idempotence from the axioms:

\begin{minted}{lean4}
namespace OrderBook
open Lattice PartialOrder

variable {α : Type} [Lattice α]

-- Commutativity of meet
theorem inf_comm (a b : α) : a ⊓ b = b ⊓ a := by
  apply le_antisymm
  · apply le_inf
    · exact inf_le_right a b
    · exact inf_le_left a b
  · apply le_inf
    · exact inf_le_right b a
    · exact inf_le_left b a

-- Idempotence of meet
theorem inf_idem (a : α) : a ⊓ a = a := by
  apply le_antisymm
  · exact inf_le_left a a
  · exact le_inf a a a (le_refl a) (le_refl a)

-- Commutativity of join
theorem sup_comm (a b : α) : a ⊔ b = b ⊔ a := by
  apply le_antisymm
  · apply le_sup
    · exact sup_le_right b a
    · exact sup_le_left b a
  · apply le_sup
    · exact sup_le_right a b
    · exact sup_le_left a b

-- Absorption: a ⊔ (a ⊓ b) = a
theorem sup_inf_self (a b : α) : a ⊔ (a ⊓ b) = a := by
  apply le_antisymm
  · apply le_sup
    · exact le_refl a
    · exact inf_le_left a b
  · exact sup_le_left a (a ⊓ b)

end OrderBook
\end{minted}

\section{The Duality Principle}

Every lattice has a \emph{dual}: swap $\sqcap \leftrightarrow \sqcup$ and
$\leq\, \leftrightarrow\, \geq$. Any theorem proved in a lattice has a dual
theorem for free.

\begin{thmbox}[Duality Principle]
  If $\varphi$ is a theorem about lattices, then so is the statement obtained
  by replacing every occurrence of $\sqcap$ with $\sqcup$, $\sqcup$ with $\sqcap$,
  $\leq$ with $\geq$, $\bot$ with $\top$, and $\top$ with $\bot$.
\end{thmbox}

\begin{minted}{lean4}
structure Dual (α : Type) where
  val : α

-- Step 1: flip ≤.  a ≤_dual b  means  b.val ≤_orig a.val.
instance {α : Type} [LE α] : LE (Dual α) where
  le a b := LE.le (α := α) b.val a.val

-- Step 2: the dual partial order. Our PartialOrder class carries its
-- own `le` field, so the flipped order is stated again here (the LE
-- instance above only supplies notation).
open OrderBook in
instance {α : Type} [PartialOrder α] : PartialOrder (Dual α) where
  le a b := PartialOrder.le (α := α) b.val a.val
  le_refl a := PartialOrder.le_refl (α := α) a.val
  le_antisymm a b h1 h2 := by
    cases a; cases b; congr 1
    exact PartialOrder.le_antisymm (α := α) _ _ h2 h1
  le_trans a b c h1 h2 := by
    cases a; cases b; cases c
    exact PartialOrder.le_trans (α := α) _ _ _ h2 h1

-- Step 3: the dual lattice swaps inf and sup.
open OrderBook in
instance {α : Type} [Lattice α] : Lattice (Dual α) where
  inf          := fun a b => ⟨Lattice.sup (α := α) a.val b.val⟩
  sup          := fun a b => ⟨Lattice.inf (α := α) a.val b.val⟩
  inf_le_left  := fun a b => Lattice.sup_le_left  (α := α) a.val b.val
  inf_le_right := fun a b => Lattice.sup_le_right (α := α) a.val b.val
  sup_le_left  := fun a b => Lattice.inf_le_left  (α := α) a.val b.val
  sup_le_right := fun a b => Lattice.inf_le_right (α := α) a.val b.val
  le_inf := fun a b c h1 h2 => by
    change Lattice.sup (α := α) b.val c.val ≤ a.val
    exact Lattice.le_sup (α := α) b.val c.val a.val h1 h2
  le_sup := fun a b c h1 h2 => by
    change LE.le (α := α) c.val (Lattice.inf (α := α) a.val b.val)
    exact Lattice.le_inf (α := α) c.val a.val b.val h1 h2
\end{minted}

\begin{notebox}
  \textbf{\lc{structure} vs \lc{def}.} Using \lc{structure Dual} rather
  than \lc{def Dual α := α} makes \lc{Dual α} and \lc{α} genuinely
  distinct types. With a transparent \lc{def}, Lean's elaborator
  sometimes picks up the wrong \lc{LE} instance for hypotheses, causing
  type mismatches even when the terms are definitionally equal. The
  \lc{structure} version forces explicit \lc{.val} unwrapping and
  \lc{⟨...⟩} rewrapping, which is more verbose but unambiguous.

  \textbf{Why \lc{cases} in \lc{le\_antisymm} and \lc{le\_trans}?}
  \lc{cases a} destructs \lc{Dual.mk v} into the underlying \lc{v : α},
  so subsequent lemmas see plain \lc{α} inequalities rather than
  \lc{Dual}-typed ones. In \lc{le\_antisymm}, \lc{congr 1} further
  reduces \lc{Dual.mk x = Dual.mk y} to \lc{x = y}.

  \textbf{Why \lc{change} in \lc{le\_inf} and \lc{le\_sup}?} After
  destructuring \lc{⟨a⟩ ⟨b⟩ ⟨c⟩}, the hypotheses \lc{h1} and \lc{h2}
  are still elaborated under the dual \lc{LE}. \lc{change} restates the
  goal entirely in terms of \lc{α}'s operations, giving Lean enough
  information to match the original lemma's type.

  \textbf{Why \lc{inf} and \lc{sup} swap} is the whole point: turning
  the Hasse diagram upside down exchanges every meet with every join.
\end{notebox}

\FloatBarrier
\section{Important Lattice Examples}

\subsection{The Two-Element Lattice $\mathbf{2}$}

\begin{minted}{lean4}
-- Bool as a lattice: false < true.
-- Core's Bool.le_antisymm / Bool.le_trans take their element
-- arguments implicitly, so we eta-expand them to match our class
-- fields, which declare the elements explicitly.
instance : OrderBook.Lattice Bool where
  le          := (· ≤ ·)   -- Bool has a built-in LE in core Lean 4
  le_refl     := Bool.le_refl
  le_antisymm := fun a b => Bool.le_antisymm
  le_trans    := fun a b c => Bool.le_trans
  inf         := (· && ·)
  sup         := (· || ·)
  inf_le_left  := by decide
  inf_le_right := by decide
  le_inf       := by decide
  sup_le_left  := by decide
  sup_le_right := by decide
  le_sup       := by decide
\end{minted}

\noindent The \lc{by decide} tactic works here because \lc{Bool} is a finite
type and all the statements are decidable.

\begin{figure}[H]
  \centering
  \begin{tikzpicture}
    \node[node] (f) at (0,0) {\texttt{false}};
    \node[node] (t) at (0,2) {\texttt{true}};
    \draw[hasseedge] (f)--(t);
    \node[right=0.3cm of t, gray, font=\scriptsize] {$\top$: \texttt{true || x = true}};
    \node[right=0.3cm of f, gray, font=\scriptsize] {$\bot$: \texttt{false \&\& x = false}};
    \node[right=0.3cm of t, gray, font=\scriptsize, yshift=-0.4cm]
      {\texttt{false} $\sqcup$ \texttt{true} $=$ \texttt{true}};
    \node[right=0.3cm of f, gray, font=\scriptsize, yshift=0.4cm]
      {\texttt{false} $\sqcap$ \texttt{true} $=$ \texttt{false}};
  \end{tikzpicture}
  \caption{The two-element lattice $\mathbf{2}$ on \texttt{Bool}.
           Meet is \texttt{\&\&}, join is \texttt{||}.}
\end{figure}

\subsection{The Powerset Lattice}

\begin{minted}{lean4}
-- Sets (α → Prop) form a lattice under subset inclusion
def Set.inter {α : Type} (s t : Set α) : Set α := fun a => s a ∧ t a
def Set.union {α : Type} (s t : Set α) : Set α := fun a => s a ∨ t a

instance {α : Type} : OrderBook.Lattice (Set α) where
  le          := Set.subset
  le_refl     := fun s a ha => ha
  le_antisymm := fun s t h1 h2 => funext fun a => propext ⟨h1 a, h2 a⟩
  le_trans    := fun s t u h1 h2 a ha => h2 a (h1 a ha)
  inf         := Set.inter
  sup         := Set.union
  inf_le_left  := fun s t a ⟨hs, _⟩ => hs
  inf_le_right := fun s t a ⟨_, ht⟩ => ht
  le_inf       := fun s t u hst hsu a ha => ⟨hst a ha, hsu a ha⟩
  sup_le_left  := fun s t a hs => Or.inl hs
  sup_le_right := fun s t a ht => Or.inr ht
  le_sup       := fun s t u hsu htu a h =>
    h.elim (hsu a) (htu a)
\end{minted}

\begin{figure}[H]
  \centering
  \begin{tikzpicture}[node distance=1.4cm]
    \node[node] (bot)  at (1.5, 0)   {$\varnothing$};
    \node[node] (x)    at (0,   1.8) {$\{x\}$};
    \node[node] (y)    at (3,   1.8) {$\{y\}$};
    \node[node] (top)  at (1.5, 3.6) {$\{x,y\}$};
    \draw[hasseedge] (bot)--(x); \draw[hasseedge] (bot)--(y);
    \draw[hasseedge] (x)--(top); \draw[hasseedge] (y)--(top);
    \node[right=0.2cm of top, gray, font=\scriptsize, text width=3.5cm]
      {$\{x\} \sqcup \{y\} = \{x,y\}$ (union)};
    \node[right=0.2cm of bot, gray, font=\scriptsize, text width=3.5cm]
      {$\{x\} \sqcap \{y\} = \varnothing$ (intersection)};
  \end{tikzpicture}
  \caption{The powerset lattice $\mathcal{P}(\{x,y\})$.
           Meet is $\cap$, join is $\cup$.
           $\{x\}$ and $\{y\}$ are incomparable.}
\end{figure}

\subsection{The Divisibility Lattice on Divisors of $n$}

For divisors of a fixed natural number $n$, divisibility is a lattice where
$\gcd$ is the meet and $\mathrm{lcm}$ is the join.

\begin{minted}{lean4}
-- On natural numbers: gcd is meet, lcm is join under divisibility order
instance : OrderBook.Lattice Nat where
  le          := (· ∣ ·)   -- divisibility
  le_refl     := Nat.dvd_refl
  le_antisymm := fun a b => Nat.dvd_antisymm
  le_trans    := fun a b c => Nat.dvd_trans
  inf         := Nat.gcd
  sup         := Nat.lcm
  inf_le_left  := Nat.gcd_dvd_left
  inf_le_right := Nat.gcd_dvd_right
  le_inf       := fun a b c => Nat.dvd_gcd
  sup_le_left  := Nat.dvd_lcm_left
  sup_le_right := Nat.dvd_lcm_right
  le_sup       := fun a b c ha hb => Nat.lcm_dvd ha hb
\end{minted}

\begin{figure}[H]
  \centering
  \begin{tikzpicture}
    \node[node] (n1)  at (3,   0) {1};
    \node[node] (n2)  at (0,   2) {2};
    \node[node] (n3)  at (6,   2) {3};
    \node[node] (n4)  at (0,   4) {4};
    \node[node] (n6)  at (4.5, 4) {6};
    \node[node] (n12) at (2.25,6) {12};
    \draw[hasseedge] (n1)--(n2);  \draw[hasseedge] (n1)--(n3);
    \draw[hasseedge] (n2)--(n4);  \draw[hasseedge] (n2)--(n6);
    \draw[hasseedge] (n3)--(n6);
    \draw[hasseedge] (n4)--(n12); \draw[hasseedge] (n6)--(n12);
    \node[gray, font=\scriptsize] at (-1.5, 3)
      {$4 \sqcap 6 = \gcd(4,6) = 2$};
    \node[gray, font=\scriptsize] at (-1.5, 3.5)
      {$4 \sqcup 6 = \mathrm{lcm}(4,6) = 12$};
    \node[gray, font=\scriptsize] at (8, 3)
      {$3 \sqcap 4 = \gcd(3,4) = 1$};
    \node[gray, font=\scriptsize] at (8, 3.5)
      {$3 \sqcup 4 = \mathrm{lcm}(3,4) = 12$};
  \end{tikzpicture}
  \caption{The divisibility lattice on divisors of $12$.
           Meet is $\gcd$, join is $\mathrm{lcm}$.
           Every incomparable pair still has a $\gcd$ and $\mathrm{lcm}$
           in the set, confirming this is a lattice.}
\end{figure}

\section{Distributive Lattices}

Not every lattice satisfies the distributive law; those that do are especially
well-behaved.

\begin{defbox}[Distributive Lattice]
  A lattice is \emph{distributive} if
  $a \sqcap (b \sqcup c) = (a \sqcap b) \sqcup (a \sqcap c)$
  (equivalently, the dual law holds too).
\end{defbox}

The powerset lattice and Boolean lattices are distributive.
Figure~\ref{fig:m3n5} shows the two smallest non-distributive lattices.

\begin{figure}[H]
  \centering
  \begin{tikzpicture}
    \begin{scope}[xshift=0cm]
      \node[font=\small\bfseries] at (1.5, 5.0) {$\mathbf{M}_3$};
      \node[node] (bot) at (1.5,0)   {$\bot$};
      \node[node] (a)   at (0,2)     {$a$};
      \node[node] (b)   at (1.5,2)   {$b$};
      \node[node] (c)   at (3,2)     {$c$};
      \node[node] (top) at (1.5,4)   {$\top$};
      \draw[hasseedge] (bot)--(a); \draw[hasseedge] (bot)--(b);
      \draw[hasseedge] (bot)--(c);
      \draw[hasseedge] (a)--(top); \draw[hasseedge] (b)--(top);
      \draw[hasseedge] (c)--(top);
      \node[gray,font=\scriptsize,text width=3cm] at (1.5,-0.8)
        {$a \sqcap (b \sqcup c) = a \sqcap \top = a$\\
         $(a \sqcap b) \sqcup (a \sqcap c) = \bot \sqcup \bot = \bot$\\
         $a \neq \bot$: \textbf{not distributive}};
    \end{scope}
    \begin{scope}[xshift=6cm]
      \node[font=\small\bfseries] at (1.5, 5.0) {$\mathbf{N}_5$};
      \node[node] (bot2) at (1.5,0)  {$\bot$};
      \node[node] (p)    at (0,1.5)  {$p$};
      \node[node] (q)    at (2.5,1.5){$q$};
      \node[node] (r)    at (2.5,3)  {$r$};
      \node[node] (top2) at (1.5,4.5){$\top$};
      \draw[hasseedge] (bot2)--(p); \draw[hasseedge] (bot2)--(q);
      \draw[hasseedge] (q)--(r);
      \draw[hasseedge] (p)--(top2); \draw[hasseedge] (r)--(top2);
      \node[gray,font=\scriptsize,text width=3cm] at (1.5,-0.8)
        {Also non-distributive.\\
         Every non-distributive\\
         lattice contains $\mathbf{M}_3$\\
         or $\mathbf{N}_5$ as a sublattice.};
    \end{scope}
  \end{tikzpicture}
  \caption{The two smallest non-distributive lattices.
           $\mathbf{M}_3$ (left) has three mutually incomparable middle
           elements; the distributive law fails because the meet
           does not distribute over the join.
           $\mathbf{N}_5$ (right) is the pentagon lattice.
           A classical theorem states that a lattice is distributive if
           and only if it contains neither as a sublattice.}
  \label{fig:m3n5}
\end{figure}

\section{Exercises}

\begin{exercisebox}[Ordering intervals by information]
  Define the type \lc{Interval} of pairs $[l, h]$ with $l, h : \mathbb{Z}$ and
  $l \leq h$. Define the \emph{information order}: $[l_1, h_1] \leq [l_2, h_2]$
  iff $l_1 \leq l_2$ and $h_2 \leq h_1$ --- that is, the second interval is
  \emph{contained} in the first, so narrower $=$ more informative $=$ higher.
  Prove that this is a partial order in Lean~4.
  Then investigate: which pairs of intervals have a meet? A join?
  (\emph{Careful}: the meet --- the smallest interval containing both --- always
  exists, but two \emph{disjoint} intervals have no join in this type, because
  there is no ``empty interval'' to represent the contradiction. Chapter~7
  completes the picture by adding one, together with unbounded ends.)
\end{exercisebox}

\begin{exercisebox}[Absorb, then join]
  Prove in Lean~4 that in any lattice:
  $a \sqcap (a \sqcup b) = a$ (absorption).
\end{exercisebox}

\begin{exercisebox}[Characterising the order]
  Prove the equivalence $a \leq b \iff a \sqcup b = b$ in Lean~4, using only
  the lattice axioms.
  \begin{minted}{lean4}
open OrderBook in
theorem le_iff_sup_eq {α : Type} [OrderBook.Lattice α] (a b : α) :
    a ≤ b ↔ a ⊔ b = b := by
  constructor
  · intro h
    apply PartialOrder.le_antisymm
    · exact Lattice.le_sup a b b h (PartialOrder.le_refl b)
    · exact Lattice.sup_le_right a b
  · intro h
    have : b = a ⊔ b := h.symm
    rw [this]
    exact Lattice.sup_le_left a b
  \end{minted}
\end{exercisebox}

\begin{exercisebox}[The product lattice]
  Given two lattices \lc{[Lattice α]} and \lc{[Lattice β]}, define a
  \lc{Lattice (α × β)} instance with the component-wise order:
  $(a_1, b_1) \leq (a_2, b_2) \iff a_1 \leq a_2 \wedge b_1 \leq b_2$,
  and component-wise meet and join:
  $(a_1, b_1) \sqcap (a_2, b_2) = (a_1 \sqcap a_2,\, b_1 \sqcap b_2)$,
  $(a_1, b_1) \sqcup (a_2, b_2) = (a_1 \sqcup a_2,\, b_1 \sqcup b_2)$.
  Prove all eight lattice axioms.
  (\emph{Hint}: each axiom reduces to the corresponding axiom applied component-wise.)
\end{exercisebox}

% ============================================================
\chapter{Complete Lattices and the Fixed-Point Theorem}

Lattices let us take meets and joins of \emph{pairs} of elements. A
\emph{complete} lattice extends this to arbitrary (possibly infinite) subsets.
This extra power buys us one of the most useful theorems in theoretical
computer science: the Knaster--Tarski fixed-point theorem.

\section{Complete Lattices}

\begin{defbox}[Complete Lattice]
  A lattice $(L, \leq)$ is \emph{complete} if every subset $S \subseteq L$
  has a supremum $\bigvee S$ and an infimum $\bigwedge S$ in $L$.
\end{defbox}

\begin{notebox}
  In particular, the empty set has a supremum (which becomes $\bot$) and
  infimum (which becomes $\top$), so every complete lattice is bounded.
  Taking $S = L$ gives $\bigvee L = \top$ and $\bigwedge L = \bot$.
\end{notebox}

We encode complete lattices using a predicate representing a subset:

\begin{minted}{lean4}
namespace OrderBook

-- We represent a "subset" as a predicate α → Prop.
-- This is the same definition as `Set α` from Chapter 1
-- (`def Set (α : Type) : Type := α → Prop`).
-- We use `abbrev` rather than `def` so that Lean treats `Pred α`
-- as completely transparent — unfolding it automatically during
-- elaboration without needing `simp` or `unfold`.  This makes the
-- `sSup`/`sInf` signatures below easier to work with.
-- `def Set` was opaque because we wanted `Set α` to have its own
-- identity for attaching instances; here we just want a readable alias.
abbrev Pred (α : Type) := α → Prop

class CompleteLattice (α : Type) extends Lattice α where
  sSup : Pred α → α    -- supremum of a set  (a "big ⊔")
  sInf : Pred α → α    -- infimum of a set   (a "big ⊓")
  le_sSup : ∀ (s : Pred α) (a : α), s a → a ≤ sSup s
  sSup_le : ∀ (s : Pred α) (u : α), (∀ a, s a → a ≤ u) → sSup s ≤ u
  sInf_le : ∀ (s : Pred α) (a : α), s a → sInf s ≤ a
  le_sInf : ∀ (s : Pred α) (l : α), (∀ a, s a → l ≤ a) → l ≤ sInf s

end OrderBook
\end{minted}

The powerset lattice is the canonical example of a complete lattice:
$\bigvee \mathcal{F} = \bigcup \mathcal{F}$ and
$\bigwedge \mathcal{F} = \bigcap \mathcal{F}$.

\begin{minted}{lean4}
-- Set α is a complete lattice
instance {α : Type} : OrderBook.CompleteLattice (Set α) where
  -- (inheriting the Lattice instance from Chapter 3)
  sSup := fun F a => ∃ s : Set α, F s ∧ s a    -- union of a family
  sInf := fun F a => ∀ s : Set α, F s → s a    -- intersection of a family
  le_sSup := fun F s hs a ha => ⟨s, hs, ha⟩
  sSup_le := fun F u hu a ⟨s, hs, ha⟩ => hu s hs a ha
  sInf_le := fun F s hs a ha => ha s hs
  le_sInf := fun F l hl a ha s hs => hl s hs a ha
  -- (le, inf, sup fields inherited)
  le          := Set.subset
  le_refl     := fun s a h => h
  le_antisymm := fun s t h1 h2 => funext fun a => propext ⟨h1 a, h2 a⟩
  le_trans    := fun s t u h1 h2 a ha => h2 a (h1 a ha)
  inf         := Set.inter; sup := Set.union
  inf_le_left  := fun s t a ⟨hs,_⟩ => hs
  inf_le_right := fun s t a ⟨_,ht⟩ => ht
  le_inf       := fun s t u hst hsu a ha => ⟨hst a ha, hsu a ha⟩
  sup_le_left  := fun s t a hs => Or.inl hs
  sup_le_right := fun s t a ht => Or.inr ht
  le_sup       := fun s t u hsu htu a h => h.elim (hsu a) (htu a)
\end{minted}

\section{Why Fixed Points Matter}

A \emph{fixed point} of $f : L \to L$ is an element $x$ where $f(x) = x$.
This is where an iterative process \emph{stabilises}: no further application
of $f$ changes anything.

\begin{itemize}
  \item \textbf{Propagator networks.} Cell values stop changing when the
    propagators reach a fixed point of ``apply all propagators once''. That
    stable state is the least fixed point --- the most conservative set of
    conclusions forced by the constraints, nothing more.

  \item \textbf{Dataflow analysis.} A compiler iterates a flow function until
    it stabilises. The result is the least fixed point of that function.

  \item \textbf{Denotational semantics.} The meaning of a recursive program
    is the least fixed point of a semantic functional.

  \item \textbf{Regular expressions.} The set matched by $a^*$ is the least
    fixed point of $S \mapsto \{\varepsilon\} \cup a \cdot S$.
\end{itemize}

\begin{notebox}
  We will return to the propagator connection in Chapter~6, where network
  execution terminates at the least fixed point of the combined propagation
  function.
\end{notebox}

\section{The Knaster--Tarski Fixed-Point Theorem}

\begin{thmbox}[Knaster--Tarski (1955)]
  Let $(L, \leq)$ be a complete lattice and $f : L \to L$ a monotone function.
  Then $f$ has a \emph{least fixed point}
  $\mu f = \bigwedge \{ x \mid f(x) \leq x \}$
  and a \emph{greatest fixed point}
  $\nu f = \bigvee \{ x \mid x \leq f(x) \}$.
\end{thmbox}

\begin{proof}
  Let $P = \{ x \in L \mid f(x) \leq x \}$ (the set of \emph{pre-fixed points})
  and let $m = \bigwedge P$.

  \textit{Step 1: $f(m) \leq m$.}
  For every $x \in P$ we have $m \leq x$, so $f(m) \leq f(x) \leq x$ by
  monotonicity. Hence $f(m)$ is a lower bound of $P$, so $f(m) \leq m$.

  \textit{Step 2: $m \leq f(m)$.}
  From Step~1, $f(f(m)) \leq f(m)$ (applying $f$ to $f(m) \leq m$ and using
  monotonicity gives $f(f(m)) \leq f(m)$, so $f(m) \in P$). Since $m$ is the
  infimum of $P$, we have $m \leq f(m)$.

  \textit{Conclusion:} $f(m) = m$, so $m$ is a fixed point. Any other fixed
  point $x^*$ satisfies $f(x^*) = x^* \leq x^*$, so $x^* \in P$, hence
  $m \leq x^*$. So $m$ is the \emph{least} fixed point.
\end{proof}

\begin{minted}{lean4}
namespace OrderBook
open CompleteLattice PartialOrder

-- Knaster-Tarski: least fixed point of a monotone function
-- on a complete lattice
theorem knaster_tarski {α : Type} [CompleteLattice α]
    (f : α → α) (hf : Monotone f) :
    ∃ (lfp : α), f lfp = lfp ∧
      ∀ (x : α), f x = x → lfp ≤ x := by
  -- Define P = {x | f(x) ≤ x} and m = ⊓ P
  let P : Pred α := fun x => f x ≤ x
  let m := sInf P
  -- Step 1: f(m) ≤ m
  have step1 : f m ≤ m := by
    apply le_sInf
    intro x hx
    -- m ≤ x, so f(m) ≤ f(x) ≤ x.  (The class states le_trans with
    -- explicit element arguments, hence the placeholders `_ _ _`.)
    exact le_trans _ _ _ (hf.map_le m x (sInf_le P x hx)) hx
  -- Step 2: m ≤ f(m)  (f(m) ∈ P since f(f(m)) ≤ f(m))
  have step2 : m ≤ f m := by
    apply sInf_le
    exact hf.map_le _ _ step1
  -- Conclude m is a fixed point
  have fixed : f m = m := le_antisymm _ _ step1 step2
  refine ⟨m, fixed, ?_⟩
  -- Minimality: every fixed point is a pre-fixed point, so m ≤ x.
  intro x hx
  apply sInf_le
  show f x ≤ x
  rw [hx]
  exact le_refl x
end OrderBook
\end{minted}

\begin{figure}[H]
  \centering
  \begin{tikzpicture}
    \node[node,fill=green!10,draw=green!50!black] (b)  at (0,0)   {$\bot$};
    \node[node,fill=green!10,draw=green!50!black] (f1) at (0,1.5) {$f(\bot)$};
    \node[node,fill=green!10,draw=green!50!black] (f2) at (0,3)   {$f^2(\bot)$};
    \node[gray,font=\small]                       (dots) at (0,4.2){$\vdots$};
    \node[node,fill=orange!15,draw=orange!70!black,
          minimum size=9mm]                       (lfp) at (0,5.4) {$\mu f$};
    \node[gray,font=\small]                       (d2)  at (0,6.5) {$\vdots$};
    \node[node]                                   (top) at (0,7.5) {$\top$};
    \draw[hasseedge] (b)--(f1);  \draw[hasseedge] (f1)--(f2);
    \draw[hasseedge] (f2)--(dots);\draw[hasseedge] (dots)--(lfp);
    \draw[hasseedge] (lfp)--(d2); \draw[hasseedge] (d2)--(top);
    \draw[-{Stealth},blue!60,thick,bend right=50]
      (b.east)  to node[right,font=\scriptsize,blue!60]{$f$} (f1.east);
    \draw[-{Stealth},blue!60,thick,bend right=50]
      (f1.east) to node[right,font=\scriptsize,blue!60]{$f$} (f2.east);
    \node[right=1.2cm of b,  gray,font=\scriptsize]{no information yet};
    \node[right=1.2cm of lfp,font=\scriptsize,
          orange!70!black,text width=3.5cm]
      {least fixed point:\\$f(\mu f) = \mu f$\\reached when chain stabilises};
    \draw[decorate,decoration={brace,amplitude=6pt},gray!60]
      (-0.7,0) -- (-0.7,5.4)
      node[midway,left=8pt,gray,font=\scriptsize,text width=2cm,align=right]
        {ascending\\Kleene chain};
  \end{tikzpicture}
  \caption{The ascending Kleene chain starting from $\bot$ under a
           monotone function $f$. Each application of $f$ moves upward.
           If the lattice satisfies the ascending chain condition the
           sequence stabilises at the least fixed point $\mu f$.}
\end{figure}

\section{Chains and Convergence}

Chains are the lattice-theoretic version of sequences in calculus, and understanding
them is essential for understanding why iterative computations terminate.

\begin{defbox}[Chain]
  A \emph{chain} in a poset $(P, \leq)$ is a subset $C \subseteq P$ where every
  pair of elements is comparable: for all $a, b \in C$, either $a \leq b$ or
  $b \leq a$. An \emph{ascending chain} is a chain written as a sequence
  $a_0 \leq a_1 \leq a_2 \leq \cdots$.
\end{defbox}

The parallel with calculus sequences is direct:

\begin{center}
\begin{tabular}{ll}
  \textbf{Calculus} & \textbf{Lattice theory} \\
  \hline
  Sequence $a_0 \leq a_1 \leq a_2 \leq \cdots$ & Ascending chain \\
  Limit $\lim_{n\to\infty} a_n$ & Supremum $\bigsqcup_n a_n$ \\
  Convergence & Chain stabilising \\
  Completeness of $\mathbb{R}$ & Completeness of the lattice \\
  Cauchy sequence & Ascending chain condition \\
\end{tabular}
\end{center}

Just as the completeness of $\mathbb{R}$ guarantees that every bounded ascending
sequence has a limit, the completeness of a lattice guarantees that every chain
has a supremum. Without completeness, the supremum might not exist in your
structure --- exactly as $\sqrt{2}$ does not exist in $\mathbb{Q}$ even though
there are rational sequences converging to it.

\subsection*{The Ascending Chain Condition}

A poset satisfies the \emph{ascending chain condition} (ACC) if every strictly
ascending chain $a_0 < a_1 < a_2 < \cdots$ is finite --- there are no infinite
strictly ascending sequences.

This is the lattice analogue of a sequence that reaches its limit in finite steps
rather than merely approaching it. When the ACC holds, the Kleene chain
$$\bot \leq f(\bot) \leq f^2(\bot) \leq \cdots$$
must eventually stabilise at some $f^n(\bot) = f^{n+1}(\bot)$, giving a fixed
point in finite iterations.

\begin{exbox}[ACC in practice]
  The flat lattice \lc{FlatNat} that we will build in Chapter~6
  ($\bot = $ unknown, exact values in the middle, $\top = $ contradiction)
  satisfies the ACC: the longest strictly ascending chains are
  $\bot < \mathrm{known}(n) < \mathrm{conflict}$, of length 3. So any
  propagator network using \lc{FlatNat} cells terminates after at most two
  updates per cell.

  The interval lattice of Chapter~7 (where narrower intervals sit
  \emph{higher}, and the ends may be unbounded) does \emph{not} satisfy the
  ACC: $(-\infty,+\infty) < [0,+\infty) < [1,+\infty) < [2,+\infty) < \cdots$
  is an infinite strictly ascending chain. This is why interval propagators
  in general need a \emph{widening} operator to force termination ---
  although, as we will see, well-behaved constraint systems often reach a
  fixed point on their own.
\end{exbox}

\subsection*{$\bot$ is the starting point, not the limit}

A common point of confusion: the supremum of the Kleene chain is \emph{not}
$\bot$. $\bot$ is the \emph{minimum} of the chain --- the smallest element in
it. The supremum is the smallest element above the whole chain, which is
wherever the iteration converges.

In the calculus analogy: for the sequence $0, \tfrac{1}{2}, \tfrac{3}{4},
\tfrac{7}{8}, \ldots$ converging to 1, the minimum of the sequence is 0 but
the supremum (limit) is 1. The starting point and the limit are different
objects. The Kleene chain behaves the same way: $\bot$ is where you start,
$\mu f$ is where you end up.

\section{Exercises}

\begin{exercisebox}[Greatest fixed point]
  Using the Knaster--Tarski proof as a template, prove the existence of the
  \emph{greatest} fixed point $\nu f = \bigvee \{ x \mid x \leq f(x) \}$.
\end{exercisebox}

\begin{exercisebox}[Ascending Kleene chain]
  For a complete lattice with a bottom element $\bot$ and a monotone $f$,
  the sequence $\bot \leq f(\bot) \leq f^2(\bot) \leq \cdots$ is called the
  \emph{ascending Kleene chain}. If the lattice satisfies the
  \emph{ascending chain condition} (every ascending chain eventually stabilises),
  this chain reaches $\mu f$ in finitely many steps.
  Prove in Lean~4 that each element of this sequence is $\leq$ the next.
\end{exercisebox}

% ============================================================
\part{Propagator Networks}

% ============================================================
\chapter{The Propagator Model}

\begin{quote}
  \itshape
  ``A propagator network is a collection of machines that compute
  values for cells by combining information from other cells.''\\
  --- G.J.\ Sussman \& A.\ Radul, \emph{The Art of the Propagator}, 2009
\end{quote}

We now have all the mathematical machinery we need. This chapter explains the
propagator model from first principles and makes precise why lattices are not
merely convenient but \emph{necessary}.

\section{Cells: Accumulating Partial Information}

The central insight of propagator networks is a shift in perspective:

\begin{center}
  \begin{tikzpicture}
    \node[draw=gray,fill=gray!10,rounded corners,minimum width=5cm,
          minimum height=0.8cm, font=\small]
      (old) {Traditional cell: stores a \textbf{single value}};
    \node[draw=blue!60!black,fill=blue!6,rounded corners,minimum width=5cm,
          minimum height=0.8cm, font=\small, below=0.6cm of old]
      (new) {Propagator cell: accumulates \textbf{partial information}};
    \draw[-{Stealth},thick,blue!60!black] (old) -- (new)
      node[midway,right=4pt,font=\small\itshape]{evolution};
  \end{tikzpicture}
\end{center}

A cell does not \emph{store a value}; it stores \emph{everything we know so
far} about a value. New information can arrive from multiple sources and be
merged in. The key requirement is that information can only \emph{increase}:
once we know something, we cannot un-know it.

This is precisely the structure of a \emph{join-semilattice with a bottom
element}: the bottom represents ``we know nothing'', and merging two pieces of
information is the join $(\sqcup)$.

\begin{defbox}[Information Order]
  A \emph{bounded join-semilattice} is a partial order $(J, \leq)$ equipped with
  a join $\sqcup$ and a bottom $\bot$ such that:
  \begin{itemize}[nosep]
    \item $a \sqcup b$ is the least upper bound of $a$ and $b$,
    \item $\bot \leq a$ for all $a$,
    \item $a \leq b$ is read as ``$b$ contains at least as much information as $a$''.
  \end{itemize}
  We call this the \emph{information order}.
\end{defbox}

\begin{minted}{lean4}
namespace Propagators

class BoundedJoinSemilattice (α : Type) extends OrderBook.PartialOrder α where
  sup : α → α → α
  bot : α
  le_sup_left  : ∀ (a b : α), a ≤ sup a b
  le_sup_right : ∀ (a b : α), b ≤ sup a b
  sup_le       : ∀ (a b c : α), a ≤ c → b ≤ c → sup a b ≤ c
  bot_le       : ∀ (a : α), bot ≤ a

-- Convenient notation
infixl:65 " ⊔ " => BoundedJoinSemilattice.sup
notation    "⊥"  => BoundedJoinSemilattice.bot

end Propagators
\end{minted}

\section{Why Information Can Only Increase}

Suppose a propagator computes that ``$x$ is between 3 and 7''. Later another
propagator deduces that ``$x$ is between 5 and 12''. Combining these we know
``$x$ is between 5 and 7''. Notice: the combined knowledge is \emph{strictly
more specific} than either input.

If we allowed information to \emph{decrease} (e.g., forgetting the $[3,7]$
constraint and keeping only $[5,12]$), we would lose soundness: our cell might
claim values that the original propagators had ruled out.

\begin{thmbox}[Monotonicity Ensures Soundness]
  In a propagator network over a bounded join-semilattice, if every propagator
  is monotone, then the network is \emph{sound}: the final cell values are
  consistent with all constraints.
\end{thmbox}

Intuitively: since each step can only add information (move up in the lattice),
the final state contains at least as much information as any intermediate state.
Nothing is lost; nothing false is added (assuming correct propagators).

\section{Propagators as Monotone Functions}

\begin{defbox}[Propagator]
  Given cells with values in $(J, \leq)$, a \emph{propagator} is a monotone
  function $f : J^n \to J$ (reading from $n$ input cells and writing to one
  output cell by merging the result).
\end{defbox}

Writing is done by merging: the output cell $c$ receives $c \leftarrow c \sqcup f(\text{inputs})$.
Since $\sqcup$ is the least upper bound, this can only increase $c$.

\begin{minted}{lean4}
namespace Propagators

-- A unary propagator: a pure function with a bundled monotonicity proof.
-- The condition fn a ≤ fn b is the standard order-theoretic definition.
-- (Equivalent to BoundedJoinSemilattice.sup (fn a) (fn b) = fn b,
--  but the ≤ form is cleaner to state and prove.)
structure Propagator1 (α β : Type)
    [BoundedJoinSemilattice α]
    [BoundedJoinSemilattice β] where
  fn       : α → β
  monotone : ∀ (a b : α), a ≤ b → fn a ≤ fn b

-- A binary propagator: monotone in both arguments.
structure Propagator2 (α β γ : Type)
    [BoundedJoinSemilattice α]
    [BoundedJoinSemilattice β]
    [BoundedJoinSemilattice γ] where
  fn       : α → β → γ
  monotone : ∀ (a1 a2 : α) (b1 b2 : β),
    a1 ≤ a2 → b1 ≤ b2 → fn a1 b1 ≤ fn a2 b2

end Propagators
\end{minted}

\begin{notebox}
  \textbf{Specification vs.\ implementation.}
  \lc{Propagator1} and \lc{Propagator2} capture what a propagator
  \emph{is} mathematically: a function with a bundled monotonicity proof.
  They are the \emph{specification}.

  In practice we use \lc{PropStep = IO Bool}: a plain \lc{IO} action that
  reads cells, computes, writes, and returns whether anything changed.
  There is no bundled proof --- the user is trusted to write monotone
  propagators. Section~\ref{sec:flatnat-monotone} below proves that
  \lc{addProp} and \lc{subProp} do in fact satisfy the \lc{Propagator2}
  specification. For the capstone projects the proofs would be analogous
  but verbose, so we omit them.
\end{notebox}

\section{Networks, Scheduling, and Quiescence}

A \emph{propagator network} is a directed graph where nodes are cells and edges
are propagators. Execution proceeds as follows:

\begin{enumerate}
  \item Initialise all cells to $\bot$.
  \item Place all propagators on a \emph{work queue}.
  \item While the queue is non-empty:
    \begin{enumerate}
      \item Take a propagator $p$ from the queue.
      \item Compute the new value $v = p(\text{input cells})$.
      \item Merge $v$ into the output cell: $c \leftarrow c \sqcup v$.
      \item If $c$ changed, add all propagators that read from $c$ back to the queue.
    \end{enumerate}
  \item The network has reached \emph{quiescence} when the queue is empty.
\end{enumerate}

\begin{figure}[H]
  \centering
  \begin{tikzpicture}[
    cell/.style={draw=blue!60!black, fill=blue!6, rounded corners=4pt,
                 minimum width=1.1cm, minimum height=0.7cm, font=\small},
    prop/.style={draw=orange!70!black, fill=orange!8, diamond,
                 minimum width=1.3cm, minimum height=0.7cm, font=\scriptsize,
                 inner sep=1pt},
    arr/.style={-{Stealth[scale=0.8]}, thick, gray!70},
  ]
    \node[cell] (cx)   at (0, 0)    {$c_x$};
    \node[cell] (cy)   at (0, -2.5) {$c_y$};
    \node[cell] (csum) at (4, -1.2) {$c_\text{sum}$};
    \node[cell] (ck)   at (4, 1.2)  {$c_k$};
    \node[cell] (cz)   at (8, -1.2) {$c_z$};
    \node[prop] (add)  at (2, -1.2) {\texttt{add}};
    \node[prop] (mul)  at (6, -1.2) {\texttt{mul}};
    \draw[arr] (cx) -- (add); \draw[arr] (cy) -- (add);
    \draw[arr] (csum) -- (mul); \draw[arr] (ck) -- (mul);
    \draw[arr] (add) -- (csum) node[midway,above,gray,font=\scriptsize]{writes};
    \draw[arr] (mul) -- (cz)   node[midway,above,gray,font=\scriptsize]{writes};
    \draw[arr,bend left=25,dashed,blue!50] (csum) to
      node[above,font=\scriptsize,blue!60]{back-prop} (add);
    \node[below=0.05cm of cx,   gray,font=\scriptsize]{val: $\bot \to 3$};
    \node[below=0.05cm of cy,   gray,font=\scriptsize]{val: $\bot \to 7$};
    \node[below=0.05cm of csum, gray,font=\scriptsize]{val: $\bot \to 10$};
    \node[above=0.05cm of ck,   gray,font=\scriptsize]{val: $\bot \to 7$};
    \node[below=0.05cm of cz,   gray,font=\scriptsize]{val: $\bot \to 70$};
    \node[cell, right=0.3cm] at (0,-4.2) {};
    \node[font=\scriptsize,right=1.3cm] at (0,-4.2) {Cell (stores a lattice value)};
    \node[prop] at (4.5,-4.2) {};
    \node[font=\scriptsize,right=5.5cm] at (0,-4.2) {Propagator (monotone function)};
  \end{tikzpicture}
  \caption{A small propagator network computing $x + y = \mathrm{sum}$ and
           $\mathrm{sum} \times k = z$ (with $x=3$, $y=7$, $k=7$).
           Cells hold lattice values initialised to $\bot$. Propagators
           read from inputs and merge results into output cells via $\sqcup$.
           The dashed arrow shows back-propagation: when $c_\text{sum}$ is
           updated the \texttt{add} propagator is re-queued.}
\end{figure}

\begin{thmbox}[Convergence]
  If the lattice satisfies the \emph{ascending chain condition}
  (every strictly ascending sequence is finite), then the network reaches
  quiescence in finite time.
\end{thmbox}

This is immediate: each cell value is non-decreasing and the chain condition
guarantees it cannot increase indefinitely.

For the Sudoku capstone the domain lattice is finite, so convergence is
guaranteed outright; for the type-inference capstone each program generates
finitely many constraints whose refinements stabilise. The interval lattice
of Chapter~7 is the cautionary tale: it fails the ACC, and we will be
careful to say why its example nevertheless terminates.

\section{The Fundamental Role of the Lattice}

Let us be precise about why we need a lattice rather than merely some arbitrary
type.

\begin{itemize}
  \item \textbf{Join ($\sqcup$) for merging.}
    When two propagators independently derive information about the same cell,
    we need to combine their results. The join gives the \emph{least informative}
    value that is \emph{at least as informative} as both. No information is lost;
    no information is invented.

  \item \textbf{Bottom ($\bot$) for initialisation.}
    Before any propagator has run, a cell contains \emph{no information}.
    The bottom element represents this ``empty'' state.

  \item \textbf{Partial order ($\leq$) for monitoring change.}
    We must know whether a cell's value actually changed after merging, to
    decide whether to re-schedule dependent propagators. The partial order
    provides this: if $c \sqcup v = c$ (the new value adds no information),
    nothing changed.

  \item \textbf{Antisymmetry for termination.}
    If a cell's value can only go up and the lattice is finite, the antisymmetry
    of $\leq$ ensures each value is visited at most once.
\end{itemize}

Any type that fails to be a join-semilattice with $\bot$ is simply not a valid
cell domain for a propagator network. The lattice structure is not a convenience
--- it is a prerequisite.

% ============================================================
\chapter{Building a Propagator Network in Lean~4}

Theory in hand, we now build a complete, runnable propagator network in Lean~4,
starting from our typeclass definitions and ending with a working scheduler.

\section{Interlude: Imperative Lean in Half a Page}

Everything so far has been pure functions and proofs. Propagator networks,
however, are stateful machines --- cells get updated in place --- so from here
on the code looks imperative. Lean makes this look like any mainstream
language via \lc{do}-notation:

\begin{minted}{lean4}
def imperativeTour : IO Unit := do
  let r ← IO.mkRef 0            -- a mutable reference holding a Nat
  let mut best := 0             -- a local mutable variable
  for i in [0:5] do             -- i runs over 0, 1, 2, 3, 4
    r.modify (· + i)            -- add i to the reference's contents
    if i > best then best := i
  let total ← r.get             -- read the reference back
  IO.println s!"total = {total}, best = {best}"

#eval imperativeTour   -- total = 10, best = 4
\end{minted}

Four things to know, in decreasing order of importance:
\begin{itemize}
  \item \textbf{\lc{do} sequences effects.} A \lc{do} block runs in a
    \emph{monad} --- here \lc{IO}, the type of actions that may touch the
    outside world. \lc{let x ← action} \emph{runs} the action and binds its
    result; plain \lc{let x := e} is an ordinary pure binding.
  \item \textbf{\lc{IO.Ref α} is a mutable memory cell} holding an \lc{α},
    with \lc{IO.mkRef}, \lc{.get}, \lc{.set}, and \lc{.modify}. Our
    propagator cells wrap exactly one of these.
  \item \textbf{\lc{mut}, \lc{for}, \lc{while}, \lc{if} inside \lc{do} are
    sugar.} The compiler rewrites them into folds and recursive calls with
    the mutable variables threaded through as accumulators --- no actual
    mutation of \lc{best} ever happens. The full desugaring story is in the
    primer of Appendix~\ref{app:quickref}.
  \item \textbf{\lc{Id.run} lets pure code use the same sugar.} Running a
    \lc{do} block in the do-nothing monad \lc{Id} and unwrapping it with
    \lc{Id.run} gives imperative-looking code with a pure result. The Sudoku
    solver in Chapter~8 uses this heavily.
\end{itemize}

That is all the imperative Lean this book needs.

\section{Mutable Cells}

Lean~4 has an \lc{IO} monad with mutable references (\lc{IO.Ref}). We use
these to implement cells.

\begin{minted}{lean4}
namespace Propagators

-- A Cell wraps an IO.Ref and a type with a BoundedJoinSemilattice
structure Cell (α : Type) where
  ref      : IO.Ref α

variable {α : Type} [BoundedJoinSemilattice α] [DecidableEq α]

-- Create a new cell initialised to ⊥
def Cell.new : IO (Cell α) := do
  let r ← IO.mkRef (BoundedJoinSemilattice.bot (α := α))
  return { ref := r }

-- Read the current value
def Cell.read (c : Cell α) : IO α := c.ref.get

-- Attempt to write new information.
-- Returns true if the cell's value increased.
def Cell.write (c : Cell α) (v : α) : IO Bool := do
  let old ← c.ref.get
  let merged := old ⊔ v
  if merged == old then
    return false          -- no new information
  else
    c.ref.set merged
    return true           -- cell updated, re-schedule dependents

end Propagators
\end{minted}

\begin{notebox}
  We require \lc{DecidableEq α} to compare \lc{merged == old}. For the
  finite lattices we use in our capstone projects this is always available.
\end{notebox}

\section{The Network and Scheduler}

\begin{minted}{lean4}
namespace Propagators

-- A propagator step is an IO action that returns
-- true if it caused any cell to change.
abbrev PropStep := IO Bool

-- Run a list of propagators to fixpoint.
-- Each pass re-runs all propagators; we stop when
-- a full pass causes no changes.
def runToFixpoint (steps : Array PropStep) : IO Unit := do
  let mut changed := true
  while changed do
    changed := false
    for step in steps do
      let c ← step
      if c then changed := true

-- A smarter scheduler that only re-runs propagators
-- whose inputs changed (using a dependency queue).
structure Network where
  steps    : Array PropStep
  -- In a production implementation, steps would be indexed
  -- and cells would hold lists of dependent propagator IDs.
  -- For clarity we use the simple "run all" scheduler here.

def Network.run (n : Network) : IO Unit :=
  runToFixpoint n.steps

end Propagators
\end{minted}

\section{Building Propagators: Numeric Example}

Let us define the \emph{flat natural number lattice}: information is either
``unknown'' ($\bot$), a specific number, or ``contradiction'' ($\top$).

\begin{minted}{lean4}
namespace Propagators

inductive FlatNat where
  | unknown      : FlatNat         -- ⊥: no information
  | known (n : Nat) : FlatNat      -- exactly n
  | conflict     : FlatNat         -- ⊤: contradiction

-- We give the order a NAME rather than an anonymous match inside the
-- instance: later proofs can then unfold it with `simp [FlatNat.le]`,
-- and the rfl-lemma `le_def` below lets simp see through `≤`.
def FlatNat.le (a b : FlatNat) : Prop :=
  match a, b with
    | .unknown, _      => True     -- unknown ≤ everything
    | _, .conflict     => True     -- everything ≤ conflict
    | .known m, .known n => m = n  -- known values only ≤ themselves
    | .conflict, .unknown => False
    | .conflict, .known _ => False
    | .known _, .unknown  => False

instance : OrderBook.PartialOrder FlatNat where
  le := FlatNat.le
  le_refl  := by intro a; cases a <;> simp [FlatNat.le]
  le_antisymm := by
    intro a b hab hba
    cases a <;> cases b <;> simp_all [FlatNat.le]
  le_trans := by
    intro a b c hab hbc
    cases a <;> cases b <;> cases c <;> simp_all [FlatNat.le]

theorem FlatNat.le_def (a b : FlatNat) : (a ≤ b) = FlatNat.le a b := rfl

instance : BoundedJoinSemilattice FlatNat where
  le          := FlatNat.le
  le_refl     := OrderBook.PartialOrder.le_refl
  le_antisymm := OrderBook.PartialOrder.le_antisymm
  le_trans    := OrderBook.PartialOrder.le_trans
  bot         := .unknown
  sup a b := match a, b with
    | .unknown, x      => x
    | x, .unknown      => x
    | .known m, .known n => if m = n then .known m else .conflict
    | _, _             => .conflict
  bot_le      := by intro a; cases a <;> simp [FlatNat.le_def, FlatNat.le]
  le_sup_left  := by
    intro a b
    cases a <;> cases b <;> simp_all [FlatNat.le_def, FlatNat.le]
    rename_i m n
    by_cases h : m = n <;> simp [FlatNat.le_def, FlatNat.le, h]
  le_sup_right := by
    intro a b
    cases a <;> cases b <;> simp_all [FlatNat.le_def, FlatNat.le]
    rename_i m n
    by_cases h : m = n <;> simp [FlatNat.le_def, FlatNat.le, h]
  -- The known/known case needs its own `if m = n` split at the end.
  sup_le       := by
    intro a b c hac hbc
    cases a <;> cases b <;> cases c <;> simp_all [FlatNat.le_def, FlatNat.le]
    rename_i m n
    by_cases h : m = n <;> simp [FlatNat.le, h]

deriving instance DecidableEq for FlatNat

end Propagators
\end{minted}

\section{Monotonicity Proofs for \texttt{addProp} and \texttt{subProp}}
\label{sec:flatnat-monotone}

The \lc{PropStep} actions \lc{addProp} and \lc{subProp} implement pure functions
on \lc{FlatNat}. We can extract those functions and prove they satisfy the
\lc{Propagator2} specification.

\textbf{Why conflict propagates.}
There is a subtlety: the pure function must return \lc{conflict} (not \lc{unknown})
when an input is \lc{conflict}, otherwise monotonicity fails.
Suppose we returned \lc{unknown} for conflicted inputs. We have
\lc{known 3 ≤ conflict} in \lc{FlatNat}, so monotonicity would demand
\lc{addFn (known 3) (known 2) ≤ addFn conflict (known 2)} --- but the left
side is \lc{known 5}, the right side would be \lc{unknown}, and
\lc{known 5} $\not\leq$ \lc{unknown}.
Returning \lc{conflict} fixes this: \lc{known 5 ≤ conflict} always holds.

\begin{minted}{lean4}
namespace Propagators

-- Pure addition: both inputs known → known sum.
-- Either input conflicted → conflict propagates.
-- Otherwise → unknown (no information yet).
def FlatNat.addFn : FlatNat → FlatNat → FlatNat
  | .known a,  .known b  => .known (a + b)
  | .conflict, _         => .conflict
  | _,         .conflict => .conflict
  | _,         _         => .unknown

-- Pure subtraction: known total - known a, with conflict on underflow.
def FlatNat.subFn (s x : FlatNat) : FlatNat :=
  match s, x with
  | .known total, .known a =>
      if total ≥ a then .known (total - a) else .conflict
  | .conflict, _ | _, .conflict => .conflict
  | _,         _               => .unknown

-- Proof: addFn is monotone.
-- Strategy: case-split all 81 constructor combinations (3 × 3 × 3 × 3).
-- Most cases are trivial:
--   • unknown ≤ anything  →  addFn unknown ... = unknown ≤ anything.
--   • anything ≤ conflict  →  addFn ... conflict = conflict, and anything ≤ conflict.
--   • known m ≤ unknown is False  →  hypothesis contradiction closes the goal.
-- The only non-trivial case is known/known on both sides, where the
-- hypotheses force the arguments (and hence both sides) to be equal.
-- simp_all with the unfolded order closes every case.
def addPropagator : Propagator2 FlatNat FlatNat FlatNat where
  fn       := FlatNat.addFn
  monotone := by
    intro a1 a2 b1 b2 ha hb
    cases a1 <;> cases a2 <;> cases b1 <;> cases b2
      <;> simp_all [FlatNat.addFn, FlatNat.le_def, FlatNat.le]

-- Proof: subFn is monotone.
-- Same strategy. After the case split, the remaining goals are stuck
-- on the `if total ≥ a` scrutinees; `split <;> simp_all` finishes them.
def subPropagator : Propagator2 FlatNat FlatNat FlatNat where
  fn       := FlatNat.subFn
  monotone := by
    intro a1 a2 b1 b2 ha hb
    cases a1 <;> cases a2 <;> cases b1 <;> cases b2
      <;> simp_all [FlatNat.subFn, FlatNat.le_def, FlatNat.le]
      <;> split <;> simp_all

-- Connection to PropStep actions:
-- addProp cx cy csum (§6.6) is addPropagator.fn wrapped in IO:
--   it reads cx and cy, computes addFn, and merges into csum.
-- subProp csum cx cy wraps subPropagator.fn similarly.
-- The Propagator2 instances certify that these IO actions are sound.

end Propagators
\end{minted}

\section{A Complete Running Example}

\begin{exbox}[Arithmetic propagation]
  We know: $x + y = 10$ and $x = 3$. A propagator should deduce $y = 7$.
\end{exbox}

\begin{minted}{lean4}
namespace Propagators

open FlatNat

-- Addition propagator: addFn wrapped in IO — read the inputs, apply
-- the pure monotone function, merge the result into the output cell.
-- When either input is unknown, addFn returns unknown = ⊥, and
-- merging ⊥ changes nothing (write returns false). When an input is
-- conflict, addFn returns conflict — so contradictions propagate.
def addProp (cx cy csum : Cell FlatNat) : PropStep := do
  let x ← cx.read
  let y ← cy.read
  csum.write (FlatNat.addFn x y)

-- Subtraction (used to deduce y = sum - x): subFn wrapped in IO.
def subProp (csum cx cy : Cell FlatNat) : PropStep := do
  let s ← csum.read
  let x ← cx.read
  cy.write (FlatNat.subFn s x)

def exampleNetwork : IO Unit := do
  let cx   ← Cell.new (α := FlatNat)
  let cy   ← Cell.new (α := FlatNat)
  let csum ← Cell.new (α := FlatNat)

  -- Seed known facts
  let _ ← cx.write (.known 3)
  let _ ← csum.write (.known 10)

  let net : Network := {
    steps := #[
      addProp cx cy csum,
      subProp csum cx cy,
      subProp csum cy cx
    ]
  }

  net.run

  let yVal ← cy.read
  match yVal with
  | .known n => IO.println s!"y = {n}"   -- prints: y = 7
  | .unknown  => IO.println "y is unknown"
  | .conflict => IO.println "contradiction!"

end Propagators
\end{minted}

\section{Why the Lattice Is Load-Bearing Here}

Notice several things:
\begin{itemize}
  \item \lc{Cell.write} can be called multiple times from different propagators
    without coordination. The join ensures the final value is consistent with
    all of them.
  \item Running \lc{addProp} when either input is \lc{unknown} correctly does
    nothing: \lc{addFn} returns $\bot$, merging $\bot$ into the target
    changes nothing, and \lc{write} returns \lc{false} --- the network
    simply waits for more information.
  \item The \lc{conflict} element propagates correctly: if two propagators
    derive incompatible values, the join produces \lc{conflict}, which then
    propagates to any cell that depends on this one.
  \item The whole system terminates because \lc{FlatNat} has height 2: after at
    most two increases, any cell stabilises.
\end{itemize}

\section{Exercises}

\begin{exercisebox}[Three-cell constraint]
  Add a \lc{mulProp} propagator for $x \times y = z$ with the \lc{FlatNat}
  lattice. Build a network with $x \times y = 12$, $x = 3$, and verify it
  deduces $y = 4$.
\end{exercisebox}

\begin{exercisebox}[Conflict detection]
  Seed a network with $x = 3$ from one source and $x = 5$ from another.
  Verify that \lc{cx.read} returns \lc{.conflict}.
\end{exercisebox}

% ============================================================
\chapter{The Interval Lattice}

Flat lattices are powerful but coarse: a cell either knows the exact value or
knows nothing. \emph{Interval lattices} let us represent partial numeric
knowledge as a range, enabling \emph{bound propagation} --- the engine behind
many real-world constraint solvers.

\section{Intervals as Partial Information}

An interval $[l, h]$ says: ``the true value lies between $l$ and $h$''. The
wider the interval, the less we know. The narrowest possible interval is a
single point $[n, n]$, meaning we know the value exactly.

Chapter~5's convention therefore dictates the order: more information $=$
\emph{higher}, so narrower intervals must sit \emph{above} wider ones, and
the join $\sqcup$ --- ``merge what two sources tell us'' --- must be the
\emph{intersection}: if one source says $x \in [3, 7]$ and another says
$x \in [5, 12]$, together they say $x \in [5, 7]$.

Two boundary cases complete the picture, and both should look familiar from
\lc{FlatNat}:
\begin{itemize}
  \item $\bot$ must mean ``no information at all''. For intervals that is
    ``the value could be any integer'': the unbounded interval
    $(-\infty, +\infty)$. This forces us to make $\pm\infty$ representable.
  \item $\top$ must mean ``contradiction''. Intersect two disjoint intervals
    --- $[0,1] \sqcup [5,6]$ --- and \emph{no} value remains: the
    \emph{empty} interval plays exactly the role of \lc{conflict}.
\end{itemize}

\begin{defbox}[Interval Lattice]
  Let the domain be the extended integers $\mathbb{Z}_\infty = \mathbb{Z}
  \cup \{-\infty, +\infty\}$. Define:
  \begin{itemize}[nosep]
    \item An \emph{interval} is either $[l, h]$ with $l \leq h$ (where $l$
          may be $-\infty$ and $h$ may be $+\infty$), or the \emph{empty
          interval} $\varnothing$.
    \item The \emph{information order}: $[l_1, h_1] \leq [l_2, h_2] \iff
          l_1 \leq l_2 \wedge h_2 \leq h_1$ (the second interval is contained
          in the first: narrower $=$ more informative $=$ higher), and
          $I \leq \varnothing$ for every interval $I$.
    \item $\bot = (-\infty, +\infty)$: no information.
          $\top = \varnothing$: contradiction --- matching \lc{FlatNat},
          where contradiction is also the \emph{top}.
    \item The \emph{join} (merge two pieces of information) is
          \emph{intersection}:
          $[l_1, h_1] \sqcup [l_2, h_2] =
          [\max(l_1,l_2),\, \min(h_1,h_2)]$, or $\varnothing$ if those
          bounds cross.
  \end{itemize}
\end{defbox}

\begin{figure}[H]
  \centering
  \begin{tikzpicture}
    % Number line at top for reference
    \draw[gray,->] (-0.5,7.8) -- (11,7.8) node[right,gray,font=\small]{value};
    \foreach \x in {0,1,...,10}
      \draw[gray] (\x,7.85)--(\x,7.75)
        node[above,gray,font=\scriptsize]{\x};
    % top = empty at the TOP of the order
    \node[red!70!black,font=\scriptsize] at (5, 6.8)
      {$\top = \varnothing$ --- contradiction (bounds crossed)};
    \filldraw[orange!80!black] (5,5.9) circle (3pt);
    \node[font=\scriptsize,orange!70!black,anchor=west] at (10.3,5.9)
      {$[5,5]$ --- exact};
    \draw[blue!90,line width=4pt] (5,4.8) -- (6,4.8);
    \node[font=\scriptsize,blue!90!black,anchor=west] at (10.3,4.8) {$[5,6]$};
    \draw[blue!75,line width=4pt] (4,3.7) -- (7,3.7);
    \node[font=\scriptsize,blue!75!black,anchor=west] at (10.3,3.7) {$[4,7]$};
    \draw[blue!55,line width=4pt] (2,2.6) -- (8,2.6);
    \node[font=\scriptsize,blue!55!black,anchor=west] at (10.3,2.6) {$[2,8]$};
    \draw[blue!35,line width=4pt] (0,1.5) -- (10,1.5);
    \node[font=\scriptsize,blue!50!black,anchor=west] at (10.3,1.5) {$[0,10]$};
    \draw[blue!20,line width=4pt,densely dashed] (-0.4,0.4) -- (10.4,0.4);
    \node[font=\scriptsize,blue!40!black,anchor=west] at (10.5,0.4)
      {$\bot = (-\infty,+\infty)$};
    \draw[-{Stealth},gray!60,thick] (-1.0,0.4)--(-1.0,5.9)
      node[midway,left,gray,font=\scriptsize,rotate=90,yshift=4pt]
        {narrowing (join $\sqcup$ goes upward)};
  \end{tikzpicture}
  \caption{The interval information order. $\bot = (-\infty,+\infty)$
           (``the value could be anything'') sits at the bottom. Going
           upward, intervals \emph{narrow}: the join $\sqcup$ merges two
           constraints by intersecting them, which is at least as narrow as
           either input. At the very top sits $\top = \varnothing$, the
           empty interval: two constraints with no common value --- a
           contradiction.}
\end{figure}

\section{Bounds: Option Int as the Extended Integers}

How do we represent $\pm\infty$? The cheapest honest encoding in Lean is
\lc{Option Int}: \lc{some n} is the finite bound $n$, and \lc{none} means
``no bound at all'' --- read it as $-\infty$ when it appears as a lower
bound and $+\infty$ when it appears as an upper bound. (\lc{Option} is
Lean's ``maybe'' type; if you know Rust's \lc{Option} or Haskell's
\lc{Maybe}, this is the same thing.)

\begin{minted}{lean4}
namespace Intervals

-- A bound on an integer. `none` encodes "no bound at all":
-- read it as -∞ when it is a lower bound and +∞ when it is an
-- upper bound.
abbrev Bound := Option Int

-- Order on LOWER bounds: none = -∞ is the loosest lower bound.
def loLe : Bound → Bound → Prop
  | none,   _      => True
  | some _, none   => False
  | some a, some b => a ≤ b

-- Order on UPPER bounds: none = +∞ is the loosest upper bound.
def hiLe : Bound → Bound → Prop
  | _,      none   => True
  | none,   some _ => False
  | some a, some b => a ≤ b

-- A lower bound and an upper bound are consistent unless both are
-- finite and crossed (lo > hi).
def ordered : Bound → Bound → Prop
  | some l, some h => l ≤ h
  | _,      _      => True

instance : (l h : Bound) → Decidable (ordered l h)
  | none,   none   => isTrue True.intro
  | none,   some _ => isTrue True.intro
  | some _, none   => isTrue True.intro
  | some l, some h => inferInstanceAs (Decidable (l ≤ h))

-- The tighter of two lower bounds (their max, with none = -∞) …
def loMax : Bound → Bound → Bound
  | none, b => b
  | a, none => a
  | some a, some b => some (max a b)

-- … and the tighter of two upper bounds (their min, with none = +∞).
def hiMin : Bound → Bound → Bound
  | none, b => b
  | a, none => a
  | some a, some b => some (min a b)

end Intervals
\end{minted}

Note that we needed \emph{two} order relations: on lower bounds \lc{none}
is the smallest thing there is; on upper bounds it is the largest. Keeping
the two roles in separate definitions (\lc{loLe} vs.\ \lc{hiLe}, \lc{loMax}
vs.\ \lc{hiMin}) keeps every later proof a brainless case split.

Everything we need about bounds follows from a dozen tiny lemmas. Each is
proved the same way: \lc{cases} on every bound in sight (an \lc{Option} has
just two constructors), then \lc{simp} to unfold the definitions, then
\lc{omega} for whatever integer arithmetic remains.

\begin{minted}{lean4}
namespace Intervals

-- Small lemmas about bounds. Each is a mechanical case split.
theorem loLe_refl (a : Bound) : loLe a a := by
  cases a <;> simp [loLe]

theorem hiLe_refl (a : Bound) : hiLe a a := by
  cases a <;> simp [hiLe]

theorem loLe_trans {a b c : Bound} : loLe a b → loLe b c → loLe a c := by
  cases a <;> cases b <;> cases c <;> simp [loLe] <;> omega

theorem hiLe_trans {a b c : Bound} : hiLe a b → hiLe b c → hiLe a c := by
  cases a <;> cases b <;> cases c <;> simp [hiLe] <;> omega

theorem loLe_antisymm {a b : Bound} : loLe a b → loLe b a → a = b := by
  cases a <;> cases b <;> simp [loLe] <;> omega

theorem hiLe_antisymm {a b : Bound} : hiLe a b → hiLe b a → a = b := by
  cases a <;> cases b <;> simp [hiLe] <;> omega

theorem loLe_max_left (a b : Bound) : loLe a (loMax a b) := by
  cases a <;> cases b <;> simp [loLe, loMax] <;> omega

theorem loLe_max_right (a b : Bound) : loLe b (loMax a b) := by
  cases a <;> cases b <;> simp [loLe, loMax] <;> omega

theorem loMax_le {a b c : Bound} :
    loLe a c → loLe b c → loLe (loMax a b) c := by
  cases a <;> cases b <;> cases c <;> simp [loLe, loMax] <;> omega

theorem hiMin_le_left (a b : Bound) : hiLe (hiMin a b) a := by
  cases a <;> cases b <;> simp [hiLe, hiMin] <;> omega

theorem hiMin_le_right (a b : Bound) : hiLe (hiMin a b) b := by
  cases a <;> cases b <;> simp [hiLe, hiMin] <;> omega

theorem le_hiMin {a b c : Bound} :
    hiLe c a → hiLe c b → hiLe c (hiMin a b) := by
  cases a <;> cases b <;> cases c <;> simp [hiLe, hiMin] <;> omega

-- If [l', h'] is a genuine interval squeezed between l and h,
-- then l and h cannot be crossed either.
theorem ordered_squeeze {l l' h' h : Bound} :
    loLe l l' → ordered l' h' → hiLe h' h → ordered l h := by
  cases l <;> cases l' <;> cases h' <;> cases h <;>
    simp [loLe, hiLe, ordered] <;> omega

end Intervals
\end{minted}

\section{The Interval Lattice in Lean}

An interval is either a pair of bounds --- carrying a proof that they are
not crossed --- or the distinguished contradiction \lc{empty}:

\begin{minted}{lean4}
namespace Intervals

inductive Interval where
  -- All integers between lo and hi (inclusive; none = unbounded).
  -- The proof field rules out crossed bounds like [8, 2].
  | range (lo hi : Bound) (ok : ordered lo hi) : Interval
  -- ⊤: no value is possible — a contradiction.
  | empty : Interval

-- ⊥: the interval (-∞, +∞) — "the value could be any integer".
def Interval.unbounded : Interval := .range none none True.intro

-- The interval [l, h] with finite ends.
def Interval.between (l h : Int) (ok : l ≤ h := by omega) : Interval :=
  .range (some l) (some h) ok

-- Information order: narrower = more informative = HIGHER.
-- b sits above a exactly when b's bounds are at least as tight.
def Interval.le (a b : Interval) : Prop :=
  match a, b with
  | _, .empty                      => True    -- ⊤ is above everything
  | .empty, _                      => False
  | .range l1 h1 _, .range l2 h2 _ => loLe l1 l2 ∧ hiLe h2 h1

-- Join = INTERSECTION: merge two constraints by keeping the tighter
-- bound on each side. If the bounds cross, the constraints are
-- contradictory and the join is ⊤ = empty.
def Interval.sup (a b : Interval) : Interval :=
  match a, b with
  | .empty, _ | _, .empty => .empty
  | .range l1 h1 _, .range l2 h2 _ =>
    if ok : ordered (loMax l1 l2) (hiMin h1 h2)
    then .range (loMax l1 l2) (hiMin h1 h2) ok
    else .empty

end Intervals
\end{minted}

Why does the \lc{range} constructor carry a proof? Without the invariant
$l \leq h$, ``crossed'' pairs like $[8, 2]$ would be legal values distinct
from \lc{empty} --- a swarm of extra contradiction elements. Worse, the
join of two disjoint intervals would have to be one of those crossed pairs
(not \lc{empty}!) to remain a \emph{least} upper bound, and the lattice
laws would quietly break. The proof field makes the illegal states
unrepresentable, and (as you will see in \lc{sup\_le\_thm}) is precisely
the fact that rescues the ``bounds crossed'' case of the join laws. In
\lc{Interval.between}, the proof is an \emph{autoparam}:
\lc{(ok : l ≤ h := by omega)} tells Lean to run \lc{omega} to fill in the
argument whenever the call site omits it, so numeric literals just work:
\lc{Interval.between 2 8}.

The partial-order laws are stated as standalone lemmas and then plugged
into the instance (the \lc{\_thm} suffix avoids confusing a lemma with the
instance field it fills):

\begin{minted}{lean4}
namespace Intervals

theorem Interval.le_refl_thm : ∀ (a : Interval), Interval.le a a := by
  intro a
  cases a with
  | empty => exact True.intro
  | range l h ok => exact ⟨loLe_refl l, hiLe_refl h⟩

theorem Interval.le_antisymm_thm :
    ∀ (a b : Interval), Interval.le a b → Interval.le b a → a = b := by
  intro a b hab hba
  cases a with
  | empty =>
    cases b with
    | empty => rfl
    | range l h ok => exact absurd hab (fun h => h)
  | range l1 h1 ok1 =>
    cases b with
    | empty => exact absurd hba (fun h => h)
    | range l2 h2 ok2 =>
      have hab' : loLe l1 l2 ∧ hiLe h2 h1 := hab
      have hba' : loLe l2 l1 ∧ hiLe h1 h2 := hba
      have hl : l1 = l2 := loLe_antisymm hab'.1 hba'.1
      have hh : h1 = h2 := hiLe_antisymm hba'.2 hab'.2
      subst hl; subst hh
      rfl   -- the `ok` proofs agree by proof irrelevance

theorem Interval.le_trans_thm :
    ∀ (a b c : Interval), Interval.le a b → Interval.le b c → Interval.le a c := by
  intro a b c hab hbc
  cases c with
  | empty => cases a <;> exact True.intro
  | range l3 h3 ok3 =>
    cases b with
    | empty => exact absurd hbc (fun h => h)
    | range l2 h2 ok2 =>
      cases a with
      | empty => exact absurd hab (fun h => h)
      | range l1 h1 ok1 =>
        have hab' : loLe l1 l2 ∧ hiLe h2 h1 := hab
        have hbc' : loLe l2 l3 ∧ hiLe h3 h2 := hbc
        exact ⟨loLe_trans hab'.1 hbc'.1, hiLe_trans hbc'.2 hab'.2⟩

instance : OrderBook.PartialOrder Interval where
  le          := Interval.le
  le_refl     := Interval.le_refl_thm
  le_antisymm := Interval.le_antisymm_thm
  le_trans    := Interval.le_trans_thm

end Intervals
\end{minted}

Three idioms are worth pausing on. First, lines like
\lc{have hab' : loLe l1 l2 ∧ hiLe h2 h1 := hab} do no work at all: once
both arguments are \lc{range}s, \lc{Interval.le a b} \emph{computes} to
that conjunction, and the \lc{have} merely re-states the hypothesis at its
unfolded type so that \lc{.1}/\lc{.2} projections apply. Second,
\lc{exact absurd hab (fun h => h)} dispatches impossible cases: there
\lc{hab}'s type computes to \lc{False}, so the identity function \emph{is}
its refutation. Third, antisymmetry ends with plain \lc{rfl} even though
the two \lc{ok} proof fields are different terms --- proofs of the same
proposition are definitionally equal in Lean (\emph{proof irrelevance}).

Now the join laws. The one interesting moment is the ``bounds crossed''
branch of \lc{sup\_le\_thm}, which is impossible precisely because the
upper bound \lc{c} is itself a well-formed interval squeezed between the
merged bounds --- this is \lc{ordered\_squeeze} doing its job:

\begin{minted}{lean4}
namespace Intervals

theorem Interval.bot_le_thm :
    ∀ (a : Interval), Interval.le Interval.unbounded a := by
  intro a
  cases a with
  | empty => exact True.intro
  | range l h ok => exact ⟨True.intro, True.intro⟩

theorem Interval.le_sup_left_thm :
    ∀ (a b : Interval), Interval.le a (Interval.sup a b) := by
  intro a b
  cases a with
  | empty => cases b <;> exact True.intro
  | range l1 h1 ok1 =>
    cases b with
    | empty => exact True.intro
    | range l2 h2 ok2 =>
      simp only [Interval.sup]
      split
      · exact ⟨loLe_max_left l1 l2, hiMin_le_left h1 h2⟩
      · exact True.intro

theorem Interval.le_sup_right_thm :
    ∀ (a b : Interval), Interval.le b (Interval.sup a b) := by
  intro a b
  cases a with
  | empty => cases b <;> exact True.intro
  | range l1 h1 ok1 =>
    cases b with
    | empty => exact True.intro
    | range l2 h2 ok2 =>
      simp only [Interval.sup]
      split
      · exact ⟨loLe_max_right l1 l2, hiMin_le_right h1 h2⟩
      · exact True.intro

theorem Interval.sup_le_thm :
    ∀ (a b c : Interval),
      Interval.le a c → Interval.le b c → Interval.le (Interval.sup a b) c := by
  intro a b c hac hbc
  cases a with
  | empty => cases c <;> exact hac
  | range l1 h1 ok1 =>
    cases b with
    | empty => cases c <;> exact hbc
    | range l2 h2 ok2 =>
      cases c with
      | empty =>
        simp only [Interval.sup]
        split <;> exact True.intro
      | range l3 h3 ok3 =>
        have hac' : loLe l1 l3 ∧ hiLe h3 h1 := hac
        have hbc' : loLe l2 l3 ∧ hiLe h3 h2 := hbc
        have hlo : loLe (loMax l1 l2) l3 := loMax_le hac'.1 hbc'.1
        have hhi : hiLe h3 (hiMin h1 h2) := le_hiMin hac'.2 hbc'.2
        simp only [Interval.sup]
        split
        · exact ⟨hlo, hhi⟩
        · -- the "crossed bounds" branch is impossible: c is a genuine
          -- interval squeezed between the merged bounds.
          rename_i hbad
          exact absurd (ordered_squeeze hlo ok3 hhi) hbad

instance : Propagators.BoundedJoinSemilattice Interval where
  le          := Interval.le
  le_refl     := Interval.le_refl_thm
  le_antisymm := Interval.le_antisymm_thm
  le_trans    := Interval.le_trans_thm
  bot         := Interval.unbounded
  sup         := Interval.sup
  bot_le      := Interval.bot_le_thm
  le_sup_left  := Interval.le_sup_left_thm
  le_sup_right := Interval.le_sup_right_thm
  sup_le       := Interval.sup_le_thm

-- DecidableEq, needed by Cell.write. Two `range`s are equal iff their
-- bounds are; the `ok` proofs never disagree (proof irrelevance).
instance : DecidableEq Interval := fun a b =>
  match a, b with
  | .empty, .empty => isTrue rfl
  | .empty, .range _ _ _ => isFalse (fun h => Interval.noConfusion h)
  | .range _ _ _, .empty => isFalse (fun h => Interval.noConfusion h)
  | .range l1 h1 _, .range l2 h2 _ =>
    if h : l1 = l2 ∧ h1 = h2 then
      isTrue (by obtain ⟨rfl, rfl⟩ := h; rfl)
    else
      isFalse (fun heq => by cases heq; exact h ⟨rfl, rfl⟩)

end Intervals
\end{minted}

\section{Interval Arithmetic Propagators}

To use \lc{Interval} with the \lc{Cell} infrastructure from Chapter~6 we
needed exactly two instances --- \lc{BoundedJoinSemilattice Interval} (so
the scheduler can merge values with $\sqcup$ and initialise cells to
$\bot$) and \lc{DecidableEq Interval} (so \lc{Cell.write} can detect
whether the value actually changed) --- and both are now in place. What
remains is the arithmetic itself, which must handle unbounded ends: a
finite number plus ``anything at all'' is again ``anything at all''.

\begin{minted}{lean4}
namespace Intervals

-- Adding bounds: if either end is unbounded, so is the sum.
def addBound : Bound → Bound → Bound
  | some a, some b => some (a + b)
  | _,      _      => none

def subBound : Bound → Bound → Bound
  | some a, some b => some (a - b)
  | _,      _      => none

theorem ordered_addBound {l1 h1 l2 h2 : Bound} :
    ordered l1 h1 → ordered l2 h2 →
    ordered (addBound l1 l2) (addBound h1 h2) := by
  cases l1 <;> cases l2 <;> cases h1 <;> cases h2 <;>
    simp [ordered, addBound] <;> omega

theorem ordered_subBound {l1 h1 l2 h2 : Bound} :
    ordered l1 h1 → ordered l2 h2 →
    ordered (subBound l1 h2) (subBound h1 l2) := by
  cases l1 <;> cases l2 <;> cases h1 <;> cases h2 <;>
    simp [ordered, subBound] <;> omega

-- [l1,h1] + [l2,h2] = [l1+l2, h1+h2]
def Interval.add : Interval → Interval → Interval
  | .empty, _ | _, .empty => .empty
  | .range l1 h1 ok1, .range l2 h2 ok2 =>
    .range (addBound l1 l2) (addBound h1 h2) (ordered_addBound ok1 ok2)

-- [l1,h1] - [l2,h2] = [l1-h2, h1-l2]
def Interval.sub : Interval → Interval → Interval
  | .empty, _ | _, .empty => .empty
  | .range l1 h1 ok1, .range l2 h2 ok2 =>
    .range (subBound l1 h2) (subBound h1 l2) (ordered_subBound ok1 ok2)

end Intervals
\end{minted}

Note the asymmetry in subtraction: the smallest possible value of $x - y$
pairs the smallest $x$ with the \emph{largest} $y$, hence
$[l_1 - h_2,\; h_1 - l_2]$. The proof obligations on the results are
discharged by the two \lc{ordered\_*} lemmas, whose hypotheses are exactly
the \lc{ok} fields of the argument intervals.

\section{A Worked Example: Bound Propagation}

\begin{exbox}[{$x + y = 10$, $x \in [2, 8]$}]
  What is the tightest interval for $y$?

  From $x + y = 10$ we get $y = 10 - x$. Since $x \in [2, 8]$, we have
  $y \in [10 - 8, 10 - 2] = [2, 8]$. The propagator should derive this
  automatically.
\end{exbox}

\begin{minted}{lean4}
namespace Intervals

-- Propagator: sum = x + y (bidirectional: also narrows x and y)
def addConstraint
    (cx cy csum : Propagators.Cell Interval) : Array Propagators.PropStep :=
  #[
    -- Forward: sum ← x + y
    do
      let x ← cx.read; let y ← cy.read
      csum.write (Interval.add x y),
    -- Backward: x ← sum - y
    do
      let s ← csum.read; let y ← cy.read
      cx.write (Interval.sub s y),
    -- Backward: y ← sum - x
    do
      let s ← csum.read; let x ← cx.read
      cy.write (Interval.sub s x)
  ]

def showLo : Bound → String
  | none   => "-∞"
  | some n => toString n

def showHi : Bound → String
  | none   => "+∞"
  | some n => toString n

def Interval.show : Interval → String
  | .empty => "∅ (contradiction)"
  | .range l h _ => s!"[{showLo l}, {showHi h}]"

def intervalExample : IO Unit := do
  let cx   ← Propagators.Cell.new (α := Interval)
  let cy   ← Propagators.Cell.new (α := Interval)
  let csum ← Propagators.Cell.new (α := Interval)

  -- x ∈ [2, 8] and x + y = 10 (i.e. sum ∈ [10, 10])
  let _ ← cx.write (Interval.between 2 8)
  let _ ← csum.write (Interval.between 10 10)

  Propagators.runToFixpoint (addConstraint cx cy csum)

  IO.println s!"x   ∈ {Interval.show (← cx.read)}"
  IO.println s!"y   ∈ {Interval.show (← cy.read)}"
  IO.println s!"sum ∈ {Interval.show (← csum.read)}"

#eval intervalExample

end Intervals
\end{minted}

Running it, the network --- forward \emph{and} backward propagators
installed --- reaches quiescence and prints:

\begin{minted}{text}
x   ∈ [2, 8]
y   ∈ [2, 8]
sum ∈ [10, 10]
\end{minted}

\begin{notebox}
  \textbf{The join in action.}
  Pass 1: the forward propagator reads $x = [2,8]$ and $y = \bot$, computes
  $x + y = (-\infty, +\infty)$, and merges it into $c_\text{sum}$ --- no
  change, since $[10,10] \sqcup \bot = [10,10]$. The backward propagator
  for $y$ computes $\mathrm{sum} - x = [10-8,\, 10-2] = [2,8]$ and merges
  it into $c_y$: $\bot \sqcup [2,8] = [2,8]$, so $c_y$ narrows.
  Pass 2: the forward propagator now computes $x + y = [4,16]$ and merges:
  $[10,10] \sqcup [4,16] = [\max(10,4), \min(10,16)] = [10,10]$. Nothing
  widens and nothing is lost --- the join is intersection, so a looser
  derived constraint disappears into a tighter stored one. The pass
  changes nothing, the queue drains, and the network is quiescent. (With a
  hull-shaped ``widening'' join instead, that same forward write would
  smear $[10,10]$ out to $[4,16]$, the backward propagators would widen
  $x$ and $y$ in response, and the network would never quiesce.)
\end{notebox}

\begin{notebox}
  \textbf{Termination, honestly.}
  Chapter~5's convergence theorem assumes the ascending chain condition,
  and the interval lattice does \emph{not} satisfy it:
  $(-\infty,+\infty) < [0,+\infty) < [1,+\infty) < \cdots$ climbs forever.
  Quiescence above is therefore a property of \emph{this constraint
  system}, not of the lattice: these propagators happen to hit their least
  fixed point after two passes. A system like $y = x + 1,\ x = y + 1$
  seeded with $x \in [0,+\infty)$ really would tighten forever; industrial
  interval solvers add a \emph{widening} operator that jumps a bound
  straight to $\pm\infty$ after a few rounds to force termination. The ACC
  is sufficient for termination, never necessary.
\end{notebox}

\section{Exercises}

\begin{exercisebox}[The hull is the meet]
  The join of the interval lattice is intersection. Show that the lattice
  also has all \emph{meets}: define \lc{Interval.inf}, the \emph{hull}
  (smallest interval containing both arguments --- the looser bound on each
  side, with \lc{Interval.inf .empty x = x}), and prove the analogue of
  \lc{inf\_le\_left}. Why does the hull, which \emph{loses} information,
  sit \emph{below} its arguments in the information order?
\end{exercisebox}

\begin{exercisebox}[Multiplication intervals]
  Implement \lc{Interval.mul}. Be careful on two fronts: multiplying
  finite intervals with negative ends requires considering all four corner
  products, and unbounded ends can flip sign --- e.g.\
  $[-2,-1] \times [0,+\infty) = (-\infty, 0]$ --- so \lc{mulBound} needs
  genuine sign analysis rather than a one-line match.
\end{exercisebox}

% ============================================================
\part{Capstone Projects}

% ============================================================
\chapter{Capstone I: A Sudoku Solver}

\section{What is Sudoku?}

Sudoku is a logic puzzle played on a $9 \times 9$ grid. The grid is divided into
nine $3 \times 3$ \emph{boxes}. Some cells are pre-filled with digits; your task
is to fill in the rest.

\begin{figure}[H]
  \centering
  \begin{tikzpicture}[scale=0.55]
    % Outer border
    \draw[line width=2pt] (0,0) rectangle (9,9);
    % Thin grid lines
    \foreach \x in {1,...,8} \draw[gray!50, line width=0.4pt] (\x,0)--(\x,9);
    \foreach \y in {1,...,8} \draw[gray!50, line width=0.4pt] (0,\y)--(9,\y);
    % Thick box borders
    \foreach \x in {3,6} \draw[line width=1.5pt] (\x,0)--(\x,9);
    \foreach \y in {3,6} \draw[line width=1.5pt] (0,\y)--(9,\y);
    % Given digits (example puzzle)
    \foreach \row/\col/\val in {
      9/1/5, 9/2/3, 9/5/7,
      8/1/6, 8/4/1, 8/5/9, 8/6/5,
      7/2/9, 7/3/8, 7/8/6,
      6/1/8, 6/5/6, 6/9/3,
      5/1/4, 5/4/8, 5/6/3, 5/9/1,
      4/1/7, 4/5/2, 4/9/6,
      3/2/6, 3/7/2, 3/8/8,
      2/4/4, 2/5/1, 2/6/9, 2/9/5,
      1/5/8, 1/8/7, 1/9/9
    } {
      \node[font=\small\bfseries] at (\col-0.5, \row-0.5) {\val};
    }
    % Solved digits (lighter) — from the puzzle's unique solution
    \foreach \row/\col/\val in {
      9/3/4, 9/4/6, 9/6/8, 9/7/9, 9/8/1, 9/9/2,
      8/2/7, 8/3/2, 8/7/3, 8/8/4, 8/9/8,
      7/1/1, 7/4/3, 7/5/4, 7/6/2, 7/7/5, 7/9/7
    } {
      \node[font=\small, gray!70] at (\col-0.5, \row-0.5) {\val};
    }
    % Box labels
    \node[font=\tiny, gray] at (1.5, 7.85) {box};
    \node[font=\tiny, gray] at (1.5, 7.5) {(top-left)};
    % Row/col annotations
    \foreach \c in {1,...,9}
      \node[font=\tiny, gray] at (\c-0.5, -0.35) {c\c};
    \foreach \r/\rl in {1/9,2/8,3/7,4/6,5/5,6/4,7/3,8/2,9/1}
      \node[font=\tiny, gray] at (-0.4, \r-0.5) {r\rl};
  \end{tikzpicture}
  \caption{A Sudoku puzzle. Bold digits are givens; light grey digits show
           a few solved cells. The thick lines divide the grid into nine
           $3 \times 3$ boxes.}
\end{figure}

\textbf{The three rules.} Every digit $1$--$9$ must appear \emph{exactly once}
in each of:
\begin{itemize}[nosep]
  \item each \textbf{row} (9 rows, one per horizontal strip),
  \item each \textbf{column} (9 columns, one per vertical strip),
  \item each \textbf{box} (9 boxes, each a $3 \times 3$ sub-grid).
\end{itemize}

That is the entire rule set. No arithmetic is involved --- the digits could
equally well be colours, letters, or symbols. The puzzle is purely about
placement constraints.

\textbf{How to solve by hand.}
The standard technique is \emph{elimination}: for each empty cell, maintain the
set of digits that could still legally go there given the filled cells in the
same row, column, and box. If a cell's candidate set shrinks to a single digit,
that digit is determined. Filling it in eliminates that digit from the candidate
sets of neighbouring cells, which may in turn determine more cells. This
cascading process is exactly what a propagator network formalises.

\begin{exbox}[A simple deduction]
  Suppose row 1 already contains $1, 2, 3, 4, 6, 7, 8, 9$ in various cells.
  The only missing digit is $5$, so the one remaining empty cell in row 1
  must be $5$ --- no other placement is consistent with the row constraint.
\end{exbox}

\begin{exbox}[Elimination across constraints]
  Suppose a cell sits in row 3, column 5, box 5 (the centre box). Row 3 already
  contains $\{1,3,4,6,7\}$, column 5 contains $\{2,3,6,8\}$, and box 5 contains
  $\{1,2,6,9\}$. The candidates for this cell are all digits \emph{not} ruled
  out by any of its three units:
  $$\{1,\ldots,9\} \setminus \bigl(\{1,3,4,6,7\} \cup \{2,3,6,8\}
    \cup \{1,2,6,9\}\bigr) = \{5\}.$$
  One candidate remains, so this cell must be $5$. Each elimination step is
  an intersection with the complement of a forbidden set --- computed with
  \lc{CandidateSet.inter}, which is exactly the join of the information
  order.
\end{exbox}

\textbf{Why propagator networks are a natural fit.}
Each constraint (row, column, box) is a propagator that reads the determined
cells in its unit and writes the resulting eliminations to undetermined cells.
The network runs until quiescence --- no propagator can make further
eliminations. This is guaranteed to terminate because candidate sets can only
shrink, and a finite set can only shrink finitely many times. Whether it
\emph{solves} the puzzle completely depends on the puzzle; harder puzzles
require search in addition to propagation.

\section{The Domain Lattice}

Sudoku is a constraint satisfaction problem: each of 81 cells must receive a
digit $1$--$9$ subject to row, column, and box uniqueness constraints.
Propagator networks handle this beautifully. We build a complete solver
using nothing but lattice-ordered propagation.

Each Sudoku cell can be modelled as a set of \emph{candidate digits}. Initially
every cell can hold any of $\{1,\ldots,9\}$. Constraints eliminate candidates.

\begin{defbox}[Candidate Set Lattice]
  Let $D = \{1,\ldots,9\}$. The \emph{candidate-set lattice} is the powerset
  $\mathcal{P}(D)$ ordered by \emph{reverse inclusion}:
  $$S \leq T \iff T \subseteq S.$$
  The bottom $\bot = D$ (no information: any digit possible),
  the top $\top = \varnothing$ (contradiction: no digit possible).
  The join is intersection: $S \sqcup T = S \cap T$
  (the smallest set of candidates consistent with both sources).
\end{defbox}

\begin{figure}[h]
  \centering
  \begin{tikzpicture}[node distance=1cm]
    \node[node,minimum size=1.1cm] (bot) {\tiny $\{1..9\}$};
    \node[node,minimum size=1.1cm] (mid1) [above left=1.1cm and 1.8cm of bot] {\tiny $\{1,2,3\}$};
    \node[node,minimum size=1.1cm] (mid2) [above=1.4cm of bot]               {\tiny $\{4,5\}$};
    \node[node,minimum size=1.1cm] (mid3) [above right=1.1cm and 1.8cm of bot]{\tiny $\{7,8,9\}$};
    \node[node,minimum size=1.1cm] (hi1)  [above left=1.1cm and 0.8cm of mid1]{\tiny $\{2\}$};
    \node[node,minimum size=1.1cm] (hi2)  [above=1.1cm of mid2]               {\tiny $\{5\}$};
    \node[node,minimum size=1.1cm] (hi3)  [above right=1.1cm and 0.8cm of mid3]{\tiny $\{8\}$};
    \node[node,minimum size=1.1cm] (top)  [above=1.4cm of hi2]                {\tiny $\varnothing$};
    \draw[hasseedge] (bot)--(mid1); \draw[hasseedge] (bot)--(mid2);
    \draw[hasseedge] (bot)--(mid3);
    \draw[hasseedge] (mid1)--(hi1); \draw[hasseedge] (mid2)--(hi2);
    \draw[hasseedge] (mid3)--(hi3);
    \draw[hasseedge] (hi1)--(top);  \draw[hasseedge] (hi2)--(top);
    \draw[hasseedge] (hi3)--(top);
    \node[right=0.2cm of top, gray, font=\small\itshape] {$\top$: contradiction};
    \node[right=0.2cm of bot, gray, font=\small\itshape] {$\bot$: no info};
  \end{tikzpicture}
  \caption{A fragment of the candidate-set lattice (reverse inclusion order).}
\end{figure}

This is precisely the \emph{dual} powerset lattice from Exercise~1.4!

\begin{minted}{lean4}
namespace Sudoku

-- Digits 1..9, represented as Fin 9 (0..8 internally, displayed as 1..9)
abbrev Digit := Fin 9

-- A candidate set: which digits are still possible for this cell
-- We represent it as a Bool array of length 9
structure CandidateSet where
  present : Fin 9 → Bool

def CandidateSet.full : CandidateSet :=
  { present := fun _ => true }

def CandidateSet.empty : CandidateSet :=
  { present := fun _ => false }

def CandidateSet.singleton (d : Digit) : CandidateSet :=
  { present := fun i => i == d }

def CandidateSet.remove (s : CandidateSet) (d : Digit) : CandidateSet :=
  { present := fun i => s.present i && !(i == d) }

def CandidateSet.inter (s t : CandidateSet) : CandidateSet :=
  { present := fun i => s.present i && t.present i }

def CandidateSet.union (s t : CandidateSet) : CandidateSet :=
  { present := fun i => s.present i || t.present i }

-- The information order: s ≤ t iff t ⊆ s (t is more specific)
def CandidateSet.le (s t : CandidateSet) : Prop :=
  ∀ i : Digit, t.present i = true → s.present i = true

-- Antisymmetry: compare the two indicator functions pointwise
-- (rename_i names the functions exposed by `cases`).
instance : OrderBook.PartialOrder CandidateSet where
  le          := CandidateSet.le
  le_refl     := fun s i h => h
  le_antisymm := fun s t h1 h2 =>
    show s = t from by
      cases s; cases t
      congr 1; funext i
      rename_i sp tp
      cases hi : tp i
      · cases hs : sp i
        · rfl
        · exact absurd (h2 i hs) (by simp [hi])
      · exact h1 i hi
  le_trans    := fun s t u hst htu i hi => hst i (htu i hi)

-- Join = intersection (combines constraints).
-- `Bool.and_eq_true : ((a && b) = true) = (a = true ∧ b = true)` is
-- an equality of propositions, so we rewrite with `▸`.
instance : Propagators.BoundedJoinSemilattice CandidateSet where
  le          := CandidateSet.le
  le_refl     := fun s i h => h
  le_antisymm := OrderBook.PartialOrder.le_antisymm (α := CandidateSet)
  le_trans    := OrderBook.PartialOrder.le_trans (α := CandidateSet)
  bot         := CandidateSet.full
  sup         := CandidateSet.inter    -- ⊔ = ∩ in information order!
  bot_le      := fun s i _ => rfl     -- full ≤ s since s ⊆ full
  le_sup_left  := fun s t i h => (Bool.and_eq_true _ _ ▸ h).1
  le_sup_right := fun s t i h => (Bool.and_eq_true _ _ ▸ h).2
  sup_le       := fun s t u hsu htu i h =>
    Bool.and_eq_true _ _ ▸ (⟨hsu i h, htu i h⟩ : _ ∧ _)

instance : DecidableEq CandidateSet :=
  fun s t =>
    if h : ∀ i, s.present i = t.present i
    then isTrue  (by cases s; cases t; congr 1; funext i; exact h i)
    else isFalse (fun heq => h (fun i => by rw [heq]))

end Sudoku
\end{minted}

\section{Representing the Grid}

\begin{minted}{lean4}
namespace Sudoku

-- 81 cells indexed by row and column (both Fin 9)
abbrev Grid := Fin 9 → Fin 9 → Propagators.Cell CandidateSet

-- The `[·]!` index operators need an `Inhabited` fallback value;
-- `IO.Ref` provides none, so we allocate one dummy cell and bring a
-- local `Inhabited` instance into scope.
def makeGrid : IO Grid := do
  let mut rows : Array (Array (Propagators.Cell CandidateSet)) := #[]
  for _ in [:9] do
    let mut row : Array (Propagators.Cell CandidateSet) := #[]
    for _ in [:9] do
      let c ← Propagators.Cell.new (α := CandidateSet)
      row := row.push c
    rows := rows.push row
  let dummy ← Propagators.Cell.new (α := CandidateSet)
  have : Inhabited (Propagators.Cell CandidateSet) := ⟨dummy⟩
  return fun r c => rows[r.val]![c.val]!

-- Seed a given digit into a cell
def seedCell (grid : Grid) (r c : Fin 9) (d : Digit) : IO Unit := do
  let _ ← (grid r c).write (CandidateSet.singleton d)

end Sudoku
\end{minted}

\section{Constraint Propagators}

Sudoku has three families of constraints: row, column, and $3\times 3$ box
uniqueness. Each is expressed as a propagator that, given a cell known to hold
digit $d$, removes $d$ from all peers.

\begin{minted}{lean4}
namespace Sudoku

-- Note the loop syntax: `for h : i in [:9]` binds, alongside i, a
-- proof `h` that i lies in the range; `h.upper : i < 9` is exactly
-- what the `Fin 9` constructor ⟨i, _⟩ needs.

-- If cell (r,c) is determined to be d, remove d from all other cells in row r
def rowPropagator (grid : Grid) (r c : Fin 9) : Propagators.PropStep := do
  let cands ← (grid r c).read
  -- Check if this cell is a singleton
  let det : Option Digit := Id.run do
    let mut found : Option Digit := none
    let mut count := 0
    for h : i in [:9] do
      let i' : Digit := ⟨i, h.upper⟩
      if cands.present i' then
        found := some i'; count := count + 1
    if count = 1 then found else none
  match det with
  | none   => return false
  | some d =>
    let mut changed := false
    for h : j in [:9] do
      let j' : Fin 9 := ⟨j, h.upper⟩
      if j' ≠ c then
        -- We re-read and merge
        let old ← (grid r j').read
        let upd := CandidateSet.inter old { present := fun i => !(i == d) }
        let c ← (grid r j').write upd
        if c then changed := true
    return changed

-- Analogous: column propagator
def colPropagator (grid : Grid) (r c : Fin 9) : Propagators.PropStep := do
  let cands ← (grid r c).read
  let det : Option Digit := Id.run do
    let mut found : Option Digit := none
    let mut count := 0
    for h : i in [:9] do
      let i' : Digit := ⟨i, h.upper⟩
      if cands.present i' then found := some i'; count := count + 1
    if count = 1 then found else none
  match det with
  | none => return false
  | some d =>
    let mut changed := false
    for h : i in [:9] do
      let i' : Fin 9 := ⟨i, h.upper⟩
      if i' ≠ r then
        let old ← (grid i' c).read
        let upd := CandidateSet.inter old { present := fun k => !(k == d) }
        let ch ← (grid i' c).write upd
        if ch then changed := true
    return changed

-- Box propagator (which 3×3 box contains (r,c)?)
def boxPropagator (grid : Grid) (r c : Fin 9) : Propagators.PropStep := do
  let cands ← (grid r c).read
  let det : Option Digit := Id.run do
    let mut found : Option Digit := none; let mut count := 0
    for h : i in [:9] do
      let i' : Digit := ⟨i, h.upper⟩
      if cands.present i' then found := some i'; count := count + 1
    if count = 1 then found else none
  match det with
  | none => return false
  | some d =>
    let boxR := r.val / 3 * 3
    let boxC := c.val / 3 * 3
    let mut changed := false
    for hr : dr in [:3] do
      for hc : dc in [:3] do
        let r' : Fin 9 := ⟨boxR + dr, by have h1 : dr < 3 := hr.upper; have := r.isLt; omega⟩
        let c' : Fin 9 := ⟨boxC + dc, by have h2 : dc < 3 := hc.upper; have := c.isLt; omega⟩
        if !(r' = r && c' = c) then
          let old ← (grid r' c').read
          let upd := CandidateSet.inter old { present := fun k => !(k == d) }
          let ch ← (grid r' c').write upd
          if ch then changed := true
    return changed

end Sudoku
\end{minted}

\section{Assembling and Running the Solver}

\begin{minted}{lean4}
namespace Sudoku

def buildNetwork (grid : Grid) : Array Propagators.PropStep := Id.run do
  let mut steps : Array Propagators.PropStep := #[]
  for hr : r in [:9] do
    for hc : c in [:9] do
      let r' : Fin 9 := ⟨r, hr.upper⟩
      let c' : Fin 9 := ⟨c, hc.upper⟩
      steps := steps.push (rowPropagator grid r' c')
      steps := steps.push (colPropagator grid r' c')
      steps := steps.push (boxPropagator grid r' c')
  return steps

-- Print the grid (showing the unique digit or '?' for unsolved cells)
def printGrid (grid : Grid) : IO Unit := do
  for hr : r in [:9] do
    let mut row := ""
    for hc : c in [:9] do
      let r' : Fin 9 := ⟨r, hr.upper⟩
      let c' : Fin 9 := ⟨c, hc.upper⟩
      let cands ← (grid r' c').read
      let mut digit := '?'
      let mut count := 0
      for h : i in [:9] do
        let i' : Digit := ⟨i, h.upper⟩
        if cands.present i' then digit := Char.ofNat (i + 49); count := count + 1
      row := row ++ (if count = 1 then s!"{digit}" else "?")
      if c % 3 == 2 && c < 8 then row := row ++ "|"
    IO.println row
    if r % 3 == 2 && r < 8 then IO.println "-----------"

-- Example: a simple Sudoku puzzle
-- (0 = unknown, digits are given clues)
def examplePuzzle : Array (Array Nat) := #[
  #[5, 3, 0, 0, 7, 0, 0, 0, 0],
  #[6, 0, 0, 1, 9, 5, 0, 0, 0],
  #[0, 9, 8, 0, 0, 0, 0, 6, 0],
  #[8, 0, 0, 0, 6, 0, 0, 0, 3],
  #[4, 0, 0, 8, 0, 3, 0, 0, 1],
  #[7, 0, 0, 0, 2, 0, 0, 0, 6],
  #[0, 6, 0, 0, 0, 0, 2, 8, 0],
  #[0, 0, 0, 4, 1, 9, 0, 0, 5],
  #[0, 0, 0, 0, 8, 0, 0, 7, 9]
]

-- The dependent `if h : 1 ≤ d ∧ d ≤ 9` hands omega the bound it
-- needs to build the `Fin 9` seed value `⟨d - 1, _⟩`.
def solveSudoku : IO Unit := do
  let grid ← makeGrid
  -- Seed known digits
  for hr : r in [:9] do
    for hc : c in [:9] do
      let d := examplePuzzle[r]![c]!
      if h : 1 ≤ d ∧ d ≤ 9 then
        let r' : Fin 9 := ⟨r, hr.upper⟩
        let c' : Fin 9 := ⟨c, hc.upper⟩
        seedCell grid r' c' ⟨d - 1, by omega⟩
  -- Build and run the propagator network
  let steps := buildNetwork grid
  Propagators.runToFixpoint steps
  -- Print the result
  IO.println "Solution:"
  printGrid grid

end Sudoku
\end{minted}

Running \lc{#eval solveSudoku}, propagation alone drives every cell to a
singleton --- this puzzle needs no search --- and prints the grid's unique
solution:

\begin{minted}{text}
Solution:
534|678|912
672|195|348
198|342|567
-----------
859|761|423
426|853|791
713|924|856
-----------
961|537|284
287|419|635
345|286|179
\end{minted}

\section{Why This Works: The Lattice Perspective}

\begin{itemize}
  \item Each cell starts at $\bot = \{1,\ldots,9\}$: no information.
  \item Every propagator only writes using \lc{CandidateSet.inter} (the join
    in the information order), which can only \emph{reduce} candidate sets ---
    i.e., move \emph{up} in the information order.
  \item The network terminates because the candidate-set lattice is finite and
    the cells only move upward.
  \item Reaching $\top = \varnothing$ in any cell signals a contradiction.
  \item A solved cell is one that has reached a singleton $\{d\}$.
\end{itemize}

The elegant insight: \emph{solving Sudoku is computing the least fixed point of
the combined propagation function}, where the partial order is the information
order on candidate sets. Knaster--Tarski guarantees it exists; our scheduler
computes it.

\begin{exercisebox}[Naked pair]
  A \emph{naked pair} occurs when two cells in the same unit both have exactly
  the same two candidates $\{a, b\}$. Then $a$ and $b$ can be removed from all
  other cells in that unit. Implement this as an additional propagator.
\end{exercisebox}

\begin{exercisebox}[Contradiction detection]
  Add a \lc{contradictionProp} that checks every cell; if any cell reaches
  $\top = \varnothing$, print a diagnostic and halt. Why is this detection
  sound (i.e., why can't a cell reach $\varnothing$ unless the puzzle is truly
  unsatisfiable)?
\end{exercisebox}

% ============================================================
\chapter{Capstone II: Type Inference via Propagation}

Our second capstone takes us into programming language theory.
Type inference --- the process of automatically deducing types of expressions
--- is naturally formulated as constraint propagation on a type lattice.

\section{A Simple Expression Language}

We work with a fragment of the lambda calculus extended with integers and
booleans.

\begin{minted}{lean4}
namespace TypeInfer

-- Expressions
inductive Expr where
  | intLit  (n : Int)        : Expr
  | boolLit (b : Bool)       : Expr
  | var     (x : String)     : Expr
  | add     (e1 e2 : Expr)   : Expr
  | ifExpr  (cond t f : Expr): Expr
  | lam     (x : String) (body : Expr) : Expr
  | app     (fn arg : Expr)  : Expr
  | letExpr (x : String) (e body : Expr) : Expr

end TypeInfer
\end{minted}

\section{The Type Lattice}

Rather than types being a simple set, we arrange them in a lattice that
captures the degree to which a type is \emph{known}.

\begin{defbox}[Flat Type Lattice]
  Define the type information domain:
  \begin{itemize}[nosep]
    \item $\bot$ (\lc{Unknown}): no type information yet.
    \item \lc{TInt}, \lc{TBool}: concrete base types.
    \item \lc{TFun s t}: function type, parameterised by the type information
          for the domain and codomain (which are themselves in the lattice).
    \item $\top$ (\lc{TypeError}): contradiction (incompatible constraints).
  \end{itemize}
  The order is: $\bot \leq$ everything, everything $\leq \top$, and
  base types are only $\leq$ themselves (and $\bot / \top$).
  Function types: \lc{TFun s1 t1 ≤ TFun s2 t2 iff s1 ≤ s2 ∧ t1 ≤ t2}.
\end{defbox}

\begin{minted}{lean4}
namespace TypeInfer

inductive Ty where
  | unknown            : Ty    -- ⊥
  | int                : Ty
  | bool               : Ty
  | fun_ (dom cod : Ty): Ty    -- dom → cod
  | typeError          : Ty    -- ⊤

-- Unification of two types: compute their join
def Ty.unify : Ty → Ty → Ty
  | .unknown, t          => t
  | t, .unknown          => t
  | .typeError, _        => .typeError
  | _, .typeError        => .typeError
  | .int, .int           => .int
  | .bool, .bool         => .bool
  | .fun_ d1 c1, .fun_ d2 c2 => .fun_ (Ty.unify d1 d2) (Ty.unify c1 c2)
  | _, _                 => .typeError   -- incompatible types

-- The information order
def Ty.le : Ty → Ty → Prop
  | .unknown, _            => True
  | _, .typeError          => True
  | .int, .int             => True
  | .bool, .bool           => True
  | .fun_ d1 c1, .fun_ d2 c2 => Ty.le d1 d2 ∧ Ty.le c1 c2
  | _, _                   => False

-- The order laws are proved as standalone lemmas by structural
-- induction on the first argument (the recursive `fun_` constructor
-- needs genuine induction hypotheses), then plugged into the
-- instances as terms.

theorem Ty.le_refl_thm : ∀ (t : Ty), Ty.le t t := by
  intro t
  induction t <;> simp [Ty.le, *]

theorem Ty.le_antisymm_thm : ∀ (a b : Ty), Ty.le a b → Ty.le b a → a = b := by
  intro a
  induction a with
  | fun_ d c ihd ihc =>
    intro b h1 h2
    cases b <;> simp_all [Ty.le]
    exact ⟨ihd _ h1.1 h2.1, ihc _ h1.2 h2.2⟩
  | _ => intro b h1 h2; cases b <;> simp_all [Ty.le]

theorem Ty.le_trans_thm : ∀ (a b c : Ty), Ty.le a b → Ty.le b c → Ty.le a c := by
  intro a
  induction a with
  | fun_ d c ihd ihc =>
    intro b c' h1 h2
    cases b <;> cases c' <;> simp_all [Ty.le]
    exact ⟨ihd _ _ h1.1 h2.1, ihc _ _ h1.2 h2.2⟩
  | _ => intro b c' h1 h2; cases b <;> cases c' <;> simp_all [Ty.le]

theorem Ty.le_sup_left_thm : ∀ (a b : Ty), Ty.le a (Ty.unify a b) := by
  intro a
  induction a with
  | fun_ d c ihd ihc =>
    intro b
    cases b <;> simp [Ty.le, Ty.unify]
    · exact ⟨Ty.le_refl_thm d, Ty.le_refl_thm c⟩
    · exact ⟨ihd _, ihc _⟩
  | _ => intro b; cases b <;> simp [Ty.le, Ty.unify]

theorem Ty.le_sup_right_thm : ∀ (a b : Ty), Ty.le b (Ty.unify a b) := by
  intro a
  induction a with
  | fun_ d c ihd ihc =>
    intro b
    cases b <;> simp [Ty.le, Ty.unify]
    · exact ⟨ihd _, ihc _⟩
  | _ => intro b; cases b <;> simp [Ty.le, Ty.unify, Ty.le_refl_thm]

theorem Ty.sup_le_thm :
    ∀ (a b c : Ty), Ty.le a c → Ty.le b c → Ty.le (Ty.unify a b) c := by
  intro a
  induction a with
  | fun_ d c ihd ihc =>
    intro b c' h1 h2
    cases b <;> cases c' <;> simp_all [Ty.le, Ty.unify]
  | _ => intro b c' h1 h2; cases b <;> cases c' <;> simp_all [Ty.le, Ty.unify]

instance : OrderBook.PartialOrder Ty where
  le          := Ty.le
  le_refl     := Ty.le_refl_thm
  le_antisymm := Ty.le_antisymm_thm
  le_trans    := Ty.le_trans_thm

instance : Propagators.BoundedJoinSemilattice Ty where
  le          := Ty.le
  le_refl     := Ty.le_refl_thm
  le_antisymm := Ty.le_antisymm_thm
  le_trans    := Ty.le_trans_thm
  bot         := .unknown
  sup         := Ty.unify
  bot_le      := fun _ => trivial
  le_sup_left  := Ty.le_sup_left_thm
  le_sup_right := Ty.le_sup_right_thm
  sup_le       := Ty.sup_le_thm

deriving instance DecidableEq for Ty

end TypeInfer
\end{minted}

Why lemmas instead of inline tactic blocks? \lc{Ty} is our first
\emph{recursive} lattice: \lc{fun\_ dom cod} contains two smaller
\lc{Ty}s, so every law needs genuine induction hypotheses about those
components, and the tidiest way to manage them is \lc{induction a with}
naming the hypotheses (\lc{ihd}, \lc{ihc}) explicitly. The wildcard arm
\lc{| \_ =>} covers the four non-recursive constructors at once.

\section{Type Constraint Generation}

We generate constraints by walking the expression tree. Each sub-expression
gets a fresh \emph{type cell}. Constraints are propagators that merge type
information into these cells. (A syntactic trap to avoid: the propagators
are appended with \lc{steps.modify (·.push ...)}, and the anonymous-function
dot must \emph{not} be followed by a space --- \lc{(· .push ...)} parses as
applying the placeholder to \lc{.push} and fails with ``function
expected''.)

\begin{minted}{lean4}
namespace TypeInfer

-- A typing context maps variable names to type cells
abbrev Context := List (String × Propagators.Cell Ty)

def Context.lookup (ctx : Context) (x : String) :
    Option (Propagators.Cell Ty) :=
  ctx.find? (fun p => p.1 == x) |>.map (·.2)

-- Generate type cells and propagators for an expression
-- Returns the type cell for this expression.
partial def inferExpr
    (ctx : Context)
    (e : Expr)
    (steps : IO.Ref (Array Propagators.PropStep)) :
    IO (Propagators.Cell Ty) := do
  match e with
  | .intLit _ =>
    let c ← Propagators.Cell.new (α := Ty)
    let _ ← c.write .int
    return c

  | .boolLit _ =>
    let c ← Propagators.Cell.new (α := Ty)
    let _ ← c.write .bool
    return c

  | .var x =>
    match ctx.lookup x with
    | some c => return c
    | none   =>
      let c ← Propagators.Cell.new (α := Ty)
      let _ ← c.write .typeError  -- unbound variable
      return c

  | .add e1 e2 =>
    let c1 ← inferExpr ctx e1 steps
    let c2 ← inferExpr ctx e2 steps
    let cOut ← Propagators.Cell.new (α := Ty)
    -- Propagator: both operands must be Int. We then re-read the
    -- operand cells and merge their join into the output — so if an
    -- operand has jumped to typeError (e.g. Bool ⊔ Int), the error
    -- reaches the result cell too.
    steps.modify (·.push do
      let _ ← c1.write .int
      let _ ← c2.write .int
      let t1 ← c1.read; let t2 ← c2.read
      cOut.write (Ty.unify t1 t2))
    return cOut

  | .ifExpr cond t f =>
    let ccond ← inferExpr ctx cond steps
    let ct    ← inferExpr ctx t    steps
    let cf    ← inferExpr ctx f    steps
    let cOut  ← Propagators.Cell.new (α := Ty)
    -- Propagator: cond must be Bool; t and f must agree; out = t
    steps.modify (·.push do
      let _ ← ccond.write .bool
      let tv ← ct.read; let fv ← cf.read
      let _ ← ct.write fv; let _ ← cf.write tv
      let tv' ← ct.read
      cOut.write tv')
    return cOut

  | .lam x body =>
    let cDom  ← Propagators.Cell.new (α := Ty)
    let ctx'  := (x, cDom) :: ctx
    let cBody ← inferExpr ctx' body steps
    let cOut  ← Propagators.Cell.new (α := Ty)
    -- Propagator: out = dom → body
    steps.modify (·.push do
      let d ← cDom.read; let b ← cBody.read
      cOut.write (.fun_ d b))
    return cOut

  | .app fn arg =>
    let cFn  ← inferExpr ctx fn  steps
    let cArg ← inferExpr ctx arg steps
    let cOut ← Propagators.Cell.new (α := Ty)
    -- Propagator: fn must be (arg → out)
    steps.modify (·.push do
      let fv ← cFn.read; let av ← cArg.read
      let ov ← cOut.read
      match fv with
      | .fun_ d r =>
        let _ ← cArg.write d   -- arg type must match domain
        let _ ← cOut.write r   -- result type from codomain
        let _ ← cFn.write (.fun_ av ov) -- refine fn type
        return false            -- (actual change tracked by writes)
      | .unknown =>
        let _ ← cFn.write (.fun_ av ov)
        return false
      | _ =>
        let _ ← cFn.write .typeError; return false)
    return cOut

  | .letExpr x e body =>
    let ce   ← inferExpr ctx  e    steps
    let ctx' := (x, ce) :: ctx
    inferExpr ctx' body steps

end TypeInfer
\end{minted}

\section{Running the Type Inferencer}

\begin{minted}{lean4}
namespace TypeInfer

def infer (e : Expr) : IO Ty := do
  let stepsRef ← IO.mkRef (Array.empty : Array Propagators.PropStep)
  let cResult  ← inferExpr [] e stepsRef
  let steps    ← stepsRef.get
  Propagators.runToFixpoint steps
  cResult.read

-- Helper to show types
def Ty.toString : Ty → String
  | .unknown       => "?"
  | .int           => "Int"
  | .bool          => "Bool"
  | .fun_ d r      => s!"({Ty.toString d} → {Ty.toString r})"
  | .typeError     => "TypeError"

-- Examples
def ex1 : Expr := .add (.intLit 1) (.intLit 2)
-- Expected: Int

def ex2 : Expr := .lam "x" (.add (.var "x") (.intLit 1))
-- Expected: (Int → Int)

def ex3 : Expr := .ifExpr (.boolLit true) (.intLit 1) (.intLit 2)
-- Expected: Int

def ex4 : Expr := .add (.boolLit true) (.intLit 1)
-- Expected: TypeError (Bool ≠ Int)

def runExamples : IO Unit := do
  for (name, e) in [("1 + 2", ex1), ("λx. x+1", ex2),
                    ("if true 1 2", ex3), ("true + 1", ex4)] do
    let ty ← infer e
    IO.println s!"{name} : {Ty.toString ty}"

end TypeInfer
\end{minted}

Expected output:
\begin{minted}{text}
1 + 2 : Int
λx. x+1 : (Int → Int)
if true 1 2 : Int
true + 1 : TypeError
\end{minted}

\section{What the Lattice Buys Us Here}

\begin{itemize}
  \item \textbf{Bidirectional propagation is free.}
    When we call \lc{cArg.write d} inside the \lc{app} propagator, we are
    propagating type information \emph{backwards}: from the function type to
    the argument. The lattice join ensures this never loses previously-inferred
    information.

  \item \textbf{Inconsistency is explicit.}
    \lc{TypeError} is the $\top$ element. When two incompatible type
    constraints are merged (\lc{Ty.unify .int .bool = .typeError}), the cell
    value jumps to $\top$ and stays there. No exception is thrown; the type
    error propagates naturally through the lattice.

  \item \textbf{Incrementality is free.}
    If we learn a new fact about $x$'s type (say, $x$ appears in another
    expression), we just write to $x$'s type cell and re-run the network.
    The lattice join handles merging old and new constraints correctly.

  \item \textbf{Ordering captures ``how much we know''.}
    $\bot = \texttt{Unknown} < \texttt{Int} < \texttt{TypeError} = \top$.
    The height of a cell's value in the lattice is a measure of how constrained
    it is. A fully-typed program has all cells at concrete types (neither
    $\bot$ nor $\top$).
\end{itemize}

\section{Exercises}

\begin{exercisebox}[Let-polymorphism]
  Lean's type system supports polymorphism. Extend \lc{Ty} with a
  \lc{TVar (name : String)} constructor and modify \lc{Ty.unify} to
  do proper unification (substituting type variables).
  How does this affect the lattice structure?
\end{exercisebox}

\begin{exercisebox}[Recursive types]
  Add a \lc{letRec} construct to \lc{Expr} and extend the inferencer to
  handle it. Note: the type cell for the recursive binding must be
  initialised to \lc{Unknown} and refined by propagation, enabling
  inference of e.g.\ \lc{let rec f x = if x = 0 then 1 else x * f (x-1)}.
\end{exercisebox}

\begin{exercisebox}[Error messages]
  When a cell reaches \lc{TypeError}, record \emph{why} (which two types
  conflicted). Augment the \lc{Ty} type so that \lc{TypeError} carries a
  \lc{String} explanation, and thread error messages through \lc{Ty.unify}.
\end{exercisebox}

% ============================================================
\appendix
\chapter{Lean~4 Quick Reference and Tactic Primer}
\label{app:quickref}

\section{Basic Types and Syntax}

\begin{minted}{lean4}
-- Type universes
#check Nat       -- Nat : Type
#check Type      -- Type : Type 1
#check Prop      -- Prop : Type

-- Function types
def f : Nat → Nat → Nat := fun a b => a + b
def g (a b : Nat) : Nat := a + b    -- same, different syntax

-- Dependent types (e.g. fixed-length vectors as functions)
def myList (α : Type) (n : Nat) : Type := Fin n → α

-- Let bindings
def h : Nat :=
  let x := 5
  let y := x + 1
  x * y
\end{minted}

\section{Typeclasses}

\begin{minted}{lean4}
-- Declaring a typeclass
class MyClass (α : Type) where
  someMethod : α → Nat
  someProperty : ∀ (a : α), someMethod a ≥ 0

-- Providing an instance
instance : MyClass Nat where
  someMethod := id
  someProperty := Nat.zero_le

-- Using an instance
def useMyClass [MyClass α] (a : α) : Nat :=
  MyClass.someMethod a
\end{minted}

\section{Tactic Primer}
\label{app:primer}

A theorem in Lean is a type; a proof is a term of that type. You can write
that term directly, or enter \emph{tactic mode} with \lc{by} and build it
interactively: at every point there is a \emph{goal} (hypotheses above the
line, target below), and each tactic transforms the goal until none remain.
Tactics are not magic --- each one \emph{elaborates to an ordinary term} in
the end, and for the most important ones this primer shows both the tactic
and the term it stands for. Every example below is compile-checked against
the same zero-dependency toolchain as the rest of the book.

\subsection*{\lc{intro} and \lc{exact}}

\lc{intro h} moves the premise of an implication (or the bound variable of
a $\forall$) into the hypotheses; \lc{exact e} closes the goal with the
term \lc{e}, which must fit exactly.

\begin{minted}{lean4}
theorem imp_demo (p q : Prop) (hq : q) : p → q ∧ q := by
  intro hp
  exact ⟨hq, hq⟩
\end{minted}

The angle brackets \lc{⟨...⟩} build a value of any one-constructor type
(here \lc{And.intro}); this proof is literally the term
\lc{fun hp => ⟨hq, hq⟩}.

\subsection*{\lc{apply}}

\lc{apply f} matches the goal against the \emph{conclusion} of \lc{f} and
turns \lc{f}'s remaining arguments into new goals --- backward reasoning.

\begin{minted}{lean4}
theorem apply_demo (a b : Nat) (h : a = b) : a ≤ b := by
  apply Nat.le_of_eq   -- goal becomes: a = b
  exact h
\end{minted}

\subsection*{\lc{rfl}}

Closes goals of the form \lc{a = a} --- \emph{up to computation}. Lean
evaluates both sides, so \lc{rfl} proves more than it looks like it should:

\begin{minted}{lean4}
example : 2 + 2 = 4 := rfl
example : (fun n : Nat => n + 1) 3 = 4 := rfl
\end{minted}

Because proofs of the same proposition are definitionally equal
(\emph{proof irrelevance}), \lc{rfl} also closes equalities between
structures that differ only in proof fields --- Chapter~7's interval
antisymmetry ends this way.

\subsection*{\lc{constructor}}

Applies the (unique) constructor of the goal's type: an \lc{And} goal
splits into two goals, an \lc{Iff} into two implications, and so on. The
bullets \lc{·} focus one subgoal each.

\begin{minted}{lean4}
example : 1 ≤ 2 ∧ 2 ≤ 3 := by
  constructor
  · omega
  · omega
\end{minted}

\subsection*{\lc{cases} --- and what it desugars to}

\lc{cases x} splits the proof into one branch per constructor of \lc{x}'s
inductive type. This is not a primitive: every inductive type \lc{T} comes
with a \emph{recursor} \lc{T.rec}, a function that says ``to define
something for all \lc{T}, handle each constructor''. \lc{cases} (and
\lc{induction}) are sugar for applying it.

\begin{minted}{lean4}
theorem bool_or_not (b : Bool) : (b || !b) = true := by
  cases b with
  | false => rfl
  | true  => rfl

-- the same proof written directly against the recursor
theorem bool_or_not' (b : Bool) : (b || !b) = true :=
  Bool.rec (motive := fun x => (x || !x) = true) rfl rfl b
\end{minted}

The \lc{motive} is the family of statements being proved --- ``what does
the goal look like as a function of \lc{x}?'' --- and the two \lc{rfl}s are
the \lc{false} and \lc{true} branches.

\subsection*{\lc{induction} --- and what it desugars to}

Same recursor, but for recursive types the branch for a recursive
constructor receives an \emph{induction hypothesis}. Note how precisely
the term mirrors the tactic script:

\begin{minted}{lean4}
theorem zero_add_tac (n : Nat) : 0 + n = n := by
  induction n with
  | zero => rfl
  | succ k ih => rw [Nat.add_succ, ih]

-- the same proof as an explicit Nat.rec term
theorem zero_add_term (n : Nat) : 0 + n = n :=
  Nat.rec (motive := fun m => 0 + m = m)
    rfl
    (fun k ih => congrArg Nat.succ ih)
    n
\end{minted}

\lc{Nat.rec} takes the base case (a proof for \lc{0}), the step case (a
function from \lc{k} and the induction hypothesis \lc{ih} to a proof for
\lc{k + 1}), and the number to recurse on. \emph{Induction is just
recursion over proofs.}

\subsection*{\lc{obtain} / \lc{rcases} --- and what they desugar to}

Both destructure a hypothesis by pattern: \lc{obtain ⟨k, hk⟩ := h} takes
\lc{h : ∃ k, P k} apart into a witness \lc{k} and a proof \lc{hk : P k}.
(\lc{rcases h with ⟨k, hk⟩} is the same with recursive patterns like
\lc{⟨a, ⟨b, c⟩⟩}.) Underneath it is the eliminator of \lc{Exists} --- or,
equivalently, a \lc{match}:

\begin{minted}{lean4}
theorem pos_of_succ {n : Nat} (h : ∃ k, n = k + 1) : 0 < n := by
  obtain ⟨k, hk⟩ := h
  subst hk
  exact Nat.succ_pos k

-- as an application of the eliminator …
theorem pos_of_succ' {n : Nat} (h : ∃ k, n = k + 1) : 0 < n :=
  Exists.elim h (fun k hk => hk ▸ Nat.succ_pos k)

-- … or as a match
theorem pos_of_succ'' {n : Nat} (h : ∃ k, n = k + 1) : 0 < n :=
  match h with
  | ⟨k, hk⟩ => hk ▸ Nat.succ_pos k
\end{minted}

(\lc{h ▸ e} rewrites the goal with the equation \lc{h} and then supplies
\lc{e}; it is the term-level cousin of \lc{rw}.)

\subsection*{\lc{rename\_i}}

Tactics like \lc{cases} sometimes introduce hypotheses with inaccessible
machine-generated names (shown with a dagger, e.g.\ \lc{a✝}).
\lc{rename\_i x} gives the most recent such hypothesis the name \lc{x} so
you can refer to it.

\begin{minted}{lean4}
theorem rename_demo (o : Option Nat) (h : o ≠ none) : ∃ n, o = some n := by
  cases o
  · exact absurd rfl h
  · rename_i n     -- name the anonymous value introduced by `cases`
    exact ⟨n, rfl⟩
\end{minted}

\subsection*{\lc{by\_cases} --- and the \lc{Decidable} \lc{if} underneath}

\lc{by\_cases h : p} splits the proof into a branch where \lc{h : p} and
one where \lc{h : ¬p}. In this book's zero-dependency setting it needs
\lc{p} to be \emph{decidable} --- there must be an algorithm answering
\lc{p} --- because it is sugar for a dependent \lc{if}:

\begin{minted}{lean4}
theorem em_nat (n : Nat) : n = 0 ∨ n ≠ 0 := by
  by_cases h : n = 0
  · exact Or.inl h
  · exact Or.inr h

theorem em_nat' (n : Nat) : n = 0 ∨ n ≠ 0 :=
  if h : n = 0 then Or.inl h else Or.inr h
\end{minted}

\subsection*{\lc{exfalso}}

Replaces any goal with \lc{False}. Useful when the hypotheses are
contradictory and the goal itself is irrelevant:

\begin{minted}{lean4}
theorem from_absurd_bounds (n : Nat) (h1 : n < 2) (h2 : 2 < n) :
    n = 42 := by
  exfalso
  omega
\end{minted}

\subsection*{\lc{subst}}

Given \lc{hab : a = b} where one side is a local variable, \lc{subst hab}
eliminates the variable, rewriting it away everywhere:

\begin{minted}{lean4}
theorem subst_demo (a b : Nat) (hab : a = b) (ha : a ≤ 5) : b ≤ 5 := by
  subst hab
  exact ha
\end{minted}

\subsection*{\lc{calc} --- and the \lc{Trans} instances underneath}

\lc{calc} writes a chain of (in)equalities the way you would on paper.
Each step is glued to the previous one by \lc{Trans.trans}, a typeclass
method --- which is why Chapter~1's chain had to be stated with \lc{≤}
(core Lean ships a \lc{Trans} instance for \lc{≤} but not for \lc{≥}):

\begin{minted}{lean4}
theorem calc_demo (a b c d : Nat)
    (h1 : a ≤ b) (h2 : b ≤ c) (h3 : c ≤ d) : a ≤ d :=
  calc a ≤ b := h1
    _ ≤ c := h2
    _ ≤ d := h3

theorem calc_demo' (a b c d : Nat)
    (h1 : a ≤ b) (h2 : b ≤ c) (h3 : c ≤ d) : a ≤ d :=
  Trans.trans (Trans.trans h1 h2) h3
\end{minted}

\subsection*{\lc{split}}

Where \lc{cases} splits on a \emph{value}, \lc{split} splits on the
control flow of the \emph{goal}: an \lc{if} (or \lc{match}) in the target
becomes one goal per branch, with the branch condition as a hypothesis.

\begin{minted}{lean4}
def absDiff (a b : Int) : Int := if a ≤ b then b - a else a - b

theorem absDiff_nonneg (a b : Int) : 0 ≤ absDiff a b := by
  simp only [absDiff]
  split <;> omega
\end{minted}

The combinator \lc{t1 <;> t2} runs \lc{t1}, then runs \lc{t2} on
\emph{every} goal it produced --- the workhorse of this book's
``case-split everything, then \lc{simp}/\lc{omega}'' proofs.

\subsection*{\lc{change}}

Replaces the goal by a \emph{definitionally equal} one --- nothing is
proved, the goal is merely re-displayed in a shape other tactics can
match:

\begin{minted}{lean4}
example : (fun n : Nat => n + 2) 3 = 5 := by
  change 3 + 2 = 5
  rfl
\end{minted}

\subsection*{\lc{simp} and \lc{simp\_all}}

\lc{simp} rewrites the goal to normal form with a database of equations
(extendable per call: \lc{simp [FlatNat.le]} also unfolds that
definition). \lc{simp\_all} does the same to the hypotheses \emph{and}
uses them as rewrite rules --- the closer for most of this book's
case-bash proofs.

\begin{minted}{lean4}
theorem simp_demo (xs : List Nat) : (xs ++ []).length = xs.length := by
  simp

theorem simp_all_demo (xs : List Nat) (h : xs = [1, 2]) :
    xs.length = 2 := by
  simp_all
\end{minted}

\lc{simp only [f]} restricts to the given lemmas --- predictable, no
database.

\subsection*{\lc{omega} and \lc{decide}}

\lc{omega} is a decision procedure for linear arithmetic over \lc{Int} and
\lc{Nat} (with \lc{+}, \lc{-}, \lc{≤}, \lc{<}, \lc{=}, \lc{min}/\lc{max},
multiplication by constants). \lc{decide} evaluates any \emph{decidable}
proposition to \lc{true} --- it works whenever there is a finite
computation to run, e.g.\ statements about all \lc{Bool}s.

\begin{minted}{lean4}
example (a b : Nat) (h : a + 2 * b = 10) (hb : 3 ≤ b) : a ≤ 4 := by
  omega

example : (15 % 4 = 3) ∧ (2 ^ 10 = 1024) := by decide
example : ∀ b : Bool, (b && !b) = false := by decide
\end{minted}

\subsection*{\lc{have} and \lc{show}}

\lc{have h : P := proof} adds an intermediate fact; \lc{show P} restates
the current goal (up to definitional equality --- a gentler \lc{change},
often used purely for readability):

\begin{minted}{lean4}
theorem have_demo (a : Nat) (h : a ≤ 3) : a ≤ 10 := by
  have h' : (3 : Nat) ≤ 10 := by omega
  show a ≤ 10
  exact Nat.le_trans h h'
\end{minted}

A \lc{have} with no proof term, like Chapter~7's
\lc{have hab' : loLe l1 l2 ∧ hiLe h2 h1 := hab}, re-types an existing
hypothesis at a definitionally equal type --- no work, just unfolding.

\subsection*{\lc{funext} and \lc{congr}}

\lc{funext} reduces \lc{f = g} to \lc{∀ x, f x = g x} (functions are equal
when they agree everywhere). \lc{congr 1} goes the other way for
constructors: it reduces \lc{⟨a₁, b₁⟩ = ⟨a₂, b₂⟩} to \lc{a₁ = a₂} and
\lc{b₁ = b₂}, closing proof fields automatically by proof irrelevance:

\begin{minted}{lean4}
theorem funext_demo : (fun n : Nat => n + 0) = (fun n : Nat => n) := by
  funext n
  rfl

structure Point where
  x : Int
  y : Int

theorem congr_demo (p q : Point) (hx : p.x = q.x) (hy : p.y = q.y) :
    p = q := by
  cases p; cases q
  congr 1   -- splits ⟨x₁,y₁⟩ = ⟨x₂,y₂⟩ into x₁ = x₂ and y₁ = y₂
\end{minted}

(Here \lc{congr 1} leaves the two component goals, which \lc{hx} and
\lc{hy} close --- after \lc{cases}, they already \emph{are} those goals,
so Lean finds them by assumption.)

\section{Programming Constructs}
\label{app:programming}

\subsection*{\lc{do}-notation, \lc{mut}, \lc{for}, \lc{while}, and \lc{Id.run}}

A \lc{do} block sequences actions in a monad. Inside it, \lc{let x ← a}
runs the action \lc{a} and binds its result; \lc{let mut} declares a
variable that later statements may reassign; \lc{for} and \lc{while} loop.
None of this is primitive --- the compiler \emph{desugars} it: statements
become nested calls of the monad's bind operator, and every \lc{mut}
variable becomes an accumulator threaded through a fold. Purely
functional code wearing an imperative costume:

\begin{minted}{lean4}
def sumTo (n : Nat) : Nat := Id.run do
  let mut acc := 0
  for i in [0:n+1] do
    acc := acc + i
  return acc

#eval sumTo 10   -- 55

-- what the mut/for pair compiles down to (morally): a fold
def sumTo' (n : Nat) : Nat :=
  (List.range (n + 1)).foldl (fun acc i => acc + i) 0

#eval sumTo' 10  -- 55
\end{minted}

\lc{Id} is the do-nothing monad (\lc{Id α = α}), and \lc{Id.run} just
unwraps it, so \lc{sumTo} is an ordinary pure function --- this is how the
Sudoku solver's \lc{Id.run do} blocks work. \lc{while} loops need a monad
that permits possible non-termination, such as \lc{IO}:

\begin{minted}{lean4}
def firstPowerAbove (bound : Nat) : IO Nat := do
  let mut p := 1
  while p ≤ bound do
    p := p * 2
  return p

#eval firstPowerAbove 100   -- 128
\end{minted}

\subsection*{\lc{IO.Ref}}

Genuine mutable state --- the heart of a propagator \lc{Cell} --- lives in
\lc{IO.Ref α}, a mutable reference allocated with \lc{IO.mkRef}:

\begin{minted}{lean4}
def counterDemo : IO Nat := do
  let r ← IO.mkRef 0     -- allocate a mutable reference holding 0
  r.set 5                -- overwrite the contents
  r.modify (· + 1)       -- apply a function to the contents
  r.get                  -- read it back

#eval counterDemo        -- 6
\end{minted}

\subsection*{\lc{partial def}}

Lean normally insists every function provably terminates. \lc{partial def}
opts a definition out of that check: it can still be compiled and run, but
Lean treats its body as opaque for proofs. We used it for \lc{inferExpr}
in Chapter~9, whose recursion structure is obvious to a human and tedious
for a termination checker:

\begin{minted}{lean4}
partial def collatzSteps (n : Nat) : Nat :=
  if n ≤ 1 then 0
  else if n % 2 == 0 then 1 + collatzSteps (n / 2)
  else 1 + collatzSteps (3 * n + 1)

#eval collatzSteps 6     -- 8
\end{minted}

Whether \emph{this} function terminates for all inputs is a famous open
problem --- exactly the kind of definition \lc{partial} exists for.

\subsection*{Structures, typeclasses, and \lc{deriving}}

A \lc{structure} bundles named fields (possibly including \emph{proof}
fields, like \lc{Interval.range}'s \lc{ok}); \lc{cases}/\lc{rcases} take
values apart and \lc{⟨...⟩} builds them. A \lc{class} is a structure whose
instances Lean finds automatically for square-bracket parameters like
\lc{[PartialOrder α]}; \lc{class B extends A} embeds one class in another
(our \lc{Lattice extends PartialOrder}). Finally,
\lc{deriving instance DecidableEq for FlatNat} asks Lean to generate an
instance mechanically --- available for simple inductives without proof
fields; Chapter~7's \lc{Interval} needed a hand-written one.

\section{Tactics at a Glance}

Every tactic used in this book, in one table. Sections above give worked
examples of each.

\begin{table}[H]
  \centering
  \small
  \begin{tabular}{ll}
    \toprule
    Tactic & Usage \\
    \midrule
    \lc{rfl}          & Prove \lc{a = a}, up to computation \\
    \lc{exact e}      & Close the goal with term \lc{e} \\
    \lc{apply f}      & Match goal against \lc{f}'s conclusion; args become goals \\
    \lc{intro h}      & Move premise / $\forall$-variable into hypotheses \\
    \lc{have h : P := e} & Add an intermediate fact \\
    \lc{show P}       & Restate the goal (definitional equality) \\
    \lc{cases x}      & One goal per constructor of \lc{x} (sugar for \lc{T.rec}) \\
    \lc{rcases h with ⟨a, b⟩} & \lc{cases} with nested destructuring patterns \\
    \lc{obtain ⟨a, b⟩ := h}   & Same as \lc{rcases}, different spelling \\
    \lc{rename\_i x}  & Name the last inaccessible hypothesis \\
    \lc{induction n}  & \lc{cases} + induction hypotheses (sugar for \lc{T.rec}) \\
    \lc{constructor}  & Apply the goal type's constructor (\lc{And}, \lc{Iff}, …) \\
    \lc{by\_cases h : p} & Split on decidable \lc{p} (sugar for dependent \lc{if}) \\
    \lc{exfalso}      & Replace the goal with \lc{False} \\
    \lc{subst h}      & Eliminate a variable using \lc{h : x = e} \\
    \lc{calc}         & Chain of steps glued by \lc{Trans.trans} \\
    \lc{rw [h]}       & Rewrite the goal with equation \lc{h} \\
    \lc{simp}         & Normalise with the simp-lemma database \\
    \lc{simp only [f]}& Simp restricted to the given lemmas \\
    \lc{simp\_all}    & Simp goal and hypotheses, using hypotheses \\
    \lc{split}        & One goal per branch of an \lc{if}/\lc{match} in the goal \\
    \lc{change P}     & Replace goal by definitionally equal \lc{P} \\
    \lc{omega}        & Decide linear arithmetic over \lc{Int}/\lc{Nat} \\
    \lc{decide}       & Evaluate a decidable proposition \\
    \lc{funext}       & \lc{f = g} from \lc{∀ x, f x = g x} \\
    \lc{congr n}      & Split constructor equality into field equalities \\
    \lc{t1 <;> t2}    & Run \lc{t2} on every goal produced by \lc{t1} \\
    \lc{·}~(bullet)   & Focus the first remaining goal \\
    \bottomrule
  \end{tabular}
  \caption{The complete tactic vocabulary of this book.}
\end{table}

\section{Unicode Input}

\begin{table}[h]
  \centering
  \begin{tabular}{lll}
    \toprule
    Symbol & Input & Meaning \\
    \midrule
    \lc{≤}  & \texttt{\textbackslash{}le}   & Less-than-or-equal \\
    \lc{≥}  & \texttt{\textbackslash{}ge}   & Greater-than-or-equal \\
    \lc{→}  & \texttt{\textbackslash{}to}   & Function/implication arrow \\
    \lc{←}  & \texttt{\textbackslash{}l}    & Reverse arrow \\
    \lc{↔}  & \texttt{\textbackslash{}iff}  & If and only if \\
    \lc{∀}  & \texttt{\textbackslash{}fa}   & For all \\
    \lc{∃}  & \texttt{\textbackslash{}ex}   & There exists \\
    \lc{⊓}  & \texttt{\textbackslash{}sqcap}& Meet (inf) \\
    \lc{⊔}  & \texttt{\textbackslash{}sqcup}& Join (sup) \\
    \lc{⊥}  & \texttt{\textbackslash{}bot}  & Bottom \\
    \lc{⊤}  & \texttt{\textbackslash{}top}  & Top \\
    \lc{∈}  & \texttt{\textbackslash{}in}   & Set membership \\
    \lc{∉}  & \texttt{\textbackslash{}notin}& Non-membership \\
    \lc{α}  & \texttt{\textbackslash{}a}    & Greek alpha \\
    \lc{β}  & \texttt{\textbackslash{}b}    & Greek beta \\
    \lc{ℕ}  & \texttt{\textbackslash{}N}    & Natural numbers \\
    \lc{∧}  & \texttt{\textbackslash{}and}  & Logical and \\
    \lc{∨}  & \texttt{\textbackslash{}or}   & Logical or \\
    \lc{¬}  & \texttt{\textbackslash{}n}    & Logical not \\
    \bottomrule
  \end{tabular}
  \caption{Lean~4 Unicode inputs used in this book.}
\end{table}

% ── Bibliography / Further Reading ──────────────────────────
\chapter*{Selected Solutions}
\addcontentsline{toc}{chapter}{Selected Solutions}

This appendix contains solutions to exercises with a single definitive
Lean~4 answer. Open-ended exercises are left to the reader. All solutions
assume the book's definitions are in scope --- \lc{open OrderBook} for
Part~I, and additionally \lc{namespace Propagators} for the Chapter~6
solutions.

% ── Chapter 1 ────────────────────────────────────────────────
\section*{Chapter 1}

\textbf{Exercise 1.1} (Reflexivity and symmetry)

\begin{minted}{lean4}
theorem symm_antisymm_eq {α : Type} (r : Rel α)
    (hs : Symmetric r) (ha : Antisymm r) :
    ∀ a b : α, r a b → a = b := by
  intro a b h
  exact ha a b h (hs a b h)
\end{minted}

\textbf{Exercise 1.3} (Product order)

\begin{minted}{lean4}
instance [OrderBook.PartialOrder α] [OrderBook.PartialOrder β] :
    OrderBook.PartialOrder (α × β) where
  le p q := PartialOrder.le p.1 q.1 ∧ PartialOrder.le p.2 q.2
  le_refl := fun ⟨a, b⟩ =>
    ⟨PartialOrder.le_refl a, PartialOrder.le_refl b⟩
  le_antisymm := fun ⟨a1,b1⟩ ⟨a2,b2⟩ ⟨h1,h2⟩ ⟨h3,h4⟩ => by
    have := PartialOrder.le_antisymm _ _ h1 h3
    have := PartialOrder.le_antisymm _ _ h2 h4
    simp_all
  le_trans := fun ⟨a1,b1⟩ ⟨a2,b2⟩ ⟨a3,b3⟩ ⟨h1,h2⟩ ⟨h3,h4⟩ =>
    ⟨PartialOrder.le_trans _ _ _ h1 h3,
     PartialOrder.le_trans _ _ _ h2 h4⟩
\end{minted}

% ── Chapter 2 ────────────────────────────────────────────────
\section*{Chapter 2}

\textbf{Exercise 2.1} (Monotone composition)

\begin{minted}{lean4}
-- NB: no [LE _] binders — our PartialOrder already provides the ≤
-- notation; a second LE instance would introduce an unrelated order
-- and the hypotheses would no longer match Monotone's.
theorem antitone_comp_antitone {α β γ : Type}
    [PartialOrder α] [PartialOrder β] [PartialOrder γ]
    {f : α → β} {g : β → γ}
    (hf : ∀ a b : α, a ≤ b → f b ≤ f a)
    (hg : ∀ a b : β, a ≤ b → g b ≤ g a) :
    Monotone (g ∘ f) where
  map_le := fun a b h => hg (f b) (f a) (hf a b h)
\end{minted}

\textbf{Exercise 2.2} (Fixed points --- ascending chain)

\begin{minted}{lean4}
def iter {α : Type} (f : α → α) : Nat → α → α
  | 0,     x => x
  | n + 1, x => f (iter f n x)

theorem iter_chain {α : Type} [PartialOrder α]
    {f : α → α} (hf : Monotone f)
    {a : α} (ha : a ≤ f a) :
    ∀ n : Nat, iter f n a ≤ iter f (n + 1) a := by
  intro n
  induction n with
  | zero => exact ha
  | succ n ih => exact hf.map_le _ _ ih
\end{minted}

\textbf{Exercise 2.3} (Identity is monotone)

\begin{minted}{lean4}
theorem Monotone.id {α : Type} [PartialOrder α] :
    Monotone (id : α → α) where
  map_le := fun _ _ h => h
\end{minted}

% ── Chapter 3 ────────────────────────────────────────────────
\section*{Chapter 3}

\textbf{Exercise 3.1} (Interval information order)

\begin{minted}{lean4}
structure Interval where
  l      : Int
  h      : Int
  l_le_h : l ≤ h

instance : PartialOrder Interval where
  -- Information order: i ≤ j iff j is contained in i
  -- (narrower = more informative = higher).
  le := fun i j => i.l ≤ j.l ∧ j.h ≤ i.h
  le_refl := fun a => ⟨Int.le_refl a.l, Int.le_refl a.h⟩
  le_antisymm := fun a b hyp₁ hyp₂ => by
    cases a; cases b
    rename_i l₁ h₁ l₁_le_h₁ l₂ h₂ l₂_le_h₂
    congr 1
    · exact Int.le_antisymm hyp₁.left hyp₂.left
    · exact Int.le_antisymm hyp₂.right hyp₁.right
    -- the l_le_h field is closed by proof irrelevance
  le_trans := fun a b c hyp₁ hyp₂ => by
    rcases a with ⟨a_l, a_h, a_le⟩
    rcases b with ⟨b_l, b_h, b_le⟩
    rcases c with ⟨c_l, c_h, c_le⟩
    have hl_ab : a_l ≤ b_l := hyp₁.left
    have hl_bc : b_l ≤ c_l := hyp₂.left
    have hh_ba : b_h ≤ a_h := hyp₁.right
    have hh_cb : c_h ≤ b_h := hyp₂.right
    exact ⟨Int.le_trans hl_ab hl_bc, Int.le_trans hh_cb hh_ba⟩
\end{minted}

Meets always exist: the hull $[\min(l_1,l_2), \max(h_1,h_2)]$ is the
greatest lower bound. Joins exist exactly when the intervals overlap ---
for disjoint intervals every candidate join would have to be an ``empty
interval'', which this type cannot represent. Chapter~7's lattice adds it
(as $\top$), plus unbounded ends (as $\bot$).

\textbf{Exercise 3.3} (Characterising the order)

\begin{minted}{lean4}
theorem le_iff_sup_eq {α : Type} [Lattice α] {a b : α} :
    a ≤ b ↔ a ⊔ b = b := by
  apply Iff.intro
  · intro h
    apply PartialOrder.le_antisymm
    · apply Lattice.le_sup
      · exact h
      · exact PartialOrder.le_refl b
    · exact Lattice.sup_le_right a b
  · intro h
    rw [← h]
    exact Lattice.sup_le_left a b
\end{minted}

\textbf{Exercise 3.4} (Product lattice)

\begin{minted}{lean4}
instance {α β : Type} [Lattice α] [Lattice β] : Lattice (α × β) where
  le            := fun a b => (a.fst ≤ b.fst) ∧ (a.snd ≤ b.snd)
  le_refl       := fun a => by
    apply And.intro
    · apply PartialOrder.le_refl
    · apply PartialOrder.le_refl
  le_antisymm {a b} := fun hyp₁ hyp₂ => by
    rcases a with ⟨l₁, r₁⟩; rcases b with ⟨l₂, r₂⟩; congr 1
    · exact PartialOrder.le_antisymm _ _ hyp₁.left hyp₂.left
    · exact PartialOrder.le_antisymm _ _ hyp₁.right hyp₂.right
  le_trans {a b c} := fun hyp₁ hyp₂ => by
    apply And.intro
    · exact PartialOrder.le_trans _ _ _ hyp₁.left hyp₂.left
    · exact PartialOrder.le_trans _ _ _ hyp₁.right hyp₂.right
  inf           := fun a b =>
    ⟨Lattice.inf a.fst b.fst, Lattice.inf a.snd b.snd⟩
  sup           := fun a b =>
    ⟨Lattice.sup a.fst b.fst, Lattice.sup a.snd b.snd⟩
  inf_le_left   := fun a b => ⟨Lattice.inf_le_left a.fst b.fst,
                                Lattice.inf_le_left a.snd b.snd⟩
  inf_le_right  := fun a b => ⟨Lattice.inf_le_right a.fst b.fst,
                                Lattice.inf_le_right a.snd b.snd⟩
  le_inf        := fun a b c hyp₁ hyp₂ =>
    ⟨Lattice.le_inf a.fst b.fst c.fst hyp₁.left hyp₂.left,
     Lattice.le_inf a.snd b.snd c.snd hyp₁.right hyp₂.right⟩
  sup_le_left   := fun a b => ⟨Lattice.sup_le_left a.fst b.fst,
                                Lattice.sup_le_left a.snd b.snd⟩
  sup_le_right  := fun a b => ⟨Lattice.sup_le_right a.fst b.fst,
                                Lattice.sup_le_right a.snd b.snd⟩
  le_sup        := fun a b c hyp₁ hyp₂ =>
    ⟨Lattice.le_sup a.fst b.fst c.fst hyp₁.left hyp₂.left,
     Lattice.le_sup a.snd b.snd c.snd hyp₁.right hyp₂.right⟩
\end{minted}

% ── Chapter 6 ────────────────────────────────────────────────
\section*{Chapter 6}

\textbf{Exercise 6.1} (Three-cell multiplication constraint)

The pattern mirrors the addition network: one forward propagator and two
backward propagators, one for each variable that can be derived from the
other two. Division by zero writes \lc{.conflict}.

\begin{minted}{lean4}
def mulProp (cx cy cz : Cell FlatNat) : PropStep := do
  let x ← cx.read
  let y ← cy.read
  match x, y with
    | .known a, .known b => cz.write (.known (a * b))
    | _,        _        => return false

-- z / x → y  (or conflict if x does not divide z)
def divProp (cz cx cy : Cell FlatNat) : PropStep := do
  let z ← cz.read
  let x ← cx.read
  match z, x with
    | .known z, .known x =>
      if x ∣ z then cy.write (.known (z / x))
               else cy.write (.conflict)
    | _, _ => return false

def mulNetwork : IO Unit := do
  let cx ← Cell.new (α := FlatNat)
  let cy ← Cell.new (α := FlatNat)
  let cz ← Cell.new (α := FlatNat)

  let _ ← cx.write (.known 3)
  let _ ← cz.write (.known 12)

  let net : Network := Network.mk #[
    mulProp cx cy cz,
    divProp cz cx cy,   -- z / x → y
    divProp cz cy cx    -- z / y → x
  ]
  net.run

  let yVal ← cy.read
  match yVal with
    | .known n  => IO.println s!"y = {n}"   -- y = 4
    | .unknown  => IO.println "y is unknown"
    | .conflict => IO.println "contradiction"

#eval mulNetwork   -- y = 4
\end{minted}

The three propagators cover all three ``solve for one variable'' cases.
Remove \lc{divProp cz cx cy} and \lc{y} stays unknown; remove
\lc{mulProp} and you lose the ability to derive \lc{cz} from both inputs.

\textbf{Exercise 6.2} (Conflict detection)

Two independent writes of different values to the same cell produce
\lc{.conflict} via the join in \lc{Cell.write}: \lc{known 3 ⊔ known 5 = conflict}.
No propagators are needed --- the conflict arises purely from the two seeds.

\begin{minted}{lean4}
def conflictExample : IO Unit := do
  let cx ← Cell.new (α := FlatNat)

  -- Two independent sources disagree about x
  let _ ← cx.write (.known 3)
  let _ ← cx.write (.known 5)

  let xVal ← cx.read
  match xVal with
    | .conflict => IO.println "contradiction (expected)"
    | .known n  => IO.println s!"x = {n}"
    | .unknown  => IO.println "x unknown"

#eval conflictExample   -- contradiction (expected)
\end{minted}

The key is in \lc{Cell.write}: it merges the new value with the old via
\lc{sup}, so the second write computes
\lc{known 3 ⊔ known 5 = conflict} and stores that.
From this point any further \lc{write} leaves \lc{conflict} unchanged
(since \lc{conflict ⊔ x = conflict} for all \lc{x}), and any propagator
reading \lc{cx} will receive \lc{conflict} and propagate it downstream.

To demonstrate downstream propagation, add a dependent cell:

\begin{minted}{lean4}
def conflictPropagates : IO Unit := do
  let cx   ← Cell.new (α := FlatNat)
  let cy   ← Cell.new (α := FlatNat)
  let csum ← Cell.new (α := FlatNat)

  let _ ← cx.write (.known 3)
  let _ ← cx.write (.known 5)   -- cx is now conflict
  let _ ← cy.write (.known 7)

  let net : Network := Network.mk #[addProp cx cy csum]
  net.run

  let s ← csum.read
  match s with
    | .conflict => IO.println "sum is conflict (expected)"
    | .known n  => IO.println s!"sum = {n}"
    | .unknown  => IO.println "sum unknown"

#eval conflictPropagates   -- sum is conflict (expected)
\end{minted}

\FloatBarrier
\chapter*{Further Reading}
\addcontentsline{toc}{chapter}{Further Reading}

\begin{itemize}
  \item \textbf{Davey \& Priestley}, \emph{Introduction to Lattices and Order},
    2nd ed., Cambridge University Press, 2002.
    The standard introductory text on lattice theory, accessible and rigorous.

  \item \textbf{Radul \& Sussman}, \emph{The Art of the Propagator}, MIT CSAIL
    Technical Report, 2009.
    The foundational paper for the propagator model used in this book.

  \item \textbf{Tate, Stepp, Lerner \& Chambers}, \emph{Equality Saturation: a
    New Approach to Optimization}, POPL 2009.
    An application of lattice-based fixed-point algorithms to compiler optimization.

  \item \textbf{Nielson, Nielson \& Hankin}, \emph{Principles of Program Analysis},
    Springer, 1999.
    Covers dataflow analysis (a lattice-based framework) and abstract interpretation.

  \item \textbf{The Lean~4 Documentation}, \url{https://lean-lang.org}.
    Official reference for Lean~4 syntax, tactics, and the core library.

  \item \textbf{Mathlib4}, \url{https://leanprover-community.github.io/mathlib4_docs/}.
    The community mathematics library for Lean~4.
    Once you have finished this book, exploring \texttt{Mathlib.Order} will
    show you how the constructions here scale to a large, verified library.
\end{itemize}

\end{document}
